% Generated mechanically from the frozen dga.tex target.
% Do not edit this generated file; edit the bound locale source instead.

% OK, start here.
%

\setcounter{chapter}{21}
\chapter{微分分次代数}



\phantomsection
\label{dga-section-phantom}


\section{引言}
\label{dga-section-introduction}

\noindent
本章讨论微分分次代数、模、范畴等。一个基本参考文献是
\cite{Keller-Deriving}，综述论文可见 \cite{Keller-survey}。

\medskip\noindent
由于 Stacks project 不在意叙述篇幅，我们先在微分分次模范畴的框架中
发展这些材料。此后，再在微分分次范畴上的微分分次模框架中重新完成
这些构造。



\section{约定}
\label{dga-section-conventions}

\noindent
本章继续采用如下约定：{\it 环}指有 $1$ 的交换环。若 $R$ 是环，
则一个 {\it $R$-代数 $A$} 是一个 $R$-模 $A$，连同一个 $R$-双线性映射
$A \times A \to A$（乘法），使得该乘法结合且有单位元。换言之，
它们是含幺结合 $R$-代数，并且结构映射 $R \to A$ 的像包含在 $A$ 的中心中。

\medskip\noindent
{\bf 符号规则。} 本章将研究通常配备微分的分次代数和分次模。
底层复形上的符号规则总是采用（或兼容于）《更多代数》第
\ref{more-algebra-section-sign-rules} 节中引入的规则。
这有时会使乘法结构以意外方式发生扭曲，特别是在考虑左模，
或者左模与右模之间的关系时。





\section{微分分次代数}
\label{dga-section-dga}


\noindent
这里只给出定义。

\begin{definition}
\label{dga-definition-dga}
设 $R$ 是交换环。一个 {\it $R$ 上的微分分次代数}是以下两者之一：
\begin{enumerate}
\item 一个 $R$-模链复形 $A_\bullet$，配备 $R$-双线性映射
$A_n \times A_m \to A_{n + m}$，$(a, b) \mapsto ab$，使得
$$
\text{d}_{n + m}(ab) = \text{d}_n(a)b + (-1)^n a\text{d}_m(b)
$$
并且 $\bigoplus A_n$ 成为含幺结合 $R$-代数；或者
\item 一个 $R$-模上链复形 $A^\bullet$，配备 $R$-双线性映射
$A^n \times A^m \to A^{n + m}$，$(a, b) \mapsto ab$，使得
$$
\text{d}^{n + m}(ab) = \text{d}^n(a)b + (-1)^n a\text{d}^m(b)
$$
并且 $\bigoplus A^n$ 成为含幺结合 $R$-代数。
\end{enumerate}
\end{definition}

\noindent
我们常简写为 $A = \bigoplus A_n$ 或 $A = \bigoplus A^n$，并把它看作
配备一个 $\mathbf{Z}$-分次及一个平方为零且满足上述 Leibniz 规则的
$R$-线性算子 $\text{d}$ 的含幺结合 $R$-代数。在这种情形下，
我们常说“设 $(A, \text{d})$ 是一个微分分次代数”。

\medskip\noindent
微分分次 $R$-代数 $A$ 上联系微分与乘法的 Leibniz 规则，恰好意味着
乘法映射给出一个上链复形映射
$$
\text{Tot}(A^\bullet \otimes_R A^\bullet) \to A^\bullet
$$
这里 $A^\bullet$ 表示 $A$ 的底层上链复形。

\begin{definition}
\label{dga-definition-homomorphism-dga}
微分分次代数的一个 {\it 同态}
$f : (A, \text{d}) \to (B, \text{d})$，是一个与分次及 $\text{d}$ 相容的
代数映射 $f : A \to B$。
\end{definition}

\begin{definition}
\label{dga-definition-cdga}
若对次数为 $n$ 的 $a$ 和次数为 $m$ 的 $b$ 有
$ab = (-1)^{nm}ba$，则称微分分次代数 $(A, \text{d})$ 是
{\it 交换的}。若此外对 $\deg(a)$ 为奇数的 $a$ 有 $a^2 = 0$，
则称 $A$ 是 {\it 严格交换的}。
\end{definition}

\noindent
下述定义一般都有意义，但或许只有在对交换微分分次代数取张量积时
才是“正确的”。

\begin{definition}
\label{dga-definition-tensor-product}
设 $R$ 是环。设 $(A, \text{d})$、$(B, \text{d})$ 是 $R$ 上的微分分次代数。
$A$ 与 $B$ 的 {\it 张量积微分分次代数}是代数 $A \otimes_R B$，
其乘法定义为
$$
(a \otimes b)(a' \otimes b') = (-1)^{\deg(a')\deg(b)} aa' \otimes bb'
$$
并配备由下述规则定义的微分 $\text{d}$：
$\text{d}(a \otimes b) = \text{d}(a) \otimes b + (-1)^m a \otimes \text{d}(b)$
其中 $m = \deg(a)$。
\end{definition}

\begin{lemma}
\label{dga-lemma-total-complex-tensor-product}
设 $R$ 是环。设 $(A, \text{d})$、$(B, \text{d})$ 是 $R$ 上的微分分次代数。
用 $A^\bullet$、$B^\bullet$ 表示底层上链复形。作为 $R$-模上链复形，我们有
$$
(A \otimes_R B)^\bullet = \text{Tot}(A^\bullet \otimes_R B^\bullet).
$$
\end{lemma}

\begin{proof}
回忆全复形在 $A^p \otimes_R B^q$ 上的微分由
$\text{d}_1^{p, q} + (-1)^p \text{d}_2^{p, q}$ 给出。
这恰好与定义 \ref{dga-definition-tensor-product} 中 $A \otimes_R B$
上的微分规则相同。
\end{proof}






\section{微分分次模}
\label{dga-section-modules}

\noindent
本章默认使用右模；左模将在第 \ref{dga-section-left-modules} 节中讨论。

\begin{definition}
\label{dga-definition-dgm}
设 $R$ 是环。设 $(A, \text{d})$ 是 $R$ 上的微分分次代数。
$A$ 上的一个（右）{\it 微分分次模} $M$ 是一个右 $A$-模 $M$，
配备分次 $M = \bigoplus M^n$ 和微分 $\text{d}$，使得
$M^n A^m \subset M^{n + m}$、$\text{d}(M^n) \subset M^{n + 1}$，并且
$$
\text{d}(ma) = \text{d}(m)a + (-1)^n m\text{d}(a)
$$
对 $a \in A$ 和 $m \in M^n$ 成立。微分分次模的一个 {\it 同态}
$f : M \to N$，是一个与分次及微分相容的 $A$-模映射。
（右）微分分次 $A$-模的范畴记为 $\text{Mod}_{(A, \text{d})}$。
\end{definition}

\noindent
注意，可以把 $M$ 看作（右）$R$-模的上链复形 $M^\bullet$。
事实上，对 $r \in R$ 有 $\text{d}(r) = 0$，且 $r$ 映到 $A$ 的一个
零次元素，因而 $\text{d}(mr) = \text{d}(m)r$。

\medskip\noindent
微分分次 $R$-代数 $A$ 上的微分分次 $R$-模 $M$ 中，联系微分与乘法的
Leibniz 规则恰好意味着乘法映射给出一个上链复形映射
$$
\text{Tot}(M^\bullet \otimes_R A^\bullet) \to M^\bullet
$$
这里 $A^\bullet$ 和 $M^\bullet$ 分别表示 $A$ 和 $M$ 的底层上链复形。

\begin{lemma}
\label{dga-lemma-dgm-abelian}
设 $(A, d)$ 是微分分次代数。范畴 $\text{Mod}_{(A, \text{d})}$
是阿贝尔范畴，并且具有任意极限和余极限。
\end{lemma}

\begin{proof}
取核或余核与取底层 $A$-模可交换。直和和余极限也是如此。
换言之，$\text{Mod}_{(A, \text{d})}$ 中的这些运算与到 $A$-模范畴的
遗忘函子可交换。乘积和极限则并非如此。事实上，若 $N_i$（$i \in I$）
是一族微分分次 $A$-模，则 $\text{Mod}_{(A, \text{d})}$ 中的乘积
$\prod N_i$ 由 $(\prod N_i)^n = \prod N_i^n$ 以及
$\prod N_i = \bigoplus_n (\prod N_i)^n$ 给出。因此，乘积确实与到
分次 $A$-模范畴的遗忘函子可交换。具有乘积和等化子的范畴具有极限；
见《范畴》引理 \ref{categories-lemma-limits-products-equalizers}。
\end{proof}

\noindent
因此，若 $(A, \text{d})$ 是 $R$ 上的微分分次代数，则有阿贝尔范畴间的
正合函子
$$
\text{Mod}_{(A, \text{d})} \longrightarrow \text{Comp}(R)
$$
。对于微分分次模 $M$，其上同调群 $H^n(M)$ 定义为相应 $R$-模复形的
上同调。因此，微分分次模的短正合列 $0 \to K \to L \to M \to 0$
产生长正合列
\begin{equation}
\label{dga-equation-les}
H^n(K) \to H^n(L) \to H^n(M) \to H^{n + 1}(K)
\end{equation}
；它由上同调模组成，见《同调代数》引理
\ref{homology-lemma-long-exact-sequence-cochain}。

\medskip\noindent
此外，从现在起，我们沿用模复形的全部术语。例如，若对所有
$k \in \mathbf{Z}$ 都有 $H^k(M) = 0$，则称微分分次 $A$-模 $M$
是 {\it 无环的}。若微分分次 $A$-模的同态 $M \to N$ 对所有
$k \in \mathbf{Z}$ 都诱导同构 $H^k(M) \to H^k(N)$，则称它为
{\it 拟同构}；其余术语依此类推。

\begin{definition}
\label{dga-definition-shift}
设 $(A, \text{d})$ 是微分分次代数。设 $M$ 是微分分次模，
其底层 $R$-模复形为 $M^\bullet$。对任意 $k \in \mathbf{Z}$，
如下定义 {\it $k$-移位模} $M[k]$：
\begin{enumerate}
\item $M[k]$ 的底层 $R$-模复形是 $M^\bullet[k]$，即
$M[k]^n = M^{n + k}$ 且 $\text{d}_{M[k]} = (-1)^k\text{d}_M$；以及
\item 作为 $A$-模，其乘法
$$
(M[k])^n \times A^m \longrightarrow (M[k])^{n + m}
$$
等于给定的乘法 $M^{n + k} \times A^m \to M^{n + k + m}$。
\end{enumerate}
对于微分分次 $A$-模的态射 $f : M \to N$，令
$f[k] : M[k] \to N[k]$ 为在底层 $A$-模上等于 $f$ 的映射。
这定义了函子
$[k] : \text{Mod}_{(A, \text{d})} \to \text{Mod}_{(A, \text{d})}$.
\end{definition}

\noindent
来验证这个选择满足 Leibniz 规则。设 $x \in M[k]^n = M^{n + k}$、
$a \in A^m$，并以 $\cdot_{M[k]}$ 表示 $M[k]$ 中的乘法，则
\begin{align*}
\text{d}_{M[k]}(x \cdot_{M[k]} a)
& =
(-1)^k \text{d}_M(xa) \\
& =
(-1)^k \text{d}_M(x) a + (-1)^{k + n + k} x \text{d}(a) \\
& =
\text{d}_{M[k]}(x) a + (-1)^n x \text{d}(a) \\
& =
\text{d}_{M[k]}(x) \cdot_{M[k]} a + (-1)^n x \cdot_{M[k]} \text{d}(a)
\end{align*}
这正是所需结论，因为作为 $M[k]$ 的齐次元素，$x$ 的次数为 $n$。
还注意到，在这些选择下，可把乘法映射看作复形映射
$$
\text{Tot}(M^\bullet[k] \otimes _R A^\bullet) \to
\text{Tot}(M^\bullet \otimes _R A^\bullet)[k] \to
M^\bullet[k]
$$
其中第一个箭头是《更多代数》第 \ref{more-algebra-section-sign-rules} 节
(\ref{more-algebra-item-shift-tensor}) 中的箭头；在本情形下它不涉及符号。
（事实上，也可由这一观察推出 Leibniz 规则成立。）

\medskip\noindent
《同调代数》第 \ref{homology-section-homotopy-shift} 节中的说明适用。
特别地，我们将在不引入符号的情况下认同所有移位 $M[k]$ 的上同调群。

\medskip\noindent
至此，我们已有足够结构来讨论 {\it 三角}；见《导出范畴》定义
\ref{derived-definition-triangle}。下一步将发展充分的理论，以陈述并证明
微分分次模的同伦范畴是三角范畴。我们先定义同伦范畴。






\section{同伦范畴}
\label{dga-section-homotopy}

\noindent
我们的同伦考虑 $A$-模结构和分次，但（当然）不要求与微分相容。

\begin{definition}
\label{dga-definition-homotopy}
设 $(A, \text{d})$ 是微分分次代数。设 $f, g : M \to N$ 是微分分次
$A$-模的同态。$f$ 与 $g$ 之间的一个 {\it 同伦}是一个 $A$-模映射
$h : M \to N$，使得
\begin{enumerate}
\item 对所有 $n$ 有 $h(M^n) \subset N^{n - 1}$；并且
\item 对所有 $x \in M$ 有
$f(x) - g(x) = \text{d}_N(h(x)) + h(\text{d}_M(x))$。
\end{enumerate}
若存在这样的同伦，则称 $f$ 与 $g$ {\it 同伦}。
\end{definition}

\noindent
因此，$h$ 与 $A$-模结构及分次相容，但不与微分相容。若 $f = g$，
且 $h$ 是定义中的同伦，则 $h$ 在 $\text{Mod}_{(A, \text{d})}$ 中
定义一个态射 $h : M \to N[-1]$。

\begin{lemma}
\label{dga-lemma-compose-homotopy}
设 $(A, \text{d})$ 是微分分次代数。设 $f, g : L \to M$ 是微分分次
$A$-模的同态，并另给同态 $a : K \to L$ 和 $c : M \to N$。
若 $A$-模映射 $h : L \to M$ 定义了 $f$ 与 $g$ 之间的同伦，
则 $c \circ h \circ a$ 定义了 $c \circ f \circ a$ 与
$c \circ g \circ a$ 之间的同伦。
\end{lemma}

\begin{proof}
立即由《同调代数》引理 \ref{homology-lemma-compose-homotopy-cochain} 得出。
\end{proof}

\noindent
这个引理使我们可以如下定义同伦范畴。

\begin{definition}
\label{dga-definition-complexes-notation}
设 $(A, \text{d})$ 是微分分次代数。{\it 同伦范畴}记为
$K(\text{Mod}_{(A, \text{d})})$；其对象是
$\text{Mod}_{(A, \text{d})}$ 的对象，其态射是微分分次 $A$-模同态的
同伦类。
\end{definition}

\noindent
记号 $K(\text{Mod}_{(A, \text{d})})$ 并非标准记号，但至少与
Stacks project 其他部分中 $K(-)$ 的用法一致。

\begin{lemma}
\label{dga-lemma-homotopy-direct-sums}
设 $(A, \text{d})$ 是微分分次代数。同伦范畴
$K(\text{Mod}_{(A, \text{d})})$ 具有直和和乘积。
\end{lemma}

\begin{proof}
略。提示：直接使用引理 \ref{dga-lemma-dgm-abelian} 中的直和和乘积。
这是可行的，因为我们已经看到这些函子与到分次 $A$-模范畴的遗忘函子
可交换，并且 $\prod$ 在阿贝尔群族所成的范畴上是正合函子。
\end{proof}







\section{锥}
\label{dga-section-cones}

\noindent
我们为微分分次模范畴引入锥。

\begin{definition}
\label{dga-definition-cone}
设 $(A, \text{d})$ 是微分分次代数。设 $f : K \to L$ 是微分分次
$A$-模的同态。$f$ 的 {\it 锥}是微分分次 $A$-模 $C(f)$；
它由 $C(f) = L \oplus K$ 给出，其分次为
$C(f)^n = L^n \oplus K^{n + 1}$，微分为
$$
d_{C(f)} =
\left(
\begin{matrix}
\text{d}_L & f \\
0 & -\text{d}_K
\end{matrix}
\right)
$$
它配备由显然映射 $L \to C(f)$ 和 $C(f) \to K$ 所诱导的典范复形态射
$i : L \to C(f)$ 和 $p : C(f) \to K[1]$。
\end{definition}

\noindent
锥三角的构造在下述意义下是函子的。

\begin{lemma}
\label{dga-lemma-functorial-cone}
设 $(A, \text{d})$ 是微分分次代数。假设
$$
\xymatrix{
K_1 \ar[r]_{f_1} \ar[d]_a & L_1 \ar[d]^b \\
K_2 \ar[r]^{f_2} & L_2
}
$$
是一个在相差同伦意义下交换的微分分次 $A$-模同态图表。
那么存在态射 $c : C(f_1) \to C(f_2)$，它产生三角态射
$$
(a, b, c) : (K_1, L_1, C(f_1), f_1, i_1, p_1) \to
(K_1, L_1, C(f_1), f_2, i_2, p_2)
$$
，后者位于 $K(\text{Mod}_{(A, \text{d})})$ 中。
\end{lemma}

\begin{proof}
设 $h : K_1 \to L_2$ 是 $f_2 \circ a$ 与 $b \circ f_1$ 之间的同伦。
用矩阵
$$
c =
\left(
\begin{matrix}
b & h \\
0 & a
\end{matrix}
\right) :
L_1 \oplus K_1 \to L_2 \oplus K_2
$$
定义 $c$。矩阵计算表明 $c$ 是微分分次模的态射。
显然 $c \circ i_1 = i_2 \circ b$，也容易验证
$p_2 \circ c = a \circ p_1$。
\end{proof}











\section{容许短正合列}
\label{dga-section-admissible}

\noindent
容许短正合列是逐项分裂正合列在微分分次模框架中的对应物。

\begin{definition}
\label{dga-definition-admissible-ses}
设 $(A, \text{d})$ 是微分分次代数。
\begin{enumerate}
\item 若存在分次 $A$-模映射 $L \to K$，它是 $K \to L$ 的左逆，
则称微分分次 $A$-模的同态 $K \to L$ 为 {\it 容许单态射}。
\item 若存在分次 $A$-模映射 $M \to L$，它是 $L \to M$ 的右逆，
则称微分分次 $A$-模的同态 $L \to M$ 为 {\it 容许满态射}。
\item 若微分分次 $A$-模的短正合列 $0 \to K \to L \to M \to 0$
作为分次 $A$-模列分裂，则称它为 {\it 容许短正合列}。
\end{enumerate}
\end{definition}

\noindent
因此，这些分裂与除微分外的所有数据相容。给定一个容许短正合列，
我们得到一个三角；这正是要求分裂与 $A$-模结构相容的原因。

\begin{lemma}
\label{dga-lemma-admissible-ses}
设 $(A, \text{d})$ 是微分分次代数。设
$0 \to K \to L \to M \to 0$ 是微分分次 $A$-模的容许短正合列。
设 $s : M \to L$ 和 $\pi : L \to K$ 是满足
$\Ker(\pi) = \Im(s)$ 的分裂。那么得到态射
$$
\delta = \pi \circ \text{d}_L \circ s : M \to K[1]
$$
，它属于 $\text{Mod}_{(A, \text{d})}$，并诱导上同调长正合列
(\ref{dga-equation-les}) 中的边界映射。
\end{lemma}

\begin{proof}
由构造，映射 $\pi \circ \text{d}_L \circ s$ 与 $A$-模结构及分次相容。
由《同调代数》引理
\ref{homology-lemma-ses-termwise-split-cochain}，它也与微分相容。
设 $R$ 是使 $A$ 成为其上微分分次代数的环。映射的等式是关于
$R$-模的陈述。因此，结论由《同调代数》引理
\ref{homology-lemma-ses-termwise-split-cochain} 和
\ref{homology-lemma-ses-termwise-split-long-cochain} 得出。
\end{proof}

\begin{lemma}
\label{dga-lemma-make-commute-map}
设 $(A, \text{d})$ 是微分分次代数。设
$$
\xymatrix{
K \ar[r]_f \ar[d]_a & L \ar[d]^b \\
M \ar[r]^g & N
}
$$
是一个相差同伦交换的微分分次 $A$-模同态图表。
\begin{enumerate}
\item 若 $f$ 是容许单态射，则 $b$ 同伦于一个使图表交换的同态。
\item 若 $g$ 是容许满态射，则 $a$ 同伦于一个使图表交换的态射。
\end{enumerate}
\end{lemma}

\begin{proof}
设 $h : K \to N$ 是 $bf$ 与 $ga$ 之间的同伦，即
$bf - ga = \text{d}h + h\text{d}$。假设分次 $A$-模映射
$\pi : L \to K$ 是 $f$ 的左逆。取
$b' = b - \text{d}h\pi - h\pi \text{d}$。
假设分次 $A$-模映射 $s : N \to M$ 是 $g$ 的右逆。
取 $a' = a + \text{d}sh + sh\text{d}$。计算略。
\end{proof}

\begin{lemma}
\label{dga-lemma-make-injective}
设 $(A, \text{d})$ 是微分分次代数。设 $\alpha : K \to L$ 是微分分次
$A$-模的同态。存在分解
$$
\xymatrix{
K \ar[r]^{\tilde \alpha} \ar@/_1pc/[rr]_\alpha &
\tilde L \ar[r]^\pi & L
}
$$
，它位于 $\text{Mod}_{(A, \text{d})}$ 中，并且
\begin{enumerate}
\item $\tilde \alpha$ 是容许单态射（见定义
\ref{dga-definition-admissible-ses}）；
\item 存在态射 $s : L \to \tilde L$，使得
$\pi \circ s = \text{id}_L$，并且 $s \circ \pi$ 同伦于
$\text{id}_{\tilde L}$。
\end{enumerate}
\end{lemma}

\begin{proof}
证明与《导出范畴》引理 \ref{derived-lemma-make-injective} 的证明相同。
具体而言，置 $\tilde L = L \oplus C(1_K)$，并使用锥构造的初等性质。
\end{proof}

\begin{lemma}
\label{dga-lemma-sequence-maps-split}
设 $(A, \text{d})$ 是微分分次代数。设
$L_1 \to L_2 \to \ldots \to L_n$ 是微分分次 $A$-模的可复合同态列。
存在交换图表
$$
\xymatrix{
L_1 \ar[r] &
L_2 \ar[r] &
\ldots \ar[r] &
L_n \\
M_1 \ar[r] \ar[u] &
M_2 \ar[r] \ar[u] &
\ldots \ar[r] &
M_n \ar[u]
}
$$
，它位于 $\text{Mod}_{(A, \text{d})}$ 中，使得每个
$M_i \to M_{i + 1}$ 都是容许单态射，且每个 $M_i \to L_i$
都是同伦等价。
\end{lemma}

\begin{proof}
$n = 1$ 的情形没有内容。引理 \ref{dga-lemma-make-injective} 是 $n = 2$
的情形。假设除 $M_n$ 外图表均已构造。对复合
$M_{n - 1} \to L_{n - 1} \to L_n$ 应用引理
\ref{dga-lemma-make-injective}，所得分解
$M_{n - 1} \to M_n \to L_n$ 即满足要求。
\end{proof}



\begin{lemma}
\label{dga-lemma-nilpotent}
设 $(A, \text{d})$ 是微分分次代数。设
$0 \to K_i \to L_i \to M_i \to 0$（$i = 1, 2, 3$）是微分分次
$A$-模的容许短正合列。设 $b : L_1 \to L_2$ 和
$b' : L_2 \to L_3$ 是微分分次模的同态，使得
$$
\vcenter{
\xymatrix{
K_1 \ar[d]_0 \ar[r] &
L_1 \ar[r] \ar[d]_b &
M_1 \ar[d]_0 \\
K_2 \ar[r] & L_2 \ar[r] & M_2
}
}
\quad\text{且}\quad
\vcenter{
\xymatrix{
K_2 \ar[d]^0 \ar[r] &
L_2 \ar[r] \ar[d]^{b'} &
M_2 \ar[d]^0 \\
K_3 \ar[r] & L_3 \ar[r] & M_3
}
}
$$
相差同伦交换。那么 $b' \circ b$ 同伦于 $0$。
\end{lemma}

\begin{proof}
由引理 \ref{dga-lemma-make-commute-map}，可把 $b$ 和 $b'$ 替换为与之同伦的
映射，使得左图的右方块交换，而右图的左方块交换。换言之，
$\Im(b) \subset \Im(K_2 \to L_2)$ 且
$\Ker(b') \supset \Im(K_2 \to L_2)$。于是，作为模映射，
$b' \circ b = 0$。
\end{proof}
















\section{有区别三角}
\label{dga-section-distinguished}

\noindent
下述引理产生我们所用的有区别三角。

\begin{lemma}
\label{dga-lemma-triangle-independent-splittings}
设 $(A, \text{d})$ 是微分分次代数。设
$0 \to K \to L \to M \to 0$ 是微分分次 $A$-模的容许短正合列。
三角
\begin{equation}
\label{dga-equation-triangle-associated-to-admissible-ses}
K \to L \to M \xrightarrow{\delta} K[1]
\end{equation}
中的 $\delta$ 如引理 \ref{dga-lemma-admissible-ses} 所述。该三角在
$K(\text{Mod}_{(A, \text{d})})$ 中相差典范同构，不依赖于引理
\ref{dga-lemma-admissible-ses} 中所作的选择。
\end{lemma}

\begin{proof}
具体地，设 $(s', \pi')$ 是引理 \ref{dga-lemma-admissible-ses} 中那种分裂的
第二个选择。我们断言 $\delta$ 与 $\delta'$ 同伦。事实上，对某些唯一的
次数为 $-1$ 的 $A$-模同态 $h : M \to K$ 和 $g : M \to K$，写成
$s' = s + \alpha \circ h$ 和 $\pi' = \pi + g \circ \beta$。
则 $g = -h$，且 $g$ 是 $\delta$ 与 $\delta'$ 之间的同伦。
相关计算见《同调代数》引理
\ref{homology-lemma-ses-termwise-split-homotopy-cochain} 的证明。
\end{proof}

\begin{definition}
\label{dga-definition-distinguished-triangle}
设 $(A, \text{d})$ 是微分分次代数。
\begin{enumerate}
\item 若 $0 \to K \to L \to M \to 0$ 是微分分次 $A$-模的容许短正合列，
则 {\it 与 $0 \to K \to L \to M \to 0$ 关联的三角}是
$K(\text{Mod}_{(A, \text{d})})$ 中的三角
(\ref{dga-equation-triangle-associated-to-admissible-ses})。
\item 若 $K(\text{Mod}_{(A, \text{d})})$ 中的三角同构于某个与微分分次
$A$-模的容许短正合列关联的三角，则称它为 {\it 有区别三角}。
\end{enumerate}
\end{definition}









\section{锥与有区别三角}
\label{dga-section-cones-and-triangles}

\noindent
设 $(A, \text{d})$ 是微分分次代数。设 $f : K \to L$ 是微分分次
$A$-模的同态。那么 $(K, L, C(f), f, i, p)$ 构成三角
$$
K \to L \to C(f) \to K[1]
$$
，它位于 $\text{Mod}_{(A, \text{d})}$ 中，因而也位于
$K(\text{Mod}_{(A, \text{d})})$ 中。一般而言，锥{\bf 并不}是有区别三角，
但两者只相差一个符号或一次旋转（任君选择）。下面给出两个精确陈述。

\begin{lemma}
\label{dga-lemma-rotate-cone}
设 $(A, \text{d})$ 是微分分次代数。设 $f : K \to L$ 是微分分次模的同态。
三角 $(L, C(f), K[1], i, p, f[1])$ 是与容许短正合列
$$
0 \to L \to C(f) \to K[1] \to 0
$$
关联的三角；该短正合列来自 $f$ 的锥的定义。
\end{lemma}

\begin{proof}
立即由定义得出。
\end{proof}

\begin{lemma}
\label{dga-lemma-rotate-triangle}
设 $(A, \text{d})$ 是微分分次代数。设 $\alpha : K \to L$ 和
$\beta : L \to M$ 定义容许短正合列
$$
0 \to K \to L \to M \to 0
$$
，它由微分分次 $A$-模组成。设 $(K, L, M, \alpha, \beta, \delta)$
为关联三角。那么三角
$$
(M[-1], K, L, \delta[-1], \alpha, \beta)
\quad\text{且}\quad
(M[-1], K, C(\delta[-1]), \delta[-1], i, p)
$$
同构。
\end{lemma}

\begin{proof}
选取分裂后，写 $L = K \oplus M$，并把 $\alpha$ 和 $\beta$ 分别认同为
自然包含映射和投影映射。由 $\delta$ 的构造，有
$$
d_B =
\left(
\begin{matrix}
d_K & \delta \\
0 & d_M
\end{matrix}
\right)
$$
另一方面，$\delta[-1] : M[-1] \to K$ 的锥由
$C(\delta[-1]) = K \oplus M$ 给出，其微分与上面的矩阵完全相同！
引理由此得证。
\end{proof}

\begin{lemma}
\label{dga-lemma-third-isomorphism}
设 $(A, \text{d})$ 是微分分次代数。设 $f_1 : K_1 \to L_1$ 和
$f_2 : K_2 \to L_2$ 是微分分次 $A$-模的同态。设
$$
(a, b, c) :
(K_1, L_1, C(f_1), f_1, i_1, p_1)
\longrightarrow
(K_1, L_1, C(f_1), f_2, i_2, p_2)
$$
为 $K(\text{Mod}_{(A, \text{d})})$ 中的任意三角态射。
若 $a$ 和 $b$ 是同伦等价，则 $c$ 也是同伦等价。
\end{lemma}

\begin{proof}
设 $a^{-1} : K_2 \to K_1$ 是微分分次 $A$-模的同态，且它在
$K(\text{Mod}_{(A, \text{d})})$ 中是 $a$ 的逆。设
$b^{-1} : L_2 \to L_1$ 是微分分次 $A$-模的同态，且它在该范畴中是
$b$ 的逆。对 $f_1 \circ a^{-1} = b^{-1} \circ f_2$ 应用引理
\ref{dga-lemma-functorial-cone}，得到态射 $c' : C(f_2) \to C(f_1)$。
若能证明 $c \circ c'$ 和 $c' \circ c$ 在
$K(\text{Mod}_{(A, \text{d})})$ 中都是同构，结论即得。因此，只需证明：
给定 $K(\text{Mod}_{(A, \text{d})})$ 中形如
$(1, 1, c) : (K, L, C(f), f, i, p)$ 的三角态射，态射 $c$ 是该范畴中的
同构。由假设，下图中的两个方块
$$
\xymatrix{
L \ar[r] \ar[d]_1 &
C(f) \ar[r] \ar[d]_c &
K[1] \ar[d]_1 \\
L \ar[r] &
C(f) \ar[r] &
K[1]
}
$$
相差同伦交换。由 $C(f)$ 的构造，两行都是容许短正合列。
于是，由引理 \ref{dga-lemma-nilpotent}，在
$K(\text{Mod}_{(A, \text{d})})$ 中有 $(c - 1)^2 = 0$。
因此，$c$ 是该范畴中的同构，其逆为 $2 - c$。
\end{proof}

\noindent
下述引理表明，同伦范畴中由锥给出的三角族与有区别三角族相差同构而相同，
至少在不计符号时如此！

\begin{lemma}
\label{dga-lemma-the-same-up-to-isomorphisms}
设 $(A, \text{d})$ 是微分分次代数。
\begin{enumerate}
\item 给定微分分次 $A$-模的容许短正合列
$0 \to K \xrightarrow{\alpha} L \to M \to 0$
，存在同伦等价 $C(\alpha) \to M$，使得图表
$$
\xymatrix{
K \ar[r] \ar[d] & L \ar[d] \ar[r] &
C(\alpha) \ar[r]_{-p} \ar[d] & K[1] \ar[d] \\
K \ar[r]^\alpha & L \ar[r]^\beta &
M \ar[r]^\delta & K[1]
}
$$
在 $K(\text{Mod}_{(A, \text{d})})$ 中定义一个三角同构。
\item 给定复形态射 $f : K \to L$，存在三角同构
$$
\xymatrix{
K \ar[r] \ar[d] & \tilde L \ar[d] \ar[r] &
M \ar[r]_{\delta} \ar[d] & K[1] \ar[d] \\
K \ar[r] & L \ar[r] &
C(f) \ar[r]^{-p} & K[1]
}
$$
，其中上方三角是与容许短正合列 $K \to \tilde L \to M$ 关联的三角。
\end{enumerate}
\end{lemma}

\begin{proof}
证明 (1)。有 $C(\alpha) = L \oplus K$；我们简单地把
$C(\alpha) \to M$ 定义为先投影到 $L$，再与 $\beta$ 复合。
由于复合 $K^{n + 1} \to L^{n + 1} \to M^{n + 1}$ 为零，这定义了
微分分次模的态射。选择满足 $\Ker(\pi) = \Im(s)$ 的分裂
$s : M \to L$ 和 $\pi : L \to K$，并照例置
$\delta = \pi \circ \text{d}_L \circ s$。取由 $(s , -\delta)$ 给出的
$M \to C(\alpha)$ 作为同伦逆。它与微分相容，因为 $\delta^n$
可刻画为唯一满足下式的映射 $M^n \to K^{n + 1}$：
$\text{d} \circ s^n - s^{n + 1} \circ \text{d} = \alpha \circ \delta^n$,
见《同调代数》引理 \ref{homology-lemma-ses-termwise-split-cochain} 的证明。
复合 $M \to C(f) \to M$ 是恒等映射。复合
$C(f) \to M \to C(f)$ 等于态射
$$
\left(
\begin{matrix}
s \circ \beta & 0 \\
-\delta \circ \beta & 0
\end{matrix}
\right)
$$
为看出它同伦于恒等映射，使用由矩阵
$$
\left(
\begin{matrix}
0 & 0 \\
\pi & 0
\end{matrix}
\right) :
C(\alpha) = L \oplus K
\to
L \oplus K = C(\alpha)
$$
给出的同伦 $h : C(\alpha) \to C(\alpha)$。容易验证
$$
\left(
\begin{matrix}
1 & 0 \\
0 & 1
\end{matrix}
\right)
-
\left(
\begin{matrix}
s \\
-\delta
\end{matrix}
\right)
\left(
\begin{matrix}
\beta & 0
\end{matrix}
\right)
=
\left(
\begin{matrix}
\text{d} & \alpha \\
0 & -\text{d}
\end{matrix}
\right)
\left(
\begin{matrix}
0 & 0 \\
\pi & 0
\end{matrix}
\right)
+
\left(
\begin{matrix}
0 & 0 \\
\pi & 0
\end{matrix}
\right)
\left(
\begin{matrix}
\text{d} & \alpha \\
0 & -\text{d}
\end{matrix}
\right)
$$
为完成 (1) 的证明，还须证明态射 $-p : C(\alpha) \to K[1]$
（见定义 \ref{dga-definition-cone}）与复合 $C(\alpha) \to M \to K[1]$
相差同伦相等。这由上文显然可见。事实上，可使用同伦逆
$(s, -\delta) : M \to C(\alpha)$，转而验证两个映射
$M \to K[1]$ 相等。注意 $p \circ (s, -\delta) = -\delta$，正合所需。

\medskip\noindent
证明 (2)。令 $\tilde f : K \to \tilde L$、$s : L \to \tilde L$
和 $\pi : L \to L$ 如引理 \ref{dga-lemma-make-injective} 中所述。
由引理 \ref{dga-lemma-functorial-cone} 和 \ref{dga-lemma-third-isomorphism}，
三角 $(K, L, C(f), i, p)$ 与
$(K, \tilde L, C(\tilde f), \tilde i, \tilde p)$ 同构。
注意，可以复合三角的同构。因此，可以用 $\tilde L$ 替换 $L$，
并用 $\tilde f$ 替换 $f$。换言之，可以假设 $f$ 是容许单态射。
此时结论由 (1) 得出。
\end{proof}







\section{同伦范畴是三角范畴}
\label{dga-section-homotopy-triangulated}

\noindent
我们先证明它是预三角范畴。

\begin{lemma}
\label{dga-lemma-homotopy-category-pre-triangulated}
设 $(A, \text{d})$ 是微分分次代数。同伦范畴
$K(\text{Mod}_{(A, \text{d})})$ 连同其自然平移函子和有区别三角，
是一个预三角范畴。
\end{lemma}

\begin{proof}
证明 TR1。由定义，任何同构于有区别三角的三角都是有区别三角。
此外，任意三角 $(K, K, 0, 1, 0, 0)$ 都是有区别三角，因为
$0 \to K \to K \to 0 \to 0$ 是容许短正合列。最后，给定微分分次
$A$-模的任意同态 $f : K \to L$，由引理
\ref{dga-lemma-the-same-up-to-isomorphisms}，三角
$(K, L, C(f), f, i, -p)$ 是有区别三角。

\medskip\noindent
证明 TR2。设 $(X, Y, Z, f, g, h)$ 是一个三角，并假设
$(Y, Z, X[1], g, h, -f[1])$ 是有区别三角。则存在容许短正合列
$0 \to K \to L \to M \to 0$，使其关联三角
$(K, L, M, \alpha, \beta, \delta)$ 同构于
$(Y, Z, X[1], g, h, -f[1])$。反向旋转可见，
$(X, Y, Z, f, g, h)$ 同构于
$(M[-1], K, L, -\delta[-1], \alpha, \beta)$。
由引理 \ref{dga-lemma-rotate-triangle}，三角
$(M[-1], K, L, \delta[-1], \alpha, \beta)$ 同构于
$(M[-1], K, C(\delta[-1]), \delta[-1], i, p)$。
把前一三角同构预复合以 $Y$ 上的 $-1$，可得
$(X, Y, Z, f, g, h)$ 同构于
$(M[-1], K, C(\delta[-1]), \delta[-1], i, -p)$。
因此，由引理 \ref{dga-lemma-the-same-up-to-isomorphisms}，它是有区别三角。
反过来，假设 $(X, Y, Z, f, g, h)$ 是有区别三角。
由引理 \ref{dga-lemma-the-same-up-to-isomorphisms}，这意味着它同构于某个
$\text{Mod}_{(A, \text{d})}$ 中的态射 $f$ 所给出的三角
$(K, L, C(f), f, i, -p)$。于是旋转后的三角
$(Y, Z, X[1], g, h, -f[1])$ 同构于
$(L, C(f), K[1], i, -p, -f[1])$，后者又同构于三角
$(L, C(f), K[1], i, p, f[1])$。由引理 \ref{dga-lemma-rotate-cone}，
这个三角是有区别三角。因此，$(Y, Z, X[1], g, h, -f[1])$
如所需是有区别三角。

\medskip\noindent
证明 TR3。设 $(X, Y, Z, f, g, h)$ 和 $(X', Y', Z', f', g', h')$
是 $K(\mathcal{A})$ 的有区别三角，并设 $a : X \to X'$ 和
$b : Y \to Y'$ 是满足 $f' \circ a = b \circ f$ 的态射。
由引理 \ref{dga-lemma-the-same-up-to-isomorphisms}，可以假设
$(X, Y, Z, f, g, h) = (X, Y, C(f), f, i, -p)$，并且
$(X', Y', Z', f', g', h') = (X', Y', C(f'), f', i', -p')$。
此时，只需对 $f, f', a, b$ 给出的交换图表应用引理
\ref{dga-lemma-functorial-cone}。
\end{proof}

\noindent
在一般地证明 TR4 之前，先证明一个特殊情形。

\begin{lemma}
\label{dga-lemma-two-split-injections}
设 $(A, \text{d})$ 是微分分次代数。假设 $\alpha : K \to L$ 和
$\beta : L \to M$ 是微分分次 $A$-模的容许单态射。那么存在有区别三角
$(K, L, Q_1, \alpha, p_1, d_1)$、
$(K, M, Q_2, \beta \circ \alpha, p_2, d_2)$ 和
$(L, M, Q_3, \beta, p_3, d_3)$，使得 TR4 对它们成立。
\end{lemma}

\begin{proof}
设 $\pi_1 : L \to K$ 和 $\pi_3 : M \to L$ 是分次 $A$-模同态，
分别为 $\alpha$ 和 $\beta$ 的左逆。于是 $K \to M$ 也是容许单态射，
其左逆为 $\pi_2 = \pi_1 \circ \pi_3$。用 $Q_1$、$Q_2$、$Q_3$
分别表示 $K \to L$、$K \to M$、$L \to M$ 的余核。
则作为分次 $A$-模，有认同
$Q_1 = \Ker(\pi_1)$、$Q_3 = \Ker(\pi_3)$ 和
$Q_2 = \Ker(\pi_2)$。又有 $L = K \oplus Q_1$ 和
$M = L \oplus Q_3$，故 $M = K \oplus Q_1 \oplus Q_3$。
注意，$\pi_2 = \pi_1 \circ \pi_3$ 在 $Q_1$ 和 $Q_3$ 上均为零。
因此 $Q_2 = Q_1 \oplus Q_3$。考虑交换图表
$$
\begin{matrix}
0 & \to & K & \to & L & \to & Q_1 & \to & 0 \\
  &     & \downarrow &     & \downarrow &     & \downarrow  & \\
0 & \to & K & \to & M & \to & Q_2 & \to & 0 \\
  &     & \downarrow &     & \downarrow &     & \downarrow  & \\
0 & \to & L & \to & M & \to & Q_3 & \to & 0
\end{matrix}
$$
该图表的各行都是容许短正合列，因而按定义确定有区别三角。此外，
上图中的向下箭头与所选分裂相容，因而定义三角态射
$$
(K \to L \to Q_1 \to K[1])
\longrightarrow
(K \to M \to Q_2 \to K[1])
$$
以及
$$
(K \to M \to Q_2 \to K[1])
\longrightarrow
(L \to M \to Q_3 \to L[1]).
$$
注意，把图表中最下方列的分裂 $Q_3 \to M$ 与 $M \to Q_2$ 复合，
就得到分裂列 $0 \to Q_1 \to Q_2 \to Q_3 \to 0$ 的一个分裂。
由此容易看出，相应有区别三角
$$
(Q_1 \to Q_2 \to Q_3 \to Q_1[1])
$$
中的态射 $\delta : Q_3 \to Q_1[1]$ 等于复合
$Q_3 \to L[1] \to Q_1[1]$。因此，我们得到公理 TR4 结论中的结构。
\end{proof}

\noindent
下面给出最终结论。

\begin{proposition}
\label{dga-proposition-homotopy-category-triangulated}
设 $(A, \text{d})$ 是微分分次代数。微分分次 $A$-模的同伦范畴
$K(\text{Mod}_{(A, \text{d})})$ 连同其自然平移函子和有区别三角，
是一个三角范畴。
\end{proposition}

\begin{proof}
已知 $K(\text{Mod}_{(A, \text{d})})$ 是预三角范畴。因此，只需证明 TR4，
而证明它可使用《导出范畴》引理
\ref{derived-lemma-easier-axiom-four}。设 $K \to L$ 和 $L \to M$
是 $K(\text{Mod}_{(A, \text{d})})$ 中的可复合态射。由引理
\ref{dga-lemma-sequence-maps-split}，可以假设 $K \to L$ 和 $L \to M$
都是容许单态射。此时结论由引理 \ref{dga-lemma-two-split-injections} 得出。
\end{proof}











\section{左模}
\label{dga-section-left-modules}

\noindent
迄今所述的一切在左微分分次模的框架中都有对应版本，
但必须留意某些符号规则。

\medskip\noindent
设 $(A, \text{d})$ 是微分分次 $R$-代数。完全类似于右模，
把一个 {\it 左微分分次 $A$-模} $M$ 定义为一个左 $A$-模 $M$，
它配备分次 $M = \bigoplus M^n$ 和微分 $\text{d}$，使得
$A^n M^m \subset M^{n + m}$、$\text{d}(M^n) \subset M^{n + 1}$，并且
$$
\text{d}(am) = \text{d}(a) m + (-1)^{\deg(a)}a \text{d}(m)
$$
对齐次元素 $a \in A$ 和 $m \in M$ 成立。如前，这一 Leibniz 规则
恰好意味着乘法定义复形映射
$$
\text{Tot}(A^\bullet \otimes_R M^\bullet) \to M^\bullet
$$
这里 $A^\bullet$ 和 $M^\bullet$ 表示 $A$ 和 $M$ 的底层 $R$-模复形。

\begin{definition}
\label{dga-definition-opposite-dga}
设 $R$ 是环。设 $(A, \text{d})$ 是 $R$ 上的微分分次代数。
{\it 反微分分次代数}是 $R$ 上的微分分次代数
$(A^{opp}, \text{d})$；作为分次 $R$-模，有 $A^{opp} = A$，
微分满足 $\text{d} = \text{d}$，而乘法由
$$
a \cdot_{opp} b = (-1)^{\deg(a)\deg(b)} b a
$$
给出，其中 $a, b \in A$ 为齐次元素。
\end{definition}

\noindent
该定义成立，因为
\begin{align*}
\text{d}(a \cdot_{opp} b)
& =
(-1)^{\deg(a)\deg(b)} \text{d}(b a) \\
& =
(-1)^{\deg(a)\deg(b)} \text{d}(b) a +
(-1)^{\deg(a)\deg(b) + \deg(b)}b\text{d}(a) \\
& =
(-1)^{\deg(a)}a \cdot_{opp} \text{d}(b) + \text{d}(a) \cdot_{opp} b
\end{align*}
，正合所需。用底层 $R$-模复形表述，这意味着图表
$$
\xymatrix{
\text{Tot}(A^\bullet \otimes_R A^\bullet)
\ar[rrr]_-{A^{opp}\text{ 的乘法}}
\ar[d]_{\text{交换性约束}} & & &
A^\bullet \ar[d]^{\text{id}} \\
\text{Tot}(A^\bullet \otimes_R A^\bullet)
\ar[rrr]^-{A\text{ 的乘法}} & & &
A^\bullet
}
$$
交换。这里，$R$-模复形的对称幺半范畴上的交换性约束见《更多代数》
第 \ref{more-algebra-section-sign-rules} 节。

\medskip\noindent
设 $(A, \text{d})$ 是 $R$ 上的微分分次代数。设 $M$ 是左微分分次
$A$-模。把 $M$ 视为右 $A^{opp}$-模时记为 $M^{opp}$，其乘法
$\cdot_{opp}$ 由规则
$$
m \cdot_{opp} a = (-1)^{\deg(a)\deg(m)} a m
$$
定义，其中 $a$ 和 $m$ 为齐次元素。它与微分相容，因为也可以用图表
$$
\xymatrix{
\text{Tot}(M^\bullet \otimes_R A^\bullet)
\ar[rrr]_-{M^{opp}\text{ 上的乘法}}
\ar[d]_{\text{交换性约束}} & & &
M^\bullet \ar[d]^{\text{id}} \\
\text{Tot}(A^\bullet \otimes_R M^\bullet)
\ar[rrr]^-{M\text{ 上的乘法}} & & &
M^\bullet
}
$$
来定义 $M^{opp}$ 上的乘法 $\cdot_{opp}$。为看出它是结合乘法，
对齐次元素 $a, b \in A$ 和 $m \in M$ 计算如下：
\begin{align*}
m \cdot_{opp} (a \cdot_{opp} b)
& =
(-1)^{\deg(a)\deg(b)} m \cdot_{opp} (ba) \\
& =
(-1)^{\deg(a)\deg(b) + \deg(ab)\deg(m)} bam \\
& =
(-1)^{\deg(a)\deg(b) + \deg(ab)\deg(m) + \deg(b)\deg(am)}
(am) \cdot_{opp} b \\
& =
(-1)^{\deg(a)\deg(b) + \deg(ab)\deg(m) + \deg(b)\deg(am) + \deg(a)\deg(m)}
(m \cdot_{opp} a) \cdot_{opp} b \\
& =
(m \cdot_{opp} a) \cdot_{opp} b
\end{align*}
当然，也可利用 $R$-模复形的对称幺半范畴上结合性约束与交换性约束之间的
相容性证明这一点。

\begin{lemma}
\label{dga-lemma-left-right}
设 $(A, \text{d})$ 是微分分次 $R$-代数。从左微分分次 $A$-模范畴到
右微分分次 $A^{opp}$-模范畴的函子 $M \mapsto M^{opp}$ 是等价。
\end{lemma}

\begin{proof}
略。
\end{proof}

\noindent
接着讨论移位。设 $(A, \text{d})$ 是微分分次代数。设 $M$ 是左微分分次
$A$-模，其底层 $R$-模复形记为 $M^\bullet$。对任意
$k \in \mathbf{Z}$，如下定义 {\it $k$-移位模} $M[k]$：
\begin{enumerate}
\item $M[k]$ 的底层 $R$-模复形是 $M^\bullet[k]$；
\item 作为 $A$-模，其乘法
$$
A^n \times (M[k])^m \longrightarrow (M[k])^{n + m}
$$
等于给定乘法 $A^n \times M^{m + k} \to M^{n + m + k}$ 的
$(-1)^{nk}$ 倍。
\end{enumerate}
来验证该选择满足 Leibniz 规则。设 $a \in A^n$、
$x \in M[k]^m = M^{m + k}$，并以 $\cdot_{M[k]}$ 表示 $M[k]$ 中的乘法，
则
\begin{align*}
\text{d}_{M[k]}(a \cdot_{M[k]} x)
& =
(-1)^{k + nk} \text{d}_M(ax) \\
& =
(-1)^{k + nk} \text{d}(a) x + (-1)^{k + nk + n} a \text{d}_M(x) \\
& =
\text{d}(a) \cdot_{M[k]} x + (-1)^{nk + n} a \text{d}_{M[k]}(x) \\
& =
\text{d}(a) \cdot_{M[k]} x + (-1)^n a \cdot_{M[k]} \text{d}_{M[k]}(x)
\end{align*}
这正是所需结论，因为作为 $A$ 的齐次元素，$a$ 的次数为 $n$。
还注意到，在这些选择下，可把乘法映射看作复形映射
$$
\text{Tot}(A^\bullet \otimes_R M^\bullet[k]) \to
\text{Tot}(A^\bullet \otimes_R M^\bullet)[k] \to
M^\bullet[k]
$$
其中第一个箭头是《更多代数》第 \ref{more-algebra-section-sign-rules} 节
(\ref{more-algebra-item-shift-tensor}) 中的箭头；在本情形下它恰好涉及
上面所选的符号。（事实上，也可由这一观察推出 Leibniz 规则成立。）

\medskip\noindent
由上述规则，有典范认同
$$
(M[k])^{opp} = M^{opp}[k]
$$
，它是右微分分次 $A^{opp}$-模之间不引入符号的认同。换言之，
引理 \ref{dga-lemma-left-right} 的等价与移位函子相容。

\medskip\noindent
上述选择要求作出如下定义。

\begin{definition}
\label{dga-definition-shift-graded-module}
设 $R$ 是环。设 $A$ 是 $\mathbf{Z}$-分次 $R$-代数。
\begin{enumerate}
\item 给定右分次 $A$-模 $M$，把 {\it 第 $k$ 移位 $A$-模} $M[k]$
定义为与之相同的右 $A$-模，但分次改为 $(M[k])^n = M^{n + k}$。
\item 给定左分次 $A$-模 $M$，把 {\it 第 $k$ 移位 $A$-模} $M[k]$
定义为分次满足 $(M[k])^n = M^{n + k}$ 的模，其乘法
$A^n \times (M[k])^m \to (M[k])^{n + m}$
等于给定乘法 $A^n \times M^{m + k} \to M^{n + m + k}$ 的
$(-1)^{nk}$ 倍。
\end{enumerate}
\end{definition}

\noindent
设 $(A, \text{d})$ 是微分分次代数。设 $f, g : M \to N$ 是左微分分次
$A$-模的同态。$f$ 与 $g$ 之间的一个 {\it 同伦}是分次 $A$-模映射
$h : M \to N[-1]$（注意移位！），使得
$$
f(x) - g(x) = \text{d}_N(h(x)) + h(\text{d}_M(x))
$$
对所有 $x \in M$ 成立。若存在这样的同伦，则称 $f$ 与 $g$
{\it 同伦}。因此，$h$ 与 $A$-模结构（$N$ 上采用移位后的结构）及分次
（$N$ 上采用移位后的分次）相容，但不与微分相容。若 $f = g$ 且 $h$
是同伦，则 $h$ 定义左微分分次 $A$-模的态射 $h : M \to N[-1]$。

\medskip\noindent
由上述规则可知，$f, g : M \to N$ 同伦，当且仅当所诱导的态射
$f^{opp}, g^{opp} : M^{opp} \to N^{opp}$ 作为右微分分次
$A^{opp}$-模同态同伦（使用同一个同伦）。

\medskip\noindent
同伦范畴、锥、容许短正合列和有区别三角，全都按与右微分分次模完全相同的
方式定义（在底层 $R$-模复形上，一切都与 $R$-模复形的构造一致）。
这样也得到命题 \ref{dga-proposition-homotopy-category-triangulated}
关于左模的对应版本；或者也可通过研究反代数上的右模推出它。



\section{张量积}
\label{dga-section-tensor-product}

\noindent
设 $R$ 是环。设 $A$ 是 $R$-代数（见第 \ref{dga-section-conventions} 节）。
给定右 $A$-模 $M$ 和左 $A$-模 $N$，有 {\it 张量积}
$$
M \otimes_A N
$$
该张量积是 $R$-模。作为 $R$-模，$M \otimes_A N$ 由符号
$x \otimes y$（$x \in M$、$y \in N$）生成，并满足关系
$$
\begin{matrix}
(x_1 + x_2) \otimes y - x_1 \otimes y - x_2 \otimes y, \\
x \otimes (y_1 + y_2) - x \otimes y_1 - x \otimes y_2, \\
xa \otimes y - x \otimes ay
\end{matrix}
$$
，其中 $a \in A$、$x, x_1, x_2 \in M$ 且 $y, y_1, y_2 \in N$。
下面列出张量积的一些性质。

\medskip\noindent
张量积关于每个变元都是右正合的；事实上，它与直和及任意余极限可交换。

\medskip\noindent
张量积 $M \otimes_A N$ 承载普遍的 $A$-双线性映射
$M \times N \to M \otimes_A N$，$(x, y) \mapsto x \otimes y$。用公式表示为
$$
\text{双线性}_A(M \times N, Q) = \Hom_R(M \otimes_A N, Q)
$$
，对任意 $R$-模 $Q$ 成立。

\medskip\noindent
若 $A$ 是 $\mathbf{Z}$-分次代数，而 $M$、$N$ 是分次 $A$-模，
则 $M \otimes_A N$ 是分次 $R$-模。此时，$M \otimes_A N$ 的第 $n$
分次部分 $(M \otimes_A N)^n$ 等于
$$
\Coker\left(
\bigoplus\nolimits_{r + t + s = n}
M^r \otimes_R A^t \otimes_R N^s \to
\bigoplus\nolimits_{p + q = n} M^p \otimes_R N^q
\right)
$$
其中映射把 $x \otimes a \otimes y$ 送到
$x \otimes ay - xa \otimes y$；这里 $x \in M^r$、$y \in N^s$、
$a \in A^t$ 且 $r + s + t = n$。在这种情形下，映射
$M \times N \to M \otimes_A N$ 是 $A$-双线性的、与分次相容的，
并在下述意义下是普遍的：
$$
\text{分次双线性}_A(M \times N, Q) =
\Hom_{\text{分次 }R\text{-模}}(M \otimes_A N, Q)
$$
，对任意分次 $R$-模 $Q$ 成立；这里采用分次双线性映射的显然概念。

\medskip\noindent
若 $(A, \text{d})$ 是微分分次代数，而 $M$ 和 $N$ 分别是左、右微分分次
$A$-模，则 $M \otimes_A N$ 是微分分次 $R$-模，其微分为
$$
\text{d}(x \otimes y) =
\text{d}(x) \otimes y + (-1)^{\deg(x)}x \otimes \text{d}(y)
$$
，其中 $x \in M$ 和 $y \in N$ 为齐次元素。在这种情形下，映射
$M \times N \to M \otimes_A N$ 是 $A$-双线性的，并与分次及微分相容，
且在下述意义下是普遍的：
$$
\text{微分分次双线性}_A(M \times N, Q) =
\Hom_{\text{Comp}(R)}(M \otimes_A N, Q)
$$
，对任意微分分次 $R$-模 $Q$ 成立；这里采用微分分次双线性映射的显然概念。






\section{Hom 复形与微分分次模}
\label{dga-section-hom-complexes}

\noindent
我们强烈建议读者跳过本节。

\medskip\noindent
设 $R$ 是环，$M^\bullet$ 是 $R$-模复形。考虑 $R$-模复形
$$
E^\bullet = \Hom^\bullet(M^\bullet, M^\bullet)
$$
，它在《更多代数》第 \ref{more-algebra-section-hom-complexes} 节中引入。
由《更多代数》引理 \ref{more-algebra-lemma-composition}，有典范复合法则
$$
\text{Tot}(E^\bullet \otimes_R E^\bullet) \to E^\bullet
$$
，它是复形映射。因此，$E^\bullet$ 连同这一乘法是微分分次 $R$-代数，
记为 $(E, \text{d})$。此外，把 $M^\bullet$ 看作
$\Hom^\bullet(R, M^\bullet)$，可见复合定义乘法
$$
\text{Tot}(E^\bullet \otimes_R M^\bullet) \to M^\bullet
$$
，它使 $M^\bullet$ 成为一个{\bf 左}微分分次 $E$-模，记为 $M$。

\begin{lemma}
\label{dga-lemma-left-module-structure}
在上述情形中，设 $A$ 是微分分次 $R$-代数。在 $M$ 上给定左 $A$-模结构，
等价于给定微分分次 $R$-代数同态 $A \to E$。
\end{lemma}

\begin{proof}
证明略。注意，这一对应不涉及符号。
\end{proof}

\noindent
继续上述讨论，并假设另给一个 $R$-模复形 $N^\bullet$。考虑《更多代数》
第 \ref{more-algebra-section-hom-complexes} 节中引入的 $R$-模复形
$\Hom^\bullet(M^\bullet, N^\bullet)$。如上可见，复合
$$
\text{Tot}(\Hom^\bullet(M^\bullet, N^\bullet) \otimes_R E^\bullet)
\to \Hom^\bullet(M^\bullet, N^\bullet)
$$
定义一个乘法，使 $\Hom^\bullet(M^\bullet, N^\bullet)$ 成为一个
{\bf 右}微分分次 $E$-模。应用引理 \ref{dga-lemma-left-module-structure}，
可得：给定左微分分次 $A$-模 $M$ 和 $R$-模复形 $N^\bullet$，
存在一个典范的右微分分次 $A$-模，其底层 $R$-模复形为
$\Hom^\bullet(M^\bullet, N^\bullet)$，且乘法
$$
\Hom^n(M^\bullet, N^\bullet) \times A^m \longrightarrow
\Hom^{n + m}(M^\bullet, N^\bullet)
$$
把满足 $f_{p, q} \in \Hom(M^{-q}, N^p)$ 的
$f = (f_{p, q})_{p + q = n}$ 及 $a \in A^m$ 送到元素
$f \cdot a = (f_{p, q} \circ a)$，其中 $f_{p, q} \circ a$ 是映射
$$
M^{-q - m} \xrightarrow{a} M^{-q} \xrightarrow{f_{p, q}} N^p, \quad
x \longmapsto f_{p, q}(ax)
$$
；这里不引入符号。用记号 $\Hom(M, N^\bullet)$ 表示这个右微分分次
$A$-模。

\begin{lemma}
\label{dga-lemma-characterize-hom}
设 $R$ 是环。设 $(A, \text{d})$ 是微分分次 $R$-代数。
设 $M'$ 是右微分分次 $A$-模，$M$ 是左微分分次 $A$-模，
$N^\bullet$ 是 $R$-模复形。那么有
$$
\Hom_{\text{Mod}_{(A, d)}}(M', \Hom(M, N^\bullet)) =
\Hom_{\text{Comp}(R)}(M' \otimes_A M, N^\bullet)
$$
，其中如第 \ref{dga-section-tensor-product} 节那样，把
$M \otimes_A M$ 看作 $R$-模复形。
\end{lemma}

\begin{proof}
来说明等式两边都对应于与微分相容的分次 $A$-双线性映射
$$
M' \times M \longrightarrow N^\bullet
$$
。第 \ref{dga-section-tensor-product} 节已经说明右边确实如此。
给定左边的一个元素 $g$，《更多代数》引理
\ref{more-algebra-lemma-compose} 中的等式确定一个 $R$-模复形映射
$g' : \text{Tot}(M' \otimes_R M) \to N^\bullet$。
换言之，得到一个与微分相容的分次 $R$-双线性映射
$g'' : M' \times M \to N^\bullet$。$g$ 的 $A$-线性立即转化为
$g''$ 的 $A$-双线性。
\end{proof}

\noindent
设 $R$、$M^\bullet$、$E^\bullet$、$E$ 和 $M$ 如上。不过，现在假设给定
一个微分分次 $R$-代数 $A$ 以及 $M$ 上的一个{\bf 右}微分分次
$A$-模结构。于是可考虑映射
$$
\text{Tot}(A^\bullet \otimes_R M^\bullet)
\xrightarrow{\psi}
\text{Tot}(A^\bullet \otimes_R M^\bullet)
\to
M^\bullet
$$
其中第一个箭头是 $R$-模复形的微分分次范畴上的交换性约束。
这对应于 $R$-模复形映射
$$
\tau : A^\bullet \longrightarrow E^\bullet
$$
。回忆 $E^n = \prod_{p + q = n} \Hom_R(M^{-q}, M^p)$，并对
$a \in A^n$ 写 $\tau(a) = (\tau_{p, q}(a))_{p + q = n}$。于是
$$
\tau_{p, q}(a) : M^{-q} \longrightarrow M^p,\quad
x \longmapsto (-1)^{\deg(a)\deg(x)}x a = (-1)^{-nq}xa
$$
正如读者根据第 \ref{dga-section-left-modules} 节的讨论所应预料的，
它与 $A$ 上的乘法不相容。事实上，有
$$
\tau(a a') = (-1)^{\deg(a)\deg(a')}\tau(a') \tau(a)
$$
由此得到下述引理。

\begin{lemma}
\label{dga-lemma-right-module-structure}
在上述情形中，设 $A$ 是微分分次 $R$-代数。在 $M$ 上给定右 $A$-模结构，
等价于给定微分分次 $R$-代数同态 $\tau : A \to E^{opp}$。
\end{lemma}

\begin{proof}
见上面的讨论，并注意，由乘法映射
$M^n \times A^m \to M^{n + m}$ 构造 $\tau$ 时使用了符号。
\end{proof}

\noindent
设 $R$、$M^\bullet$、$E^\bullet$、$E$、$A$ 和 $M$ 如上，并如引理中那样
在 $M$ 上给定右微分分次 $A$-模结构。在这种情形下，存在一个典范的
左微分分次 $A$-模，其底层 $R$-模复形为
$\Hom^\bullet(M^\bullet, N^\bullet)$。具体而言，乘法可取为
\begin{align*}
\text{Tot}(A^\bullet \otimes_R \Hom^\bullet(M^\bullet, N^\bullet))
& \xrightarrow{\psi}
\text{Tot}(\Hom^\bullet(M^\bullet, N^\bullet) \otimes_R A^\bullet) \\
& \xrightarrow{\tau}
\text{Tot}(\Hom^\bullet(M^\bullet, N^\bullet) \otimes_R
\Hom^\bullet(M^\bullet, M^\bullet)) \\
& \to
\text{Tot}(\Hom^\bullet(M^\bullet, N^\bullet)
\end{align*}
第一个箭头使用 $R$-模复形范畴上的交换性约束，第二个箭头如上所述，
第三个箭头是 Hom 复形的复合法则。每个映射都是复形映射，因而所得映射
也是复形映射。事实上，这一构造使
$\Hom^\bullet(M^\bullet, N^\bullet)$ 成为左微分分次 $A$-模
（乘法的结合性可用对称幺半结构证明，也可使用下面的公式直接计算）。
把该乘法明确写为
$$
A^n \times \Hom^m(M^\bullet, N^\bullet) \longrightarrow
\Hom^{n + m}(M^\bullet, N^\bullet)
$$
。它把 $a \in A^n$ 以及满足 $f_{p, q} \in \Hom(M^{-q}, N^p)$ 的
$f = (f_{p, q})_{p + q = m}$ 送到元素 $a \cdot f$；后者的分量为
$$
(-1)^{nm}f_{p, q} \circ \tau_{-q, q + n}(a) =
(-1)^{nm - n(q + n)}f_{p, q} \circ a =
(-1)^{np + n} f_{p, q} \circ a
$$
，它位于 $\Hom_R(M^{-q - n}, N^p)$ 中；其中 $f_{p, q} \circ a$ 是映射
$$
M^{-q - n} \xrightarrow{a} M^{-q} \xrightarrow{f_{p, q}} N^p,\quad
x \longmapsto f_{p, q}(xa)
$$
。这里确实引入了符号 $(-1)^{np + n}$。用记号
$\Hom(M, N^\bullet)$ 表示这个左微分分次 $A$-模。

\begin{lemma}
\label{dga-lemma-characterize-hom-other-side}
设 $R$ 是环。设 $(A, \text{d})$ 是微分分次 $R$-代数。
设 $M$ 是右微分分次 $A$-模，$M'$ 是左微分分次 $A$-模，
$N^\bullet$ 是 $R$-模复形。那么有
$$
\Hom_{\text{左微分分次 }A\text{-模}}(M', \Hom(M, N^\bullet)) =
\Hom_{\text{Comp}(R)}(M \otimes_A M', N^\bullet)
$$
，其中如第 \ref{dga-section-tensor-product} 节那样，把
$M \otimes_A M'$ 看作 $R$-模复形。
\end{lemma}

\begin{proof}
来说明等式两边都对应于与微分相容的分次 $A$-双线性映射
$$
M \times M' \longrightarrow N^\bullet
$$
。第 \ref{dga-section-tensor-product} 节已经说明右边确实如此。
给定左边的一个元素 $g$，《更多代数》引理
\ref{more-algebra-lemma-compose} 中的等式确定复形映射
$g' : \text{Tot}(M' \otimes_R M) \to N^\bullet$。
预复合以交换性约束，得到
$$
\text{Tot}(M \otimes_R M') \xrightarrow{\psi}
\text{Tot}(M' \otimes_R M) \xrightarrow{g'}
N^\bullet
$$
，它对应于与微分相容的分次 $R$-双线性映射
$g'' : M \times M' \to N^\bullet$。$g$ 的 $A$-线性立即转化为
$g''$ 的 $A$-双线性。具体地，设 $x \in M^e$、
$x' \in (M')^{e'}$、$a \in A^n$。一方面有
\begin{align*}
g''(x, ax')
& =
(-1)^{e(n + e')} g'(ax' \otimes x) \\
& =
(-1)^{e(n + e')} g(ax')(x) \\
& =
(-1)^{e(n + e')} (a \cdot g(x'))(x) \\
& =
(-1)^{e(n + e') + n(n + e + e') + n} g(x')(xa)
\end{align*}
另一方面有
$$
g''(xa, x') = (-1)^{(e + n)e'} g'(x' \otimes xa) =
(-1)^{(e + n)e'} g(x')(xa) 
$$
对指数作一个简单的模 $2$ 计算，便知两者相同。
\end{proof}

\begin{remark}
\label{dga-remark-evaluation-map-left}
设 $R$ 是环，$A$ 是微分分次 $R$-代数，$M$ 是左微分分次 $A$-模，
$N^\bullet$ 是 $R$-模复形。上述构造先产生右微分分次 $A$-模
$\Hom(M, N^\bullet)$，再产生左微分分次 $A$-模
$\Hom(\Hom(M, N^\bullet), N^\bullet)$。我们断言存在求值映射
$$
ev : M \longrightarrow \Hom(\Hom(M, N^\bullet), N^\bullet)
$$
，它位于左微分分次 $A$-模范畴中。由引理
\ref{dga-lemma-characterize-hom}，为定义它，只需构造与分次及微分相容的
$A$-双线性配对
$$
\Hom(M, N^\bullet) \times M \longrightarrow N^\bullet
$$
。为此取
$$
(f, x) \longmapsto f(x)
$$
它与分次、微分及 $A$-双线性的相容性留给读者验证。$ev$ 在底层
$R$-模复形上的映射是《更多代数》条目
(\ref{more-algebra-item-evaluation})。
\end{remark}

\begin{remark}
\label{dga-remark-evaluation-map-right}
设 $R$ 是环，$A$ 是微分分次 $R$-代数，$M$ 是右微分分次 $A$-模，
$N^\bullet$ 是 $R$-模复形。上述构造先产生左微分分次 $A$-模
$\Hom(M, N^\bullet)$，再产生右微分分次 $A$-模
$\Hom(\Hom(M, N^\bullet), N^\bullet)$。我们断言存在求值映射
$$
ev : M \longrightarrow \Hom(\Hom(M, N^\bullet), N^\bullet)
$$
，它位于右微分分次 $A$-模范畴中。由引理
\ref{dga-lemma-characterize-hom}，为定义它，只需构造与分次及微分相容的
$A$-双线性配对
$$
M \times \Hom(M, N^\bullet) \longrightarrow N^\bullet
$$
。为此取
$$
(x, f) \longmapsto (-1)^{\deg(x)\deg(f)}f(x)
$$
它与分次、微分及 $A$-双线性的相容性留给读者验证。$ev$ 在底层
$R$-模复形上的映射是《更多代数》条目
(\ref{more-algebra-item-evaluation})。
\end{remark}

\begin{remark}
\label{dga-remark-shift-dual}
设 $R$ 是环，$A$ 是微分分次 $R$-代数，$M^\bullet$ 和 $N^\bullet$
是 $R$-模复形。设 $k \in \mathbf{Z}$，考虑 $R$-模复形同构
$$
\Hom^\bullet(M^\bullet, N^\bullet)[-k]
\longrightarrow
\Hom^\bullet(M^\bullet[k], N^\bullet)
$$
，它在《更多代数》条目 (\ref{more-algebra-item-shift-hom}) 中定义。
若 $M^\bullet$ 具有左（分别地，右）微分分次 $A$-模结构，
则这是右（分别地，左）微分分次 $A$-模的映射（模结构采用本节的定义）。
验证略。提醒读者：左分次 $A$-模的移位上的 $A$-模结构用一个符号定义；
见定义 \ref{dga-definition-shift-graded-module}。
\end{remark}






\section{代数上的投射模}
\label{dga-section-projectives-over-algebras}

\noindent
本节讨论代数上的投射模，类似于《代数》第
\ref{algebra-section-projective} 节。本节或许应移到别处。

\medskip\noindent
设 $R$ 是环，$A$ 是 $R$-代数；约定见第 \ref{dga-section-conventions} 节。
显然 $A$ 是投射右 $A$-模，因为对任意右 $A$-模 $M$ 都有
$\Hom_A(A, M) = M$（因而 $\Hom_A(A, -)$ 正合）。反之，设 $P$ 是投射
右 $A$-模。通过选择 $P$ 在 $A$ 上的一组生成元
$\{p_i\}_{i \in I}$，可以选取满射 $\bigoplus_{i \in I} A \to P$。
由于 $P$ 投射，该满射有左逆，故 $P$ 同构于某个自由模的直和项，
与交换情形完全相同（《代数》引理
\ref{algebra-lemma-characterize-projective}）。

\medskip\noindent
我们得到：
\begin{enumerate}
\item $A$-模范畴有足够多投射对象；
\item $A$ 是投射 $A$-模；
\item 每个 $A$-模都是若干个 $A$ 的直和的商；
\item 每个投射 $A$-模都是若干个 $A$ 的直和的一个直和项。
\end{enumerate}



\section{分次代数上的投射模}
\label{dga-section-projectives-over-graded-algebras}

\noindent
本节讨论分次代数上的投射分次模，类似于《代数》第
\ref{algebra-section-projective} 节。

\medskip\noindent
设 $R$ 是环。设 $A$ 是 $R$ 上的 $\mathbf{Z}$-分次代数。
约定见第 \ref{dga-section-conventions} 节。用 $\text{Mod}_A$ 表示分次右
$A$-模范畴。对整数 $k$，用 $A[k]$ 表示 $A$ 的移位。
对于分次右 $A$-模，有
$$
\Hom_{\text{Mod}_A}(A[k], M) = M^{-k}
$$
由于函子 $M \mapsto M^{-k}$ 在 $\text{Mod}_A$ 上正合，
$A[k]$ 是 $\text{Mod}_A$ 的投射对象。反之，假设 $P$ 是
$\text{Mod}_A$ 的投射对象。选择 $P$ 作为 $A$-模的一组齐次生成元，
可以找到满射
$$
\bigoplus\nolimits_{i \in I} A[k_i] \longrightarrow P
$$
因此，$\text{Mod}_A$ 的投射对象是若干个移位 $A[k]$ 的直和的一个直和项。

\medskip\noindent
我们得到：
\begin{enumerate}
\item 分次 $A$-模范畴有足够多投射对象；
\item 对每个 $k \in \mathbf{Z}$，$A[k]$ 都是投射 $A$-模；
\item 每个分次 $A$-模都是若干个模 $A[k]$（$k$ 可变）的直和的商；
\item 每个投射 $A$-模都是若干个模 $A[k]$（$k$ 可变）的直和的一个直和项。
\end{enumerate}




\section{投射模与微分分次代数}
\label{dga-section-projective-over-differential-graded}

\noindent
若 $(A, \text{d})$ 是微分分次代数，而 $P$ 是
$\text{Mod}_{(A, \text{d})}$ 的对象，则称 {\it $P$ 作为分次 $A$-模是投射的}，
有时也称 {\it $P$ 是分次投射的}；这意指 $P$ 是第
\ref{dga-section-projectives-over-graded-algebras} 节中分次 $A$-模的阿贝尔范畴
$\text{Mod}_A$ 的投射对象。

\begin{lemma}
\label{dga-lemma-target-graded-projective}
设 $(A, \text{d})$ 是微分分次代数。设 $M \to P$ 是微分分次 $A$-模的
满同态。若 $P$ 作为分次 $A$-模是投射的，则 $M \to P$ 是容许满态射。
\end{lemma}

\begin{proof}
这立即由定义得出。
\end{proof}

\begin{lemma}
\label{dga-lemma-hom-from-shift-free}
设 $(A, d)$ 是微分分次代数。那么，对任意微分分次 $A$-模 $M$，有
$$
\Hom_{\text{Mod}_{(A, \text{d})}}(A[k], M) =
\Ker(\text{d} : M^{-k} \to M^{-k + 1})
$$
并且
$$
\Hom_{K(\text{Mod}_{(A, \text{d})})}(A[k], M) = H^{-k}(M)
$$
。
\end{lemma}

\begin{proof}
立即由定义得出。
\end{proof}







\section{代数上的内射模}
\label{dga-section-modules-noncommutative}

\noindent
本节讨论代数上的内射模，类似于《更多代数》第
\ref{more-algebra-section-injectives-modules} 节。本节或许应移到别处。

\medskip\noindent
设 $R$ 是环，$A$ 是 $R$-代数；约定见第 \ref{dga-section-conventions} 节。
对右 $A$-模 $M$，置
$$
M^\vee = \Hom_\mathbf{Z}(M, \mathbf{Q}/\mathbf{Z})
$$
并借助乘法 $(a f)(x) = f(xa)$ 把它看作左 $A$-模。事实上，
$((ab)f)(x) = f(xab) = (bf)(xa) = (a(bf))(x)$。反之，若 $M$ 是左
$A$-模，则 $M^\vee$ 是右 $A$-模。由于 $\mathbf{Q}/\mathbf{Z}$
是内射阿贝尔群（《更多代数》引理
\ref{more-algebra-lemma-injective-abelian}），函子 $M \mapsto M^\vee$
是正合的（《更多代数》引理 \ref{more-algebra-lemma-vee-exact}）。
此外，对所有模 $M$，求值映射 $M \to (M^\vee)^\vee$ 是单射
（《更多代数》引理 \ref{more-algebra-lemma-ev-injective}）。

\medskip\noindent
我们断言 $A^\vee$ 是内射右 $A$-模。事实上，给定右 $A$-模 $N$，有
$$
\Hom_A(N, A^\vee) =
\Hom_A(N, \Hom_\mathbf{Z}(A, \mathbf{Q}/\mathbf{Z})) = N^\vee
$$
；结论由函子 $N \mapsto N^\vee$ 的正合性得出。第二个等式成立，是因为
$$
\Hom_\mathbf{Z}(N, \Hom_\mathbf{Z}(A, \mathbf{Q}/\mathbf{Z})) =
\Hom_\mathbf{Z}(N \otimes_\mathbf{Z} A, \mathbf{Q}/\mathbf{Z})
$$
，见《代数》引理 \ref{algebra-lemma-hom-from-tensor-product}。
在这个模中，$A$-线性恰好挑出双线性映射
$\varphi : N \times A \to \mathbf{Q}/\mathbf{Z}$，它们在
$x \otimes a$ 和 $xa \otimes 1$ 上取相同值；换言之，它们来自
$N^\vee$ 的元素。

\medskip\noindent
最后，对每个右 $A$-模 $M$，可选择满射
$\bigoplus_{i \in I} A \to M^\vee$，从而得到单射
$M \to (M^\vee)^\vee \to \prod_{i \in I} A^\vee$。

\medskip\noindent
我们得到：
\begin{enumerate}
\item $A$-模范畴有足够多内射对象；
\item $A^\vee$ 是内射 $A$-模；以及
\item 每个 $A$-模都单射到若干个 $A^\vee$ 的乘积中。
\end{enumerate}





\section{分次代数上的内射模}
\label{dga-section-modules-noncommutative-graded}

\noindent
本节讨论分次代数上的内射分次模，类似于《更多代数》第
\ref{more-algebra-section-injectives-modules} 节。

\medskip\noindent
设 $R$ 是环。设 $A$ 是 $R$ 上的 $\mathbf{Z}$-分次代数。
约定见第 \ref{dga-section-conventions} 节。若 $M$ 是分次 $R$-模，置
$$
M^\vee =
\bigoplus\nolimits_{n \in \mathbf{Z}}
\Hom_\mathbf{Z}(M^{-n}, \mathbf{Q}/\mathbf{Z}) =
\bigoplus\nolimits_{n \in \mathbf{Z}} (M^{-n})^\vee
$$
，把它视为分次 $R$-模（$R$ 在齐次部分上的作用不带符号）。
若 $M$ 具有左分次 $A$-模结构，则如下定义 $M^\vee$ 上的右分次
$A$-模结构：令 $a \in A^m$ 按
$$
(M^{-n})^\vee \to (M^{-n - m})^\vee, \quad
f \mapsto f \circ a
$$
作用，如第 \ref{dga-section-hom-complexes} 节所述。若 $M$ 具有右分次
$A$-模结构，则如下定义 $M^\vee$ 上的左分次 $A$-模结构：
令 $a \in A^n$ 按
$$
(M^{-m})^\vee \to (M^{-m - n})^\vee, \quad
f \mapsto (-1)^{nm}f \circ a
$$
作用，如第 \ref{dga-section-hom-complexes} 节所述（这个符号是被迫的，
因为研究微分分次模时我们希望使用同一个公式；若只关心分次模，
则可在此省略该符号）。在（左或右）分次 $A$-模范畴上，函子
$M \mapsto M^\vee$ 正合（在各分次部分上验证）。此外，有单射的求值映射
$$
ev : M \longrightarrow (M^\vee)^\vee, \quad
ev^n = (-1)^n \text{ 倍的求值映射 }M^n \to ((M^n)^\vee)^\vee
$$
，它是分次 $R$-模的映射；见《更多代数》条目
(\ref{more-algebra-item-evaluation})。若 $M$ 是左（分别地，右）$A$-模，
则这个求值映射是左（分别地，右）$A$-模同态；见评注
\ref{dga-remark-evaluation-map-left} 和 \ref{dga-remark-evaluation-map-right}。
最后，给定 $k \in \mathbf{Z}$，有典范同构
$$
M^\vee[-k] \longrightarrow (M[k])^\vee
$$
；它是使用一个符号的分次 $R$-模同构。若 $M$ 是左（分别地，右）
$A$-模，则它是右（分别地，左）$A$-模的同构。见评注
\ref{dga-remark-shift-dual}。

\medskip\noindent
我们断言 $A^\vee$ 是分次右 $A$-模范畴 $\text{Mod}_A$ 的内射对象。
事实上，给定分次右 $A$-模 $N$，有
$$
\Hom_{\text{Mod}_A}(N, A^\vee) =
\Hom_{\text{Comp}(\mathbf{Z})}(N \otimes_A A, \mathbf{Q}/\mathbf{Z})) =
(N^0)^\vee
$$
；这是引理 \ref{dga-lemma-characterize-hom} 应用于所有微分均为零的情形。
结论由函子 $N \mapsto (N^0)^\vee = (N^\vee)^0$ 的正合性得出。

\medskip\noindent
最后，对每个分次右 $A$-模 $M$，可以选择分次左 $A$-模的满射
$$
\bigoplus\nolimits_{i \in I} A[k_i] \to M^\vee
$$
，其中 $A[k_i]$ 表示 $A$ 由 $k_i \in \mathbf{Z}$ 给出的移位。
为此，选择 $M^\vee$ 的齐次生成元。这样得到单射
$$
M \to (M^\vee)^\vee \to \prod A[k_i]^\vee = \prod A^\vee[-k_i]
$$
注意，上式中的乘积是分次模范畴中的乘积（换言之，先在每个次数上取乘积，
再对各部分取直和）。

\medskip\noindent
我们得到：
\begin{enumerate}
\item 分次 $A$-模范畴有足够多内射对象；
\item 对每个 $k \in \mathbf{Z}$，模 $A^\vee[k]$ 都是内射的；以及
\item 每个 $A$-模都单射到分次模范畴中若干个移位 $A^\vee[k]$ 的乘积中。
\end{enumerate}





\section{内射模与微分分次代数}
\label{dga-section-modules-noncommutative-differential-graded}

\noindent
若 $(A, \text{d})$ 是微分分次代数，而 $I$ 是
$\text{Mod}_{(A, \text{d})}$ 的对象，则称 {\it $I$ 作为分次 $A$-模是内射的}，
有时也称 {\it $I$ 是分次内射的}；这意指 $I$ 是分次 $A$-模的阿贝尔范畴
$\text{Mod}_A$ 的内射对象。

\begin{lemma}
\label{dga-lemma-source-graded-injective}
设 $(A, \text{d})$ 是微分分次代数。设 $I \to M$ 是微分分次 $A$-模的
单同态。若 $I$ 是分次内射的，则 $I \to M$ 是容许单态射。
\end{lemma}

\begin{proof}
这立即由定义得出。
\end{proof}

\noindent
设 $(A, \text{d})$ 是微分分次代数。若 $M$ 是左（分别地，右）微分分次
$A$-模，则
$$
M^\vee = \Hom^\bullet(M^\bullet, \mathbf{Q}/\mathbf{Z})
$$
连同第 \ref{dga-section-modules-noncommutative-graded} 节中构造的 $A$-模结构，
由第 \ref{dga-section-hom-complexes} 节的讨论，是右（分别地，左）微分分次
$A$-模。由评注 \ref{dga-remark-evaluation-map-left} 和
\ref{dga-remark-evaluation-map-right}，第
\ref{dga-section-modules-noncommutative-graded} 节的求值映射
$$
M \longrightarrow (M^\vee)^\vee
$$
是左（分别地，右）微分分次 $A$-模的同态。

\begin{lemma}
\label{dga-lemma-map-into-dual}
设 $(A, \text{d})$ 是微分分次代数。若 $M$ 是左微分分次 $A$-模，
$N$ 是右微分分次 $A$-模，则
\begin{align*}
\Hom_{\text{Mod}_{(A, \text{d})}}(N, M^\vee)
& =
\Hom_{\text{Comp}(\mathbf{Z})}(N \otimes_A M, \mathbf{Q}/\mathbf{Z}) \\
& =
\text{微分分次双线性}_A(N \times M, \mathbf{Q}/\mathbf{Z})
\end{align*}
\end{lemma}

\begin{proof}
第一个等式是引理 \ref{dga-lemma-characterize-hom}，第二个等式已在第
\ref{dga-section-tensor-product} 节中证明。
\end{proof}

\begin{lemma}
\label{dga-lemma-hom-into-shift-dual-free}
设 $(A, \text{d})$ 是微分分次代数。那么有
$$
\Hom_{\text{Mod}_{(A, \text{d})}}(M, A^\vee[k]) =
\Ker(\text{d} : (M^\vee)^k \to (M^\vee)^{k + 1})
$$
并且
$$
\Hom_{K(\text{Mod}_{(A, \text{d})})}(M, A^\vee[k]) = H^k(M^\vee)
$$
；它们是关于微分分次 $A$-模 $M$ 的函子等式。
\end{lemma}

\begin{proof}
这由上面的讨论显然可见。
\end{proof}















\section{P-分辨}
\label{dga-section-P-resolutions}

\noindent
本节是《导出范畴》第 \ref{derived-section-unbounded} 节的对应版本。

\medskip\noindent
设 $(A, \text{d})$ 是微分分次代数。设 $P$ 是微分分次 $A$-模。
若存在滤过
$$
0 = F_{-1}P \subset F_0P \subset F_1P \subset \ldots \subset P
$$
，它由微分分次子模组成，并且
\begin{enumerate}
\item $P = \bigcup F_pP$；
\item 包含映射 $F_iP \to F_{i + 1}P$ 是容许单态射；
\item 作为微分分次 $A$-模，商 $F_{i + 1}P/F_iP$ 同构于若干个
$A[k]$ 的直和，
\end{enumerate}
则称 $P$ {\it 具有性质 (P)}。事实上，条件 (2) 是条件 (3) 的推论；
见引理 \ref{dga-lemma-target-graded-projective}。此外，读者可以验证，
作为分次 $A$-模，$P$ 同构于若干个 $A$ 的移位的直和。

\begin{lemma}
\label{dga-lemma-property-P-sequence}
设 $(A, \text{d})$ 是微分分次代数。设 $P$ 是微分分次 $A$-模。
若 $F_\bullet$ 是性质 (P) 中的滤过，则得到容许短正合列
$$
0 \to
\bigoplus\nolimits F_iP \to
\bigoplus\nolimits F_iP \to P \to 0
$$
，它由微分分次 $A$-模组成。
\end{lemma}

\begin{proof}
第二个映射是各包含映射的直和。第一个映射在源的直和项 $F_iP$ 上，
等于恒等映射 $F_iP \to F_iP$ 与包含映射
$F_iP \to F_{i + 1}P$ 的负映射之和。选择分次 $A$-模同态
$s_i : F_{i + 1}P \to F_iP$，使它们是包含映射的左逆。
复合后，对所有 $j > i$ 得到映射 $s_{j, i} : F_jP \to F_iP$。
于是，第一个箭头的一个左逆把 $x \in F_jP$ 映到
$\bigoplus F_iP$ 中的
$(s_{j, 0}(x), s_{j, 1}(x), \ldots, s_{j, j - 1}(x), 0, \ldots)$。
\end{proof}

\noindent
下述引理表明，具有性质 (P) 的微分分次模是 K-内射模的对偶概念
（即在某种意义下它们是 K-投射的）。见《导出范畴》定义
\ref{derived-definition-K-injective}。

\begin{lemma}
\label{dga-lemma-property-P-K-projective}
设 $(A, \text{d})$ 是微分分次代数。设 $P$ 是具有性质 (P) 的微分分次
$A$-模。那么
$$
\Hom_{K(\text{Mod}_{(A, \text{d})})}(P, N) = 0
$$
对所有无环微分分次 $A$-模 $N$ 成立。
\end{lemma}

\begin{proof}
将使用 $K(\text{Mod}_{(A, \text{d})})$ 是三角范畴这一事实
（命题 \ref{dga-proposition-homotopy-category-triangulated}）。
设 $F_\bullet$ 是 $P$ 上性质 (P) 中的滤过。引理
\ref{dga-lemma-property-P-sequence} 的短正合列产生一个有区别三角。
因此，由《导出范畴》引理
\ref{derived-lemma-representable-homological}，只需证明
$$
\Hom_{K(\text{Mod}_{(A, \text{d})})}(F_iP, N) = 0
$$
对所有无环微分分次 $A$-模 $N$ 和所有 $i$ 成立。每个微分分次模
$F_iP$ 都具有由容许单态射给出的有限滤过，其分次部分是若干个移位
$A[k]$ 的直和。因此，只需证明
$$
\Hom_{K(\text{Mod}_{(A, \text{d})})}(A[k], N) = 0
$$
对所有无环微分分次 $A$-模 $N$ 和所有 $k$ 成立。这由引理
\ref{dga-lemma-hom-from-shift-free} 得出。
\end{proof}

\begin{lemma}
\label{dga-lemma-good-quotient}
设 $(A, \text{d})$ 是微分分次代数。设 $M$ 是微分分次 $A$-模。
存在具有下列性质的微分分次 $A$-模同态 $P \to M$：
\begin{enumerate}
\item $P \to M$ 是满射；
\item $\Ker(\text{d}_P) \to \Ker(\text{d}_M)$ 是满射；以及
\item $P$ 位于容许短正合列 $0 \to P' \to P \to P'' \to 0$ 中，
其中 $P'$、$P''$ 是若干个 $A$ 的移位的直和。
\end{enumerate}
\end{lemma}

\begin{proof}
令 $P_k$ 为自由 $A$-模，其生成元 $x,y$ 的次数分别为 $k$ 和 $k+1$。
通过置 $\text{d}(x) = y$、$\text{d}(y) = 0$，在 $P_k$ 上定义微分分次
$A$-模结构。对每个元素 $m \in M^k$，都有同态 $P_k \to M$，
它把 $x$ 映到 $m$，把 $y$ 映到 $\text{d}(m)$。因此，若干个 $P_k$
的直和到 $M$ 有一个满射。这显然产生具有性质 (1) 和 (3) 的
$P \to M$。为得到性质 (2)，注意：若
$m \in \Ker(\text{d}_M)$ 的次数为 $k$，则有把 $1$ 映到 $m$ 的映射
$A[k] \to M$。因此，加入若干个 $A$ 的移位的直和即可满足 (2)。
\end{proof}

\begin{lemma}
\label{dga-lemma-resolve}
设 $(A, \text{d})$ 是微分分次代数。设 $M$ 是微分分次 $A$-模。
存在微分分次 $A$-模同态 $P \to M$，使得
\begin{enumerate}
\item $P \to M$ 是拟同构；并且
\item $P$ 具有性质 (P)。
\end{enumerate}
\end{lemma}

\begin{proof}
置 $M=M_0$。归纳选取短正合列
$$
0 \to M_{i + 1} \to P_i \to M_i \to 0
$$
，其中映射 $P_i \to M_i$ 按引理 \ref{dga-lemma-good-quotient} 选取。
这给出一个“分辨”
$$
\ldots \to P_2 \xrightarrow{f_2} P_1 \xrightarrow{f_1} P_0 \to M \to 0
$$
然后置
$$
P = \bigoplus\nolimits_{i \geq 0} P_i
$$
，把它视为 $A$-模，其分次由
$P^n = \bigoplus_{a + b = n} P_{-a}^b$ 给出，微分则如双复形关联全复形的
构造那样，由
$$
\text{d}_P(x) = f_{-a}(x) + (-1)^a \text{d}_{P_{-a}}(x)
$$
给出，其中 $x \in P_{-a}^b$。在这些约定下，$P$ 确实是微分分次
$A$-模。回忆每个 $P_i$ 都有两步滤过
$0 \to P_i' \to P_i \to P_i'' \to 0$，置
$$
F_{2i}P = \bigoplus\nolimits_{i \geq j \geq 0} P_j
\subset
\bigoplus\nolimits_{i \geq 0} P_i = P
$$
并把 $P'_{i + 1}$ 加到 $F_{2i}P$ 中得到 $F_{2i + 1}$。
这些都是微分分次子模，且相继的商是若干个 $A$ 的移位的直和。
由引理 \ref{dga-lemma-target-graded-projective}，包含映射
$F_iP \to F_{i + 1}P$ 是容许单态射。最后，还须证明由增广
$P_0 \to M$ 给出的映射 $P \to M$ 是拟同构。这由《同调代数》引理
\ref{homology-lemma-good-resolution-gives-qis} 得出。
\end{proof}





\section{I-分辨}
\label{dga-section-I-resolutions}

\noindent
本节是 P-分辨一节的对偶。

\medskip\noindent
设 $(A, \text{d})$ 是微分分次代数。设 $I$ 是微分分次 $A$-模。
若存在滤过
$$
I = F_0I \supset F_1I \supset F_2I \supset \ldots \supset 0
$$
，它由微分分次子模组成，并且
\begin{enumerate}
\item $I = \lim I/F_pI$；
\item 映射 $I/F_{i + 1}I \to I/F_iI$ 是容许满态射；
\item 作为微分分次 $A$-模，商 $F_iI/F_{i + 1}I$ 同构于第
\ref{dga-section-modules-noncommutative-differential-graded} 节构造的若干个
$A^\vee[k]$ 的直积，
\end{enumerate}
则称 $I$ {\it 具有性质 (I)}。事实上，条件 (2) 是条件 (3) 的推论；见引理
\ref{dga-lemma-source-graded-injective}。读者可以验证，作为分次模，$I$ 同构于
若干个 $A^\vee[k]$ 的直积。

\begin{lemma}
\label{dga-lemma-property-I-sequence}
设 $(A, \text{d})$ 是微分分次代数。设 $I$ 是微分分次 $A$-模。
若 $F_\bullet$ 是性质 (I) 中的滤过，则得到容许短正合列
$$
0 \to I \to
\prod\nolimits I/F_iI \to
\prod\nolimits I/F_iI \to 0
$$
，它由微分分次 $A$-模组成。
\end{lemma}

\begin{proof}
略。提示：这是引理 \ref{dga-lemma-property-P-sequence} 的对偶。
\end{proof}

\noindent
下述引理表明，具有性质 (I) 的微分分次模是 K-内射模的对应概念。见
《导出范畴》定义 \ref{derived-definition-K-injective}。

\begin{lemma}
\label{dga-lemma-property-I-K-injective}
设 $(A, \text{d})$ 是微分分次代数。设 $I$ 是具有性质 (I) 的微分分次
$A$-模。那么
$$
\Hom_{K(\text{Mod}_{(A, \text{d})})}(N, I) = 0
$$
对所有无环微分分次 $A$-模 $N$ 成立。
\end{lemma}

\begin{proof}
将使用 $K(\text{Mod}_{(A, \text{d})})$ 是三角范畴这一事实
（命题 \ref{dga-proposition-homotopy-category-triangulated}）。设 $F_\bullet$
是 $I$ 上性质 (I) 中的滤过。引理 \ref{dga-lemma-property-I-sequence}
的短正合列产生一个有区别三角。因此，由《导出范畴》引理
\ref{derived-lemma-representable-homological}，只需证明
$$
\Hom_{K(\text{Mod}_{(A, \text{d})})}(N, I/F_iI) = 0
$$
对所有无环微分分次 $A$-模 $N$ 和所有 $i$ 成立。每个微分分次模
$I/F_iI$ 都具有由容许单态射给出的有限滤过，其分次部分是若干个
$A^\vee[k]$ 的直积。因此，只需证明
$$
\Hom_{K(\text{Mod}_{(A, \text{d})})}(N, A^\vee[k]) = 0
$$
对所有无环微分分次 $A$-模 $N$ 和所有 $k$ 成立。这由引理
\ref{dga-lemma-hom-into-shift-dual-free} 以及 $(-)^\vee$ 是正合函子这一事实得出。
\end{proof}

\begin{lemma}
\label{dga-lemma-good-sub}
设 $(A, \text{d})$ 是微分分次代数。设 $M$ 是微分分次 $A$-模。
存在具有下列性质的微分分次 $A$-模同态 $M \to I$：
\begin{enumerate}
\item $M \to I$ 是单射；
\item $\Coker(\text{d}_M) \to \Coker(\text{d}_I)$ 是单射；以及
\item $I$ 位于容许短正合列 $0 \to I' \to I \to I'' \to 0$ 中，
其中 $I'$、$I''$ 是若干个 $A^\vee$ 的移位的直积。
\end{enumerate}
\end{lemma}

\begin{proof}
将使用第 \ref{dga-section-modules-noncommutative-differential-graded} 节构造的
函子 $N \mapsto N^\vee$（从左微分分次模到右微分分次模，以及从右微分
分次模到左微分分次模）及其全部性质。对每个 $k \in \mathbf{Z}$，令
$Q_k$ 为自由左 $A$-模，其生成元 $x,y$ 的次数分别为 $k$ 和 $k+1$。
通过置 $\text{d}(x) = y$、$\text{d}(y) = 0$，在 $Q_k$ 上定义左微分分次
$A$-模结构。完全仿照引理 \ref{dga-lemma-good-quotient} 的证明，可得左微分分次
$A$-模的满射
$$
\bigoplus\nolimits_{i \in I} Q_{k_i} \longrightarrow M^\vee
$$
。于是可以考虑单射
$$
M \to (M^\vee)^\vee \to (\bigoplus\nolimits_{i \in I} Q_{k_i})^\vee =
\prod\nolimits_{i \in I} I_{k_i}
$$
，其中 $I_k = Q_{-k}^\vee$ 是“对偶”右微分分次 $A$-模。此外，短正合列
$0 \to A[-k - 1] \to Q_k \to A[-k] \to 0$ 产生短正合列
$0 \to A^\vee[k] \to I_k \to A^\vee[k + 1] \to 0$.

\medskip\noindent
上一段的结果给出具有性质 (1) 和 (3) 的 $M \to I$。为得到性质 (2)，
设 $\overline{m} \in \Coker(\text{d}_M)$ 是次数为 $k$ 的非零元素。
选取映射 $\lambda : M^k \to \mathbf{Q}/\mathbf{Z}$，使它在
$\Im(M^{k - 1} \to M^k)$ 上为零，而在 $m$ 上不为零。由引理
\ref{dga-lemma-hom-into-shift-dual-free}，这对应于在 $m$ 上不为零的微分分次
$A$-模同态 $M \to A^\vee[k]$。因此，加入若干个 $A^\vee$ 的移位的副本
的直积，即可满足 (2)。
\end{proof}

\begin{lemma}
\label{dga-lemma-right-resolution}
设 $(A, \text{d})$ 是微分分次代数。设 $M$ 是微分分次 $A$-模。
存在微分分次 $A$-模同态 $M \to I$，使得
\begin{enumerate}
\item $M \to I$ 是拟同构；并且
\item $I$ 具有性质 (I)。
\end{enumerate}
\end{lemma}

\begin{proof}
置 $M=M_0$。归纳选取短正合列
$$
0 \to M_i \to I_i \to M_{i + 1} \to 0
$$
，其中映射 $M_i \to I_i$ 按引理 \ref{dga-lemma-good-sub} 选取。这给出一个“分辨”
$$
0 \to M \to I_0 \xrightarrow{f_0} I_1 \xrightarrow{f_1} I_1 \to \ldots
$$
以 $I$ 表示微分分次 $A$-模，其分次部分为
$$
I^n = \prod\nolimits_{i \geq 0} I^{n - i}_i
$$
，微分定义为
$$
\text{d}_I(x) = f_i(x) + (-1)^i \text{d}_{I_i}(x)
$$
，其中 $x \in I_i^{n-i}$。在这些约定下，$I$ 确实是微分分次 $A$-模。
回忆每个 $I_i$ 都有两步滤过 $0 \to I_i' \to I_i \to I_i'' \to 0$，置
$$
F_{2i}I^n = \prod\nolimits_{j \geq i} I^{n - j}_j
\subset
\prod\nolimits_{i \geq 0} I^{n - i}_i = I^n
$$
，并把因子 $I'_{i+1}$ 加到 $F_{2i}I$ 中得到 $F_{2i+1}I$。这些都是微分分次
子模，且相继的商是若干个 $A^\vee$ 的移位的直积。由引理
\ref{dga-lemma-source-graded-injective}，包含映射 $F_{i+1}I \to F_iI$ 是容许
单态射。最后，还须证明由增广 $M \to I_0$ 给出的映射 $M \to I$
是拟同构。这由《同调代数》引理
\ref{homology-lemma-good-right-resolution-gives-qis} 得出。
\end{proof}






\section{导出范畴}
\label{dga-section-derived}

\noindent
回忆，无环微分分次模和微分分次模之间的拟同构这两个概念均有定义
（见第 \ref{dga-section-modules} 节）。

\begin{lemma}
\label{dga-lemma-acyclic}
设 $(A, \text{d})$ 是微分分次代数。由无环模组成的
$K(\text{Mod}_{(A, \text{d})})$ 的满子范畴 $\text{Ac}$，是
$K(\text{Mod}_{(A, \text{d})})$ 的严格满且饱和的三角子范畴。相应的
$K(\text{Mod}_{(A, \text{d})})$ 的饱和乘法系（见《导出范畴》引理
\ref{derived-lemma-operations}）是拟同构类 $\text{Qis}$。特别地，局部化函子
$$
Q : K(\text{Mod}_{(A, \text{d})}) \to
\text{Qis}^{-1}K(\text{Mod}_{(A, \text{d})})
$$
的核是 $\text{Ac}$。此外，函子 $H^0$ 通过 $Q$ 分解。
\end{lemma}

\begin{proof}
由同调长正合列 (\ref{dga-equation-les})，知道 $H^0$ 是同调函子。
$H^0$ 的核是无环对象的子范畴，而在所有 $H^i$ 上诱导同构的箭头正是
拟同构。因此，本引理是《导出范畴》引理
\ref{derived-lemma-acyclic-general} 的特例。

\medskip\noindent
集合论说明。《导出范畴》命题
\ref{derived-proposition-construct-localization} 中的局部化构造，假定给定的
三角范畴是“小”的，即其底层对象族构成一个集合。令 $V_\alpha$ 是包含
$(A, \text{d})$ 的部分宇宙（如《集合》第
\ref{sets-section-sets-hierarchy} 节），并且 $\alpha$ 的共尾数大于
$\aleph_0$（见《集合》命题
\ref{sets-proposition-exist-ordinals-large-cofinality}）。于是可以考虑由
$V_\alpha$ 中的微分分次 $A$-模组成的范畴
$\text{Mod}_{(A, \text{d}), \alpha}$。直接核查可见，命题
\ref{dga-proposition-homotopy-category-triangulated} 的证明中使用的全部构造
都能在 $\text{Mod}_{(A, \text{d}), \alpha}$ 内进行（因为至多只需取微分
分次模的有限直和）。因此得到三角范畴
$\text{Qis}_\alpha^{-1}K(\text{Mod}_{(A, \text{d}), \alpha})$。
下文将看到，若 $\beta>\alpha$，则过渡函子
$$
\text{Qis}_\alpha^{-1}K(\text{Mod}_{(A, \text{d}), \alpha})
\longrightarrow
\text{Qis}_\beta^{-1}K(\text{Mod}_{(A, \text{d}), \beta})
$$
是全忠实的，因为商范畴中的态射集由同伦范畴中从 P-分辨出发的映射计算
（引理 \ref{dga-lemma-resolve} 的证明中构造 P-分辨时，要取可数直和以及由
给定模的子集作指标的直和）。因此，读者应把引理中的范畴视为这些子范畴
的并。
\end{proof}

\noindent
考虑到前一引理证明末尾的集合论说明，按如下方式定义导出范畴。

\begin{definition}
\label{dga-definition-unbounded-derived-category}
设 $(A, \text{d})$ 是微分分次代数。令 $\text{Ac}$ 和 $\text{Qis}$ 如引理
\ref{dga-lemma-acyclic} 所述。$(A, \text{d})$ 的{\it 导出范畴}是三角范畴
$$
D(A, \text{d}) =
K(\text{Mod}_{(A, \text{d})})/\text{Ac} =
\text{Qis}^{-1}K(\text{Mod}_{(A, \text{d})}).
$$
以 $H^0 : D(A, \text{d}) \to \text{Mod}_R$ 表示唯一的函子，使得它与
商函子的复合正好给出上文定义的函子 $H^0$。
\end{definition}

\noindent
下面是此前所称的、计算导出范畴中态射集的引理。

\begin{lemma}
\label{dga-lemma-hom-derived}
设 $(A, \text{d})$ 是微分分次代数。设 $M$ 和 $N$ 是微分分次 $A$-模。
\begin{enumerate}
\item 设 $P \to M$ 是引理 \ref{dga-lemma-resolve} 中的 P-分辨。那么
$$
\Hom_{D(A, \text{d})}(M, N) =
\Hom_{K(\text{Mod}_{(A, \text{d})})}(P, N)
$$
\item 设 $N \to I$ 是引理 \ref{dga-lemma-right-resolution} 中的 I-分辨。那么
$$
\Hom_{D(A, \text{d})}(M, N) =
\Hom_{K(\text{Mod}_{(A, \text{d})})}(M, I)
$$
\end{enumerate}
\end{lemma}

\begin{proof}
设 $P \to M$ 如 (1) 中所述。由于 $P \to M$ 是拟同构，由导出范畴的定义有
$$
\Hom_{D(A, \text{d})}(P, N) = \Hom_{D(A, \text{d})}(M, N)
$$
。$D(A, \text{d})$ 中的态射 $f:P\to N$ 等于 $s^{-1}f'$，其中
$f':P\to N'$ 是态射，而 $s:N\to N'$ 是拟同构。选择有区别三角
$$
N \to N' \to Q \to N[1]
$$
。由于 $s$ 是拟同构，$Q$ 是无环的。因此，由引理
\ref{dga-lemma-property-P-K-projective}，对所有 $k$ 都有
$\Hom_{K(\text{Mod}_{(A, \text{d})})}(P, Q[k]) = 0$。由于
$\Hom_{K(\text{Mod}_{(A, \text{d})})}(P, -)$
是上同调的，故 $f':P\to N'$ 可唯一提升为态射 $f:P\to N$。这就完成了证明。

\medskip\noindent
(2) 的证明是 (1) 的对偶：使用引理 \ref{dga-lemma-property-I-K-injective}
代替引理 \ref{dga-lemma-property-P-K-projective} 即可。
\end{proof}

\begin{lemma}
\label{dga-lemma-derived-products}
设 $(A, \text{d})$ 是微分分次代数。那么
\begin{enumerate}
\item $D(A, \text{d})$ 既有直和，也有直积；
\item 直和可由微分分次模的直和得到；
\item 直积可由微分分次模的直积得到。
\end{enumerate}
\end{lemma}

\begin{proof}
将使用如下事实：$\text{Mod}_{(A, \text{d})}$ 是具有任意直和与直积的
阿贝尔范畴，并且它们在 $K(\text{Mod}_{(A, \text{d})})$ 中产生直和与
直积。见引理 \ref{dga-lemma-dgm-abelian} 和
\ref{dga-lemma-homotopy-direct-sums}。

\medskip\noindent
设 $M_j$ 是一族微分分次 $A$-模。考虑分次直和
$M=\bigoplus M_j$；它是具有显然的……的微分分次 $A$-模。对微分分次
$A$-模 $N$，选取拟同构 $N\to I$，其中 $I$ 是具有性质 (I) 的微分分次
$A$-模；见引理 \ref{dga-lemma-right-resolution}。利用引理
\ref{dga-lemma-hom-derived}，有
\begin{align*}
\Hom_{D(A, \text{d})}(M, N)
& =
\Hom_{K(A, \text{d})}(M, I) \\
& =
\prod \Hom_{K(A, \text{d})}(M_j, I) \\
& =
\prod \Hom_{D(A, \text{d})}(M_j, N)
\end{align*}
由此得到 $D(A, \text{d})$ 中的直和，正如引理第 (2) 项所述。

\medskip\noindent
设 $M_j$ 是一族微分分次 $A$-模。考虑微分分次 $A$-模的直积
$M=\prod M_j$。对微分分次 $A$-模 $N$，选取拟同构 $P\to N$，
其中 $P$ 是具有性质 (P) 的微分分次 $A$-模；见引理
\ref{dga-lemma-resolve}。利用引理 \ref{dga-lemma-hom-derived}，有
\begin{align*}
\Hom_{D(A, \text{d})}(N, M)
& =
\Hom_{K(A, \text{d})}(P, M) \\
& =
\prod \Hom_{K(A, \text{d})}(P, M_j) \\
& =
\prod \Hom_{D(A, \text{d})}(N, M_j)
\end{align*}
由此得到 $D(A, \text{d})$ 中的直和，正如引理第 (3) 项所述。
\end{proof}

\begin{remark}
\label{dga-remark-P-resolution}
设 $R$ 是环，$(A, \text{d})$ 是微分分次 $R$-代数。利用 P-分辨，有时可将
关于 $D(A, \text{d})$ 中一般对象的命题约化为关于 $A[k]$ 的命题。
具体地，设 $T$ 是 $D(A, \text{d})$ 中对象的一个性质，并假设
\begin{enumerate}
\item 若 $K_i$（$i\in I$）是 $D(A, \text{d})$ 中的一族对象，且对所有
$i\in I$ 都有 $T(K_i)$，则 $T(\bigoplus K_i)$ 成立；
\item 若 $K\to L\to M\to K[1]$ 是 $D(A, \text{d})$ 中的有区别三角，
且 $T$ 对其中两个对象成立，则它对第三个对象也成立；以及
\item 对所有 $k\in\mathbf{Z}$，$T(A[k])$ 成立。
\end{enumerate}
那么 $T$ 对 $D(A, \text{d})$ 的所有对象成立。这由引理
\ref{dga-lemma-property-P-sequence} 和 \ref{dga-lemma-resolve} 即得。
\end{remark}







\section{典范 delta-函子}
\label{dga-section-canonical-delta-functor}

\noindent
设 $(A, \text{d})$ 是微分分次代数。考虑函子
$\text{Mod}_{(A, \text{d})} \to K(\text{Mod}_{(A, \text{d})})$。
一般而言，这个函子{\bf 不是} $\delta$-函子。然而，函子
$\text{Mod}_{(A, \text{d})} \to D(A, \text{d})$ 的确是一个
$\delta$-函子。为说明这一点，必须定义与阿贝尔范畴
$\text{Mod}_{(A, \text{d})}$ 中的短正合列
$$
0 \to K \xrightarrow{a} L \xrightarrow{b} M \to 0
$$
相联系的态射 $\delta$。考虑态射 $a$ 的锥 $C(a)$。有
$C(a)=L\oplus K$，并把 $q:C(a)\to M$ 定义为先投影到 $L$，再复合 $b$。
于是得到微分分次 $A$-模同态
$$
q : C(a) \longrightarrow M.
$$
显然 $q\circ i=b$，其中 $i$ 如定义 \ref{dga-definition-cone} 所述。注意，由于
$a$ 是单射，$q$ 的核可与 $\text{id}_K$ 的锥等同，而后者是无环的。
故 $q$ 是拟同构。由引理 \ref{dga-lemma-the-same-up-to-isomorphisms}，三角
$$
(K, L, C(a), a, i, -p)
$$
是 $K(\text{Mod}_{(A, \text{d})})$ 中的有区别三角。由于局部化函子
$K(\text{Mod}_{(A, \text{d})}) \to D(A, \text{d})$ 是正合的，
$(K,L,C(a),a,i,-p)$ 也是 $D(A, \text{d})$ 中的有区别三角。由于 $q$
是拟同构，它在 $D(A, \text{d})$ 中是同构。因此推出
$$
(K, L, M, a, b, -p \circ q^{-1})
$$
是 $D(A, \text{d})$ 中的有区别三角。这引出下述引理。

\begin{lemma}
\label{dga-lemma-derived-canonical-delta-functor}
设 $(A, \text{d})$ 是微分分次代数。上面定义的函子
$\text{Mod}_{(A, \text{d})} \to D(A, \text{d})$ 自然具有
$\delta$-函子结构，其中
$$
\delta_{K \to L \to M} = - p \circ q^{-1}
$$
，而 $p$ 和 $q$ 如上文所述。
\end{lemma}

\begin{proof}
已经看到，对复形的每个短正合列，这个选择都会给出有区别三角。
还须证明该构造的函子性；见《导出范畴》定义
\ref{derived-definition-delta-functor}。稍作推导即可由引理
\ref{dga-lemma-functorial-cone} 得出。可与《导出范畴》引理
\ref{derived-lemma-derived-canonical-delta-functor} 比较。
\end{proof}

\begin{lemma}
\label{dga-lemma-homotopy-colimit}
设 $(A, \text{d})$ 是微分分次代数，$M_n$ 是一个微分分次模系统。那么，
$D(A, \text{d})$ 中的导出余极限 $\text{hocolim} M_n$ 由微分分次模
$\colim M_n$ 表示。
\end{lemma}

\begin{proof}
置 $M=\colim M_n$。由《导出范畴》引理
\ref{derived-lemma-compute-colimit}（应用于其底层阿贝尔群复形），有微分分次
模的正合列
$$
0 \to \bigoplus M_n \to \bigoplus M_n \to M \to 0
$$
。由引理 \ref{dga-lemma-derived-products}，这些直和也是 $D(\mathcal{A})$ 中的
直和。因此，结论由《导出范畴》定义
\ref{derived-definition-derived-colimit} 中导出余极限的定义，以及复形的短正合列
给出有区别三角这一事实（引理
\ref{dga-lemma-derived-canonical-delta-functor}）得出。
\end{proof}








\section{线性范畴}
\label{dga-section-linear}

\noindent
这里只给出定义。

\begin{definition}
\label{dga-definition-linear-category}
设 $R$ 是环。{\it $R$-线性范畴 $\mathcal{A}$} 是这样的范畴：每个态射集
都带有 $R$-模结构，并且对 $x,y,z\in\Ob(\mathcal{A})$，复合律
$$
\Hom_\mathcal{A}(y, z) \times \Hom_\mathcal{A}(x, y)
\longrightarrow
\Hom_\mathcal{A}(x, z)
$$
是 $R$-双线性的。
\end{definition}

\noindent
因此，复合确定 $R$-模的 $R$-线性映射
$$
\Hom_\mathcal{A}(y, z) \otimes_R \Hom_\mathcal{A}(x, y)
\longrightarrow
\Hom_\mathcal{A}(x, z)
$$
。注意，这里并不假设 $R$-线性范畴是加性范畴。

\begin{definition}
\label{dga-definition-functor-linear-categories}
设 $R$ 是环。{\it $R$-线性范畴之间的函子}，也称为
{\it $R$-线性函子}，是函子 $F:\mathcal{A}\to\mathcal{B}$，使得对
$\mathcal{A}$ 的所有对象 $x,y$，映射
$F : \Hom_\mathcal{A}(x, y) \to \Hom_\mathcal{B}(F(x), F(y))$
是 $R$-模同态。
\end{definition}







\section{分次范畴}
\label{dga-section-graded}

\noindent
这里只给出一些定义。

\begin{definition}
\label{dga-definition-graded-category}
设 $R$ 是环。{\it $R$ 上的分次范畴 $\mathcal{A}$} 是这样的范畴：
每个态射集都带有分次 $R$-模结构，并且对
$x,y,z\in\Ob(\mathcal{A})$，复合是 $R$-双线性的，并诱导分次 $R$-模同态
$$
\Hom_\mathcal{A}(y, z) \otimes_R \Hom_\mathcal{A}(x, y)
\longrightarrow
\Hom_\mathcal{A}(x, z)
$$
（即保持次数的同态）。
\end{definition}

\noindent
在这种情形下，以 $\Hom_\mathcal{A}^i(x,y)$ 表示分次对象
$\Hom_\mathcal{A}(x,y)$ 的次数 $i$ 部分，从而
$$
\Hom_\mathcal{A}(x, y) =
\bigoplus\nolimits_{i \in \mathbf{Z}} \Hom_\mathcal{A}^i(x, y)
$$
是它按分次部分的直和分解。

\begin{definition}
\label{dga-definition-functor-graded-categories}
设 $R$ 是环。{\it $R$ 上分次范畴之间的函子}，也称为
{\it 分次函子}，是函子 $F:\mathcal{A}\to\mathcal{B}$，使得对
$\mathcal{A}$ 的所有对象 $x,y$，映射
$F : \Hom_\mathcal{A}(x, y) \to \Hom_\mathcal{A}(F(x), F(y))$
是分次 $R$-模同态。
\end{definition}

\noindent
给定一个分次范畴，往往关心其对应的、态射次数为 $0$ 的“通常”范畴。
正式定义如下。

\begin{definition}
\label{dga-definition-H0-of-graded-category}
设 $R$ 是环，$\mathcal{A}$ 是 $R$ 上的分次范畴。令
{\it $\mathcal{A}^0$} 为与 $\mathcal{A}$ 具有相同对象的范畴，并置
$$
\Hom_{\mathcal{A}^0}(x, y) = \Hom^0_\mathcal{A}(x, y)
$$
，即取 $\mathcal{A}$ 的态射分次模的次数 $0$ 部分。
\end{definition}

\begin{definition}
\label{dga-definition-graded-direct-sum}
设 $R$ 是环，$\mathcal{A}$ 是 $R$ 上的分次范畴。若 $\mathcal{A}$ 中的直和
$(x,y,z,i,j,p,q)$（记号同《同调代数》评注
\ref{homology-remark-direct-sum}）中的 $i,j,p,q$ 都是次数 $0$ 的齐次态射，
则称它为{\it 分次直和}。
\end{definition}

\begin{example}[分次对象的分次范畴]
\label{dga-example-graded-category-graded-objects}
设 $\mathcal{B}$ 是加性范畴。回忆，在《同调代数》定义
\ref{homology-definition-graded} 中已经定义了 $\mathcal{B}$ 的分次对象范畴
$\text{Gr}(\mathcal{B})$。本例将构造 $R=\mathbf{Z}$ 上的分次范畴
$\text{Gr}^{gr}(\mathcal{B})$，使其对应范畴
$\text{Gr}^{gr}(\mathcal{B})^0$ 恢复 $\text{Gr}(\mathcal{B})$。
取 $\mathcal{B}$ 的分次对象作为 $\text{Gr}^{gr}(\mathcal{B})$ 的对象。
给定 $\mathcal{B}$ 的分次对象 $A=(A^i)$ 和 $B=(B^i)$，置
$$
\Hom_{\text{Gr}^{gr}(\mathcal{B})}(A, B) =
\bigoplus\nolimits_{n \in \mathbf{Z}} \Hom^n(A, B)
$$
，其中次数 $n$ 的分次部分是从 $A$ 到 $B$ 的次数 $n$ 齐次映射所组成的
阿贝尔群。具体地，
$$
\Hom^n(A, B) = \prod\nolimits_{p + q = n} \Hom_\mathcal{B}(A^{-q}, B^p)
$$
（注意指标的反转，并注意这里是直积而不是直和）。换言之，从 $A$ 到 $B$
的次数 $n$ 态射 $f$ 可视为系统 $f=(f_{p,q})$，其中
$p,q\in\mathbf{Z}$、$p+q=n$，且 $f_{p,q}:A^{-q}\to B^p$ 是
$\mathcal{B}$ 的态射。给定 $\mathcal{B}$ 的分次对象 $A,B,C$，
$\text{Gr}^{gr}(\mathcal{B})$ 中态射的复合通过映射
$$
\Hom^m(B, C) \times \Hom^n(A, B) \longrightarrow \Hom^{n + m}(A, C)
$$
定义，即取分次对象齐次映射的普通复合 $(g,f)\mapsto g\circ f$。按分量写为
$$
(g \circ f)_{p, r} = g_{p, q} \circ f_{-q, r}
$$
，其中 $q$ 满足 $p+q=m$ 和 $-q+r=n$。
\end{example}

\begin{example}[分次模的分次范畴]
\label{dga-example-gm-gr-cat}
设 $A$ 是环 $R$ 上的 $\mathbf{Z}$-分次代数。构造 $R$ 上的分次范畴
$\text{Mod}^{gr}_A$，使其对应范畴 $(\text{Mod}^{gr}_A)^0$ 是分次
$A$-模范畴。取右分次 $A$-模作为 $\text{Mod}^{gr}_A$ 的对象（见第
\ref{dga-section-projectives-over-algebras} 节）。给定分次 $A$-模 $L$ 和 $M$，置
$$
\Hom_{\text{Mod}^{gr}_A}(L, M) =
\bigoplus\nolimits_{n \in \mathbf{Z}} \Hom^n(L, M)
$$
，其中 $\Hom^n(L,M)$ 是次数 $n$ 齐次的右 $A$-模映射 $L\to M$ 的集合，
即对所有 $i\in\mathbf{Z}$ 都有 $f(L^i)\subset M^{i+n}$。按分量写，
$$
\Hom^n(L, M) \subset \prod\nolimits_{p + q = n} \Hom_R(L^{-q}, M^p)
$$
（注意指标的反转）是由满足下式的 $f=(f_{p,q})$ 组成的子集：
$$
f_{p, q}(m a) = f_{p - i, q + i}(m)a
$$
，其中 $a\in A^i$、$m\in L^{-q-i}$。对分次 $A$-模 $K,L,M$，通过映射
$$
\Hom^m(L, M) \times \Hom^n(K, L) \longrightarrow \Hom^{n + m}(K, M)
$$
定义 $\text{Mod}^{gr}_A$ 中的复合，即取右 $A$-模映射的普通复合
$(g,f)\mapsto g\circ f$。
\end{example}

\begin{remark}
\label{dga-remark-graded-shift-functors}
设 $R$ 是环，$\mathcal{D}$ 是 $R$-线性范畴，并带有由
$n\in\mathbf{Z}$ 作指标的一族 $R$-线性函子
$[n]:\mathcal{D}\to\mathcal{D}$、$x\mapsto x[n]$，使得
$[n]\circ[m]=[n+m]$ 且 $[0]=\text{id}_\mathcal{D}$（作为函子相等）。
由此可以构造 $R$ 上与 $\mathcal{D}$ 具有相同对象的分次范畴
$\mathcal{D}^{gr}$，置
$$
\Hom_{\mathcal{D}^{gr}}(x, y) =
\bigoplus\nolimits_{n \in \mathbf{Z}} \Hom_\mathcal{D}(x, y[n])
$$
，其中 $x,y$ 是 $\mathcal{D}$ 的对象。注意
$(\mathcal{D}^{gr})^0=\mathcal{D}$（见定义
\ref{dga-definition-H0-of-graded-category}）。此外，分次范畴
$\mathcal{D}^{gr}$ 继承 $R$-线性分次函子 $[n]$，满足
$[n]\circ[m]=[n+m]$ 和 $[0]=\text{id}_{\mathcal{D}^{gr}}$，并具有性质
$$
\Hom_{\mathcal{D}^{gr}}(x, y[n]) = \Hom_{\mathcal{D}^{gr}}(x, y)[n]
$$
；这是与态射复合相容的分次 $R$-模同构。

\medskip\noindent
反过来，设给定 $R$ 上的分次范畴 $\mathcal{A}$，它带有一族 $R$-线性
分次函子 $[n]$，满足 $[n]\circ[m]=[n+m]$ 和
$[0]=\text{id}_\mathcal{A}$，并且还带有同构
$$
\Hom_\mathcal{A}(x, y[n]) = \Hom_\mathcal{A}(x, y)[n]
$$
；这是与态射复合相容的分次 $R$-模同构。读者容易证明
$\mathcal{A}=(\mathcal{A}^0)^{gr}$。

\medskip\noindent
下面是上述关系 $\mathcal{D}\leftrightarrow\mathcal{A}$ 的两个例子：
\begin{enumerate}
\item 设 $\mathcal{B}$ 是加性范畴。若
$\mathcal{D}=\text{Gr}(\mathcal{B})$，则
$\mathcal{A}=\text{Gr}^{gr}(\mathcal{B})$，如例
\ref{dga-example-graded-category-graded-objects} 所述。
\item 若 $A$ 是分次环，且 $\mathcal{D}=\text{Mod}_A$ 是分次右 $A$-模
范畴，则 $\mathcal{A}=\text{Mod}^{gr}_A$；见例
\ref{dga-example-gm-gr-cat}。
\end{enumerate}
\end{remark}






\section{微分分次范畴}
\label{dga-section-dga-categories}

\noindent
注意，若 $R$ 是环，则 $R$ 是自身之上的微分分次代数（当然取
$R=R^0$）。此时，微分分次 $R$-模与 $R$-模复形是一回事。特别地，
给定两个微分分次 $R$-模 $M$ 和 $N$，以 $M\otimes_R N$ 表示如下微分
分次 $R$-模：先取与 $M,N$ 对应的 $R$-模复形的张量积所得到的双复形，
再取其关联全复形。

\begin{definition}
\label{dga-definition-dga-category}
设 $R$ 是环。{\it $R$ 上的微分分次范畴 $\mathcal{A}$} 是这样的范畴：
每个态射集都带有微分分次 $R$-模结构，并且对
$x,y,z\in\Ob(\mathcal{A})$，复合是 $R$-双线性的，并诱导微分分次
$R$-模同态
$$
\Hom_\mathcal{A}(y, z) \otimes_R \Hom_\mathcal{A}(x, y)
\longrightarrow
\Hom_\mathcal{A}(x, z)
$$
。
\end{definition}

\noindent
定义中的最后一个条件表示：若
$f\in\Hom_\mathcal{A}^n(x,y)$ 和 $g\in\Hom_\mathcal{A}^m(y,z)$
分别是次数 $n$ 和 $m$ 的齐次态射，则在
$\Hom_\mathcal{A}^{n+m+1}(x,z)$ 中有
$$
\text{d}(g \circ f) = \text{d}(g) \circ f + (-1)^mg \circ \text{d}(f)
$$
。这由双复形全复形上的微分符号法则得出；见《同调代数》定义
\ref{homology-definition-associated-simple-complex}。

\begin{definition}
\label{dga-definition-functor-dga-categories}
设 $R$ 是环。{\it $R$ 上微分分次范畴之间的函子}是函子
$F:\mathcal{A}\to\mathcal{B}$，使得对 $\mathcal{A}$ 的所有对象 $x,y$，
映射
$F : \Hom_\mathcal{A}(x, y) \to \Hom_\mathcal{A}(F(x), F(y))$
是微分分次 $R$-模同态。
\end{definition}

\noindent
给定一个微分分次范畴，往往关心其对应的复形范畴和同伦范畴。正式定义如下。

\begin{definition}
\label{dga-definition-homotopy-category-of-dga-category}
设 $R$ 是环，$\mathcal{A}$ 是 $R$ 上的微分分次范畴。定义
\begin{enumerate}
\item {\it $\mathcal{A}$ 的复形范畴}\footnote{这可能不是标准术语。}
为范畴 $\text{Comp}(\mathcal{A})$；它与 $\mathcal{A}$ 具有相同对象，并且
$$
\Hom_{\text{Comp}(\mathcal{A})}(x, y) =
\Ker(d : \Hom^0_\mathcal{A}(x, y) \to \Hom^1_\mathcal{A}(x, y))
$$
\item {\it $\mathcal{A}$ 的同伦范畴}为范畴 $K(\mathcal{A})$；它与
$\mathcal{A}$ 具有相同对象，并且
$$
\Hom_{K(\mathcal{A})}(x, y) = H^0(\Hom_\mathcal{A}(x, y))
$$
\end{enumerate}
\end{definition}

\noindent
这里对符号 $K(\mathcal{A})$ 的用法并不标准，但至少与 Stacks 项目其他章中
$K(-)$ 的用法相容。

\begin{definition}
\label{dga-definition-dg-direct-sum}
设 $R$ 是环，$\mathcal{A}$ 是 $R$ 上的微分分次范畴。若 $\mathcal{A}$
中的直和 $(x,y,z,i,j,p,q)$（记号同《同调代数》评注
\ref{homology-remark-direct-sum}）中的 $i,j,p,q$ 都是次数 $0$ 的齐次闭态射，
即 $\text{d}(i)=0$ 等，则称它为{\it 微分分次直和}。
\end{definition}

\begin{lemma}
\label{dga-lemma-functorial}
设 $R$ 是环。$R$ 上微分分次范畴之间的函子
$F:\mathcal{A}\to\mathcal{B}$ 诱导函子
$\text{Comp}(\mathcal{A}) \to \text{Comp}(\mathcal{B})$
和 $K(\mathcal{A})\to K(\mathcal{B})$。
\end{lemma}

\begin{proof}
略。
\end{proof}

\begin{example}[复形的微分分次范畴]
\label{dga-example-category-complexes}
设 $\mathcal{B}$ 是加性范畴。构造 $R=\mathbf{Z}$ 上的微分分次范畴
$\text{Comp}^{dg}(\mathcal{B})$，使其对应的复形范畴是
$\text{Comp}(\mathcal{B})$，对应的同伦范畴是 $K(\mathcal{B})$。
取 $\mathcal{B}$ 的复形作为 $\text{Comp}^{dg}(\mathcal{B})$ 的对象。
给定 $\mathcal{B}$ 的复形 $A^\bullet$ 和 $B^\bullet$，有时也以
$A^\bullet$ 和 $B^\bullet$ 表示 $\mathcal{B}$ 中相应的分次对象
（即忘掉微分）。采用这一记号上的滥用，置
$$
\Hom_{\text{Comp}^{dg}(\mathcal{B})}(A^\bullet, B^\bullet) =
\Hom_{\text{Gr}^{gr}(\mathcal{B})}(A^\bullet, B^\bullet) =
\bigoplus\nolimits_{n \in \mathbf{Z}} \Hom^n(A, B)
$$
，把它视为分次 $\mathbf{Z}$-模，记号和定义同例
\ref{dga-example-graded-category-graded-objects}。换言之，第 $n$ 个分次部分是
分次对象之间次数 $n$ 的齐次态射所组成的阿贝尔群
$$
\Hom^n(A^\bullet, B^\bullet) =
\prod\nolimits_{p + q = n} \Hom_\mathcal{B}(A^{-q}, B^p)
$$
。注意指标的反转，并注意这里是直积而不是直和。对次数为 $n$ 的元素
$f\in\Hom^n(A^\bullet,B^\bullet)$，置
$$
\text{d}(f) = \text{d}_B \circ f - (-1)^n f \circ \text{d}_A
$$
。符号与《代数进阶》第 \ref{more-algebra-section-sign-rules} 节完全相同。
为使此式有意义，把 $\text{d}_B$ 和 $\text{d}_A$ 看作 $\mathcal{B}$ 的
分次对象之间次数 $1$ 的齐次映射，并在右侧使用范畴
$\text{Gr}^{gr}(\mathcal{B})$ 中的复合。按分量写，若
$f=(f_{p,q})$ 且 $f_{p,q}:A^{-q}\to B^p$，则
\begin{equation}
\label{dga-equation-differential-hom-complex}
\text{d}(f_{p, q}) =
\text{d}_B \circ f_{p, q} - (-1)^{p + q} f_{p, q} \circ \text{d}_A 
\end{equation}
注意，此表达式的第一项属于 $\Hom_\mathcal{B}(A^{-q},B^{p+1})$，
第二项属于 $\Hom_\mathcal{B}(A^{-q-1},B^p)$。读者可以验证：
\begin{enumerate}
\item $\text{d}$ 的平方为零；
\item $\Hom^n(A^\bullet,B^\bullet)$ 中的元素 $f$ 满足
$\text{d}(f)=0$，当且仅当 $\mathcal{B}$ 的分次对象之间的态射
$f:A^\bullet\to B^\bullet[n]$ 实际上是复形映射；
\item 特别地，$\text{Comp}^{dg}(\mathcal{B})$ 的复形范畴等于
$\text{Comp}(\mathcal{B})$；
\item (2) 中由 $f$ 定义的复形态射与零同伦，当且仅当存在
$g\in\Hom^{n-1}(A^\bullet,B^\bullet)$ 使得 $f=\text{d}(g)$；
\item 特别地，得到典范同构
$$
\Hom_{K(\mathcal{B})}(A^\bullet, B^\bullet)
\longrightarrow
H^0(\Hom_{\text{Comp}^{dg}(\mathcal{B})}(A^\bullet, B^\bullet))
$$
，并且 $\text{Comp}^{dg}(\mathcal{B})$ 的同伦范畴等于 $K(\mathcal{B})$。
\end{enumerate}
给定复形 $A^\bullet,B^\bullet,C^\bullet$，通过分次范畴
$\text{Gr}^{gr}(\mathcal{B})$ 中的复合 $(g,f)\mapsto g\circ f$ 定义复合
$$
\Hom^m(B^\bullet, C^\bullet) \times \Hom^n(A^\bullet, B^\bullet)
\longrightarrow
\Hom^{n + m}(A^\bullet, C^\bullet)
$$
；见例 \ref{dga-example-graded-category-graded-objects}。这定义微分分次模的映射
$$
\Hom_{\text{Comp}^{dg}(\mathcal{B})}(B^\bullet, C^\bullet)
\otimes_R
\Hom_{\text{Comp}^{dg}(\mathcal{B})}(A^\bullet, B^\bullet)
\longrightarrow
\Hom_{\text{Comp}^{dg}(\mathcal{B})}(A^\bullet, C^\bullet)
$$
，正如定义 \ref{dga-definition-dga-category} 所要求的，因为
\begin{align*}
\text{d}(g \circ f) & =
\text{d}_C \circ g \circ f - (-1)^{n + m} g \circ f \circ \text{d}_A \\
& =
\left(\text{d}_C \circ g - (-1)^m g \circ \text{d}_B\right) \circ f +
(-1)^m g \circ \left(\text{d}_B \circ f - (-1)^n f \circ \text{d}_A\right) \\
& =
\text{d}(g) \circ f + (-1)^m g \circ \text{d}(f)
\end{align*}
，这正是所需等式。
\end{example}

\begin{lemma}
\label{dga-lemma-additive-functor-induces-dga-functor}
设 $F:\mathcal{B}\to\mathcal{B}'$ 是加性范畴之间的加性函子。那么 $F$
诱导例 \ref{dga-example-category-complexes} 中微分分次范畴之间的函子
$$
F : \text{Comp}^{dg}(\mathcal{B}) \to \text{Comp}^{dg}(\mathcal{B}')
$$
，并在复形范畴和同伦范畴上诱导通常的函子。
\end{lemma}

\begin{proof}
略。
\end{proof}

\begin{example}[微分分次模的微分分次范畴]
\label{dga-example-dgm-dg-cat}
设 $(A, \text{d})$ 是环 $R$ 上的微分分次代数。构造 $R$ 上的微分分次
范畴 $\text{Mod}^{dg}_{(A, \text{d})}$，使其复形范畴是
$\text{Mod}_{(A, \text{d})}$，同伦范畴是
$K(\text{Mod}_{(A, \text{d})})$。取微分分次 $A$-模作为
$\text{Mod}^{dg}_{(A, \text{d})}$ 的对象。给定微分分次 $A$-模 $L,M$，置
$$
\Hom_{\text{Mod}^{dg}_{(A, \text{d})}}(L, M) =
\Hom_{\text{Mod}^{gr}_A}(L, M) = \bigoplus \Hom^n(L, M)
$$
，把它视为分次 $R$-模，其中右侧如例 \ref{dga-example-gm-gr-cat} 所定义。
换言之，第 $n$ 个分次部分 $\Hom^n(L,M)$ 是次数 $n$ 齐次的右 $A$-模映射
所组成的 $R$-模。对元素 $f\in\Hom^n(L,M)$，置
$$
\text{d}(f) = \text{d}_M \circ f - (-1)^n f \circ \text{d}_L
$$
。为使此式有意义，把 $\text{d}_M$ 和 $\text{d}_L$ 看作分次 $R$-模映射，
并采用分次 $R$-模映射的复合。显然，作为分次 $R$-模映射，
$\text{d}(f)$ 是次数 $n+1$ 的齐次映射；它还是 $A$-线性的，因为
\begin{align*}
\text{d}(f)(xa)
& =
\text{d}_M(f(x) a) - (-1)^n f (\text{d}_L(xa)) \\
& =
\text{d}_M(f(x)) a + (-1)^{\deg(x) + n} f(x) \text{d}(a) 
- (-1)^n f(\text{d}_L(x)) a - (-1)^{n + \deg(x)} f(x) \text{d}(a) \\
& = \text{d}(f)(x) a
\end{align*}
，这正是所需等式（注意，若 $\Hom$ 上微分的定义中没有这个符号，计算就
不能成立）。按分量表示 $f$ 的微分时，也有与例
\ref{dga-example-category-complexes} 类似的公式。读者可以仿照例
\ref{dga-example-category-complexes} 验证：
\begin{enumerate}
\item $\text{d}$ 的平方为零；
\item $\Hom^n(L,M)$ 中的元素 $f$ 满足 $\text{d}(f)=0$，当且仅当
$f:L\to M[n]$ 是微分分次 $A$-模同态；
\item 特别地，$\text{Mod}^{dg}_{(A, \text{d})}$ 的复形范畴是
$\text{Mod}_{(A, \text{d})}$；
\item (2) 中由 $f$ 定义的同态与零同伦，当且仅当存在
$g\in\Hom^{n-1}(L,M)$ 使得 $f=\text{d}(g)$；
\item 特别地，得到典范同构
$$
\Hom_{K(\text{Mod}_{(A, \text{d})})}(L, M)
\longrightarrow
H^0(\Hom_{\text{Mod}^{dg}_{(A, \text{d})}}(L, M))
$$
，并且 $\text{Mod}^{dg}_{(A, \text{d})}$ 的同伦范畴是
$K(\text{Mod}_{(A, \text{d})})$。
\end{enumerate}
给定微分分次 $A$-模 $K,L,M$，通过齐次右 $A$-模映射的复合
$(g,f)\mapsto g\circ f$ 定义复合
$$
\Hom^m(L, M) \times \Hom^n(K, L) \longrightarrow \Hom^{n + m}(K, M)
$$
。这定义微分分次模的映射
$$
\Hom_{\text{Mod}^{dg}_{(A, \text{d})}}(L, M) \otimes_R
\Hom_{\text{Mod}^{dg}_{(A, \text{d})}}(K, L) \longrightarrow
\Hom_{\text{Mod}^{dg}_{(A, \text{d})}}(K, M)
$$
，正如定义 \ref{dga-definition-dga-category} 所要求的，因为
\begin{align*}
\text{d}(g \circ f) & =
\text{d}_M \circ g \circ f - (-1)^{n + m} g \circ f \circ \text{d}_K \\
& =
\left(\text{d}_M \circ g - (-1)^m g \circ \text{d}_L\right) \circ f +
(-1)^m g \circ \left(\text{d}_L \circ f - (-1)^n f \circ \text{d}_K\right) \\
& =
\text{d}(g) \circ f + (-1)^m g \circ \text{d}(f)
\end{align*}
，这正是所需等式。
\end{example}

\begin{lemma}
\label{dga-lemma-homomorphism-induces-dga-functor}
设 $\varphi:(A,\text{d})\to(E,\text{d})$ 是微分分次代数同态。那么
$\varphi$ 诱导例 \ref{dga-example-dgm-dg-cat} 中微分分次范畴之间的函子
$$
F :
\text{Mod}^{dg}_{(E, \text{d})}
\longrightarrow
\text{Mod}^{dg}_{(A, \text{d})}
$$
，并在微分分次模范畴和同伦范畴上诱导显然的限制函子。
\end{lemma}

\begin{proof}
略。
\end{proof}

\begin{lemma}
\label{dga-lemma-construction}
设 $R$ 是环，$\mathcal{A}$ 是 $R$ 上的微分分次范畴，$x$ 是
$\mathcal{A}$ 的对象。令
$$
(E, \text{d}) = \Hom_\mathcal{A}(x, x)
$$
为 $x$ 的自同态微分分次 $R$-代数。令 $E$ 通过 $\mathcal{A}$ 中的复合作用
在 $\Hom_\mathcal{A}(x,y)$ 上，得到微分分次范畴之间的函子
$$
\mathcal{A} \longrightarrow \text{Mod}^{dg}_{(E, \text{d})},\quad
y \longmapsto \Hom_\mathcal{A}(x, y)
$$
。应用引理 \ref{dga-lemma-functorial}，此函子诱导函子
$$
\text{Comp}(\mathcal{A}) \to \text{Mod}_{(A, \text{d})}
\quad\text{且}\quad
K(\mathcal{A}) \to K(\text{Mod}_{(A, \text{d})})
$$
。
\end{lemma}

\begin{proof}
本引理不证自明。
\end{proof}









\section{得到三角范畴}
\label{dga-section-review}

\noindent
本节讨论《导出范畴》命题
\ref{derived-proposition-homotopy-category-triangulated} 和命题
\ref{dga-proposition-homotopy-category-triangulated} 的证明论证所适用的最一般框架。

\medskip\noindent
设 $R$ 是环，$\mathcal{A}$ 是 $R$ 上的微分分次范畴。为使论证成立，
对 $\mathcal{A}$ 施加下列公理：
\begin{enumerate}
\item[(A)] $\mathcal{A}$ 有零对象以及两个对象的微分分次直和
（如定义 \ref{dga-definition-dg-direct-sum}）。
\item[(B)] 存在微分分次范畴之间的函子
$[n]:\mathcal{A}\longrightarrow\mathcal{A}$，使得
$[0]=\text{id}_\mathcal{A}$ 且 $[n+m]=[n]\circ[m]$，并且给定与复合相容的
微分分次 $R$-模同构
$$
\Hom_\mathcal{A}(x, y[n]) = \Hom_\mathcal{A}(x, y)[n]
$$
。
\end{enumerate}

\noindent
对上述微分分次范畴 $\mathcal{A}$，规定：
\begin{enumerate}
\item 若在底层分次范畴中存在同构 $y\cong x\oplus z$，使得
$x\to z$ 和 $y\to z$ 是（余）投影，则称 $\text{Comp}(\mathcal{A})$
中的态射列 $x\to y\to z$ 为{\it 容许短正合列}。
\item 若 $\text{Comp}(\mathcal{A})$ 中的态射 $x\to y$ 可延拓为容许短正合列
$x\to y\to z$，则称它为{\it 容许单态射}。
\item 若 $\text{Comp}(\mathcal{A})$ 中的态射 $y\to z$ 可延拓为容许短正合列
$x\to y\to z$，则称它为{\it 容许满态射}。
\end{enumerate}
下一个引理说明，在公理 (A) 和 (B) 成立时，容许短正合列给出一个三角。

\begin{lemma}
\label{dga-lemma-get-triangle}
设 $\mathcal{A}$ 是满足公理 (A)、(B) 的微分分次范畴。给定容许短正合列
$x\to y\to z$，得到（见证明）$\text{Comp}(\mathcal{A})$ 中的三角
$$
x \to y \to z \to x[1]
$$
，其性质为：$z[-1]\to x\to y\to z\to x[1]$ 中任意两个相邻箭头的复合
在 $K(\mathcal{A})$ 中为零。
\end{lemma}

\begin{proof}
选择图表
$$
\xymatrix{
x \ar[rr]_1 \ar[rd]_a & & x \\
& y \ar[ru]_\pi \ar[rd]^b & \\
z \ar[rr]^1 \ar[ru]^s & & z
}
$$
，它给出容许短正合列定义中的分次对象同构 $y\cong x\oplus z$。在此情形下，
下列等式成立：
\begin{enumerate}
\item $1=\pi a$，从而 $\text{d}(\pi)a=0$；
\item $1=bs$，从而 $b\text{d}(s)=0$；
\item $1=a\pi+sb$，从而 $a\text{d}(\pi)+\text{d}(s)b=0$；
\item $\pi s=0$，从而 $\text{d}(\pi)s+\pi\text{d}(s)=0$；
\item $\text{d}(s)=a\pi\text{d}(s)$，因为
$\text{d}(s)=(a\pi+sb)\text{d}(s)$ 且 $b\text{d}(s)=0$；
\item $\text{d}(\pi)=\text{d}(\pi)sb$，因为
$\text{d}(\pi)=\text{d}(\pi)(a\pi+sb)$ 且 $\text{d}(\pi)a=0$；
\item $\text{d}(\pi\text{d}(s))=0$，因为将它后复合单态射 $a$，得到
$\text{d}(a\pi\text{d}(s))=\text{d}(\text{d}(s))=0$；以及
\item $\text{d}(\text{d}(\pi)s)=0$，因为由 (4)，它是
$\text{d}(\pi\text{d}(s))$ 的负值，而后者由 (7) 为 $0$。
\end{enumerate}
这里反复使用了 $\text{d}(a)=0$、$\text{d}(b)=0$ 和 $\text{d}(1)=0$。
由 (7)，
$$
\delta = \pi \text{d}(s) = - \text{d}(\pi) s : z \to x[1]
$$
是 $\text{Comp}(\mathcal{A})$ 中的态射。由 (5)，复合
$a\delta=a\pi\text{d}(s)=\text{d}(s)$ 与零同伦。由 (6)，复合
$\delta b=-\text{d}(\pi)sb=\text{d}(-\pi)$ 与零同伦。
\end{proof}

\noindent
除公理 (A)、(B) 外，还需要一个关于锥存在性的公理。将全部条件正式写成如下情形。

\begin{situation}
\label{dga-situation-ABC}
这里 $R$ 是环，$\mathcal{A}$ 是 $R$ 上满足公理 (A)、(B) 以及下列公理的
微分分次范畴：
\begin{enumerate}
\item[(C)] 给定次数为 $0$ 且满足 $\text{d}(f)=0$ 的箭头 $f:x\to y$，
$\text{Comp}(\mathcal{A})$ 中存在容许短正合列 $y\to c(f)\to x[1]$，
使引理 \ref{dga-lemma-get-triangle} 中的映射 $x[1]\to y[1]$ 等于 $f[1]$。
\end{enumerate}
\end{situation}

\noindent
称 $c(f)$ 为态射 $f$ 的一个{\it 锥}。若 (A)、(B)、(C) 成立，
则锥在一种较弱的意义下具有函子性。

\begin{lemma}
\label{dga-lemma-cone}
\begin{slogan}
同伦范畴是三角范畴。本引理证明三角范畴公理的一部分。
\end{slogan}
在情形 \ref{dga-situation-ABC} 中，设图表
$$
\xymatrix{
x_1 \ar[r]_{f_1} \ar[d]_a & y_1 \ar[d]^b \\
x_2 \ar[r]^{f_2} & y_2
}
$$
在 $\text{Comp}(\mathcal{A})$ 中在同伦意义下交换。则存在态射
$c:c(f_1)\to c(f_2)$，它在 $K(\mathcal{A})$ 中给出三角态射
$$
(a, b, c) : (x_1, y_1, c(f_1)) \to (x_1, y_1, c(f_1))
$$
。
\end{lemma}

\begin{proof}
假设意味着存在次数为 $-1$ 的态射 $h:x_1\to y_2$，使得
$\text{d}(h)=bf_1-f_2a$。选择与态射 $y_i\to c(f_i)\to x_i[1]$ 相容的
分次对象同构 $c(f_i)=y_i\oplus x_i[1]$。以
$a_i:y_i\to c(f_i)$、$b_i:c(f_i)\to x_i[1]$、
$s_i:x_i[1]\to c(f_i)$ 和 $\pi_i:c(f_i)\to y_i$ 表示这些给定态射。
回忆，$x_i[1]\to y_i[1]$ 由 $\pi_i\text{d}(s_i)$ 给出。由公理 (C)，这意味着
$$
f_i = \pi_i \text{d}(s_i) = - \text{d}(\pi_i) s_i
$$
（利用移位函子 $[1]$ 将 $\Hom(x_i,y_i)$ 与
$\Hom(x_i[1],y_i[1])$ 识别）。置
$c=a_2b\pi_1+s_2ab_1+a_2hb$。利用引理 \ref{dga-lemma-get-triangle}
证明中得到的等式，有
\begin{align*}
\text{d}(c)
& =
a_2 b \text{d}(\pi_1) + \text{d}(s_2) a b_1 + a_2 \text{d}(h) b_1 \\
& =
- a_2 b f_1 b_1 + a_2 f_2 a b_1 + a_2 (b f_1 - f_2 a) b_1 \\
& = 0
\end{align*}
（这里特别使用了
$\text{d}(\pi_1) = \text{d}(\pi_1) s_1 b_1 = f_1 b_1$ 以及
$\text{d}(s_2)=a_2\pi_2\text{d}(s_2)=a_2f_2$）。因此，$c$ 是
$\mathcal{A}$ 中次数为 $0$ 的态射 $c:c(f_1)\to c(f_2)$，并与给定态射
$y_i\to c(f_i)\to x_i[1]$ 相容。
\end{proof}

\noindent
在情形 \ref{dga-situation-ABC} 中，若存在容许短正合列 $x'\to y'\to z'$，
使得 $K(\mathcal{A})$ 中的三角 $(x,y,z,f,g,h)$ 作为三角同构于引理
\ref{dga-lemma-get-triangle} 构造的三角
$(x',y',z',x'\to y',y'\to z',\delta)$，则称它为{\it 有区别三角}。
下文将证明
$$
\boxed{
K(\mathcal{A})\text{ 是三角范畴}
}
$$
这个结果虽没有想象中那么一般，却适用于 Stacks 项目迄今所讨论情形的若干
自然推广。下面给出一些例子：
\begin{enumerate}
\item 设 $(X,\mathcal{O}_X)$ 是环化空间，$(A,d)$ 是微分分次
$\mathcal{O}_X$-代数层，$\mathcal{A}$ 是微分分次 $A$-模的微分分次范畴。
那么 $K(\mathcal{A})$ 是三角范畴。
\item 设 $(\mathcal{C},\mathcal{O})$ 是环化位点，$(A,d)$ 是微分分次
$\mathcal{O}$-代数层，$\mathcal{A}$ 是微分分次 $A$-模的微分分次范畴。
那么 $K(\mathcal{A})$ 是三角范畴。见《微分分次层》命题
\ref{sdga-proposition-homotopy-category-triangulated}。
\item 另外两个风格不同的例子见《例》第
\ref{examples-section-nongraded-differential-graded} 节。
\end{enumerate}

\noindent
下面这个简单引理是该构造的关键。

\begin{lemma}
\label{dga-lemma-id-cone-null}
在情形 \ref{dga-situation-ABC} 中，给定 $\mathcal{A}$ 的任意对象 $x$ 以及
恒等态射 $1_x:x\to x$ 的锥 $C(1_x)$，$C(1_x)$ 上的恒等态射与零同伦。
\end{lemma}

\begin{proof}
考虑公理 (C) 给出的容许短正合列
$$
\xymatrix{
x \ar@<0.5ex>[r]^a  &
C(1_x) \ar@<0.5ex>[l]^{\pi} \ar@<0.5ex>[r]^b &
x[1]\ar@<0.5ex>[l]^s
}
$$
。由引理 \ref{dga-lemma-get-triangle}，在移位下识别 Hom 集，有
$1_x=\pi d(s)=-d(\pi)s$，其中把 $s$ 看作
$\Hom_{\mathcal{A}}^{-1}(x,C(1_x))$ 中的态射。因此，利用引理
\ref{dga-lemma-get-triangle} 的公式 (5)，有 $a=a\pi d(s)=d(s)$；利用引理
\ref{dga-lemma-get-triangle} 的公式 (6)，有 $b=-d(\pi)sb=-d(\pi)$。故
$$
1_{C(1_x)} = a\pi + sb = d(s)\pi - sd(\pi) = d(s\pi)
$$
，因为 $s$ 的次数为 $-1$。
\end{proof}

\noindent
上述引理的更一般版本将见于引理 \ref{dga-lemma-cone-homotopy}。下述引理是
引理 \ref{dga-lemma-make-commute-map} 的对应版本。

\begin{lemma}
\label{dga-lemma-homo-change}
在情形 \ref{dga-situation-ABC} 中，给定 $\text{Comp}(\mathcal{A})$ 中在同伦
意义下交换的图表
$$
\xymatrix{x\ar[r]^f\ar[d]_a & y\ar[d]^b\\
z\ar[r]^g & w}
$$
。那么
\begin{enumerate}
\item 若 $f$ 是容许单态射，则 $b$ 同伦于某个使图表交换的态射 $b'$。
\item 若 $g$ 是容许满态射，则 $a$ 同伦于某个使图表交换的态射 $a'$。
\end{enumerate}
\end{lemma}

\begin{proof}
为证明 (1)，注意假设蕴含：存在 $\Hom_{\mathcal{A}}(x,w)$ 中次数为 $-1$
的某个 $h$，使得 $bf-ga=d(h)$。由于 $f$ 是容许单态射，在范畴
$\mathcal{A}$ 中存在次数为 $0$ 的态射 $\pi:y\to x$。令
$b'=b-d(h\pi)$。于是
\begin{align*}
b'f = bf - d(h\pi)f
= &
bf - d(h\pi f) \quad (\text{因为 }d(f) = 0) \\
= &
bf-d(h) \\
= &
ga
\end{align*}
，这正是所需等式。(2) 的证明略。
\end{proof}

\noindent
下述引理是引理 \ref{dga-lemma-make-injective} 的对应版本。

\begin{lemma}
\label{dga-lemma-factor}
在情形 \ref{dga-situation-ABC} 中，设 $\alpha:x\to y$ 是
$\text{Comp}(\mathcal{A})$ 中的态射。那么在 $\text{Comp}(\mathcal{A})$
中存在分解
$$
\xymatrix{
x \ar[r]^{\tilde{\alpha}}  &
\tilde{y} \ar@<0.5ex>[r]^{\pi} &
y\ar@<0.5ex>[l]^s
}
$$
，使得
\begin{enumerate}
\item $\tilde{\alpha}$ 是容许单态射，且
$\pi\tilde{\alpha}=\alpha$；
\item $\text{Comp}(\mathcal{A})$ 中存在态射 $s:y\to\tilde{y}$，
使得 $\pi s=1_y$，且 $s\pi$ 与 $1_{\tilde{y}}$ 同伦。
\end{enumerate}
\end{lemma}

\begin{proof}
由公理 (A)，可令 $\tilde y$ 为 $y$ 与 $C(1_x)$ 的微分分次直和；即存在图表
$$
\xymatrix@C=3pc{
y \ar@<0.5ex>[r]^s  &
y\oplus C(1_x) \ar@<0.5ex>[l]^{\pi} \ar@<0.5ex>[r]^{p} &
C(1_x)\ar@<0.5ex>[l]^t
}
$$
，其中所有态射次数均为零，且都在 $\text{Comp}(\mathcal{A})$ 中。
令 $\tilde y=y\oplus C(1_x)$。于是 $1_{\tilde y}=s\pi+tp$。现考虑图表
$$
\xymatrix{
x \ar[r]^{\tilde{\alpha}}  &
\tilde{y} \ar@<0.5ex>[r]^{\pi} &
y\ar@<0.5ex>[l]^s
}
$$
，其中 $\tilde\alpha$ 由态射 $x\xrightarrow{\alpha}y$ 以及下述容许短正合列
中的自然态射 $x\to C(1_x)$ 诱导：
$$
\xymatrix{
x \ar@<0.5ex>[r]  &
C(1_x) \ar@<0.5ex>[l] \ar@<0.5ex>[r] &
x[1]\ar@<0.5ex>[l]
}
$$
因此，该图表中次数为 $0$ 的态射 $C(1_x)\to x$ 与零态射 $y\to x$
共同诱导次数为 $0$ 的态射 $\beta:\tilde y\to x$。$\tilde\alpha$ 是容许
单态射，因为它位于容许短正合列
$$
\xymatrix{
x\ar[r]^{\tilde{\alpha}} &
\tilde{y} \ar[r] &
x[1]
}
$$
中。此外，由 $\tilde\alpha$ 的构造，$\pi\tilde\alpha=\alpha$；由第一个
图表，$\pi s=1_y$。只须证明 $s\pi$ 与 $1_{\tilde y}$ 同伦。把 $1_x$
写成某个次数为 $-1$ 的映射的微分 $d(h)$。于是，最后一个断言由下式得出：
\begin{align*}
1_{\tilde{y}} - s\pi
= &
tp \\
= &
t(dh)p\quad\text{（由引理 \ref{dga-lemma-id-cone-null}）} \\
= &
d(thp)
\end{align*}
，因为 $dt=dp=0$，且 $t$ 的次数为零。
\end{proof}

\noindent
下述引理是引理 \ref{dga-lemma-sequence-maps-split} 的对应版本。

\begin{lemma}
\label{dga-lemma-analogue-sequence-maps-split}
在情形 \ref{dga-situation-ABC} 中，设
$x_1\to x_2\to\ldots\to x_n$ 是 $\text{Comp}(\mathcal{A})$ 中一列可复合
态射。那么 $\text{Comp}(\mathcal{A})$ 中存在交换图表
$$
\xymatrix{x_1\ar[r] & x_2\ar[r] & \ldots\ar[r] & x_n\\
y_1\ar[r]\ar[u] & y_2\ar[r]\ar[u] & \ldots\ar[r] & y_n\ar[u]}
$$
，使每个 $y_i\to y_{i+1}$ 都是容许单态射，而每个 $y_i\to x_i$ 都是同伦等价。
\end{lemma}

\begin{proof}
$n=1$ 的情形是平凡的：只须取 $y_1=x_1$，而 $x_1$ 上的恒等态射当然是
同伦等价。$n=2$ 的情形由引理 \ref{dga-lemma-factor} 给出。假设已经把图表构造
到 $x_{n-1}$。对复合 $y_{n-1}\to x_{n-1}\to x_n$ 应用引理
\ref{dga-lemma-factor}，得到 $y_n$。于是 $y_{n-1}\to y_n$ 是容许单态射，
$y_n\to x_n$ 是同伦等价。
\end{proof}

\noindent
下述引理是引理 \ref{dga-lemma-nilpotent} 的对应版本。

\begin{lemma}
\label{dga-lemma-triseq}
在情形 \ref{dga-situation-ABC} 中，设 $x_i\to y_i\to z_i$ 是
$\mathcal{A}$ 中的态射（$i=1,2,3$），并且 $x_2\to y_2\to z_2$ 是容许
短正合列。设 $b:y_1\to y_2$ 和 $b':y_2\to y_3$ 是
$\text{Comp}(\mathcal{A})$ 中的态射，使得
$$
\vcenter{
\xymatrix{
x_1 \ar[d]_0 \ar[r] &
y_1 \ar[r] \ar[d]_b &
z_1 \ar[d]_0 \\
x_2 \ar[r] & y_2 \ar[r] & z_2
}
}
\quad\text{且}\quad
\vcenter{
\xymatrix{
x_2 \ar[d]^0 \ar[r] &
y_2 \ar[r] \ar[d]^{b'} &
z_2 \ar[d]^0 \\
x_3 \ar[r] & y_3 \ar[r] & z_3
}
}
$$
中的两个图表均在同伦意义下交换。那么 $b'\circ b$ 与 $0$ 同伦。
\end{lemma}

\begin{proof}
由引理 \ref{dga-lemma-homo-change}，可分别用与 $b,b'$ 同伦的映射
$\tilde b,\tilde b'$ 替代它们，使左图的右方块交换，右图的左方块交换。
设对 $\mathcal{A}$ 中次数为 $-1$ 的态射 $h,h'$，有
$b=\tilde b+d(h)$、$b'=\tilde b'+d(h')$。于是
$$
b'b = \tilde{b}'\tilde{b} + d(\tilde{b}'h + h'\tilde{b} + h'd(h))
$$
，因为 $d(\tilde b)=d(\tilde b')=0$。也就是说，$b'b$ 与
$\tilde b'\tilde b$ 同伦。现证明 $\tilde b'\tilde b=0$。由于
$x_2\xrightarrow{f}y_2\xrightarrow{g}z_2$ 是容许短正合列，存在次数为 $0$
的态射 $\pi:y_2\to x_2$ 和 $s:z_2\to y_2$，使得
$\text{id}_{y_2}=f\pi+sg$。因此
$$
\tilde{b}'\tilde{b} = \tilde{b}'(f\pi + sg)\tilde{b} = 0
$$
，因为两个交换方块分别蕴含 $g\tilde b=0$ 和 $\tilde b'f=0$。
\end{proof}

\noindent
下述引理是引理 \ref{dga-lemma-triangle-independent-splittings} 的对应版本。

\begin{lemma}
\label{dga-lemma-analogue-triangle-independent-splittings}
在情形 \ref{dga-situation-ABC} 中，设
$0\to x\to y\to z\to 0$ 是 $\text{Comp}(\mathcal{A})$ 中的容许短正合列。
以引理 \ref{dga-lemma-get-triangle} 中定义的 $\delta:z\to x[1]$ 构成三角
$$
\xymatrix{x\ar[r] & y\ar[r] & z\ar[r]^{\delta} & x[1]}
$$
。在 $K(\mathcal{A})$ 中，除典范同构外，它不依赖于引理
\ref{dga-lemma-get-triangle} 中所作的选择。
\end{lemma}

\begin{proof}
假设 $\delta$ 由分裂
$$
\xymatrix{
x \ar@<0.5ex>[r]^{a} &
y \ar@<0.5ex>[r]^b\ar@<0.5ex>[l]^{\pi} &
z \ar@<0.5ex>[l]^s
}
$$
定义，而 $\delta'$ 由以 $\pi',s'$ 代替 $\pi,s$ 的分裂定义。那么
$$
s'-s = (a\pi + sb)(s'-s) = a\pi s'
$$
，因为 $bs'=bs=1_z$ 且 $\pi s=0$。类似地，
$$
\pi' - \pi = (\pi' - \pi)(a\pi + sb) = \pi'sb
$$
由于按引理 \ref{dga-lemma-get-triangle} 的构造，$\delta=\pi d(s)$ 且
$\delta'=\pi'd(s')$，可计算
$$
\delta' = \pi'd(s') = (\pi + \pi'sb)d(s + a\pi s') = \delta + d(\pi s')
$$
；这里使用了 $\pi a=1_x$、$ba=0$，以及由引理
\ref{dga-lemma-get-triangle} 的公式 (5) 得出的
$\pi'sbd(s')=\pi'sba\pi d(s')=0$。
\end{proof}

\noindent
下述引理是引理 \ref{dga-lemma-rotate-cone} 的对应版本。

\begin{lemma}
\label{dga-lemma-restate-axiom-c}
在情形 \ref{dga-situation-ABC} 中，设 $f:x\to y$ 是
$\text{Comp}(\mathcal{A})$ 中的态射。三角 $(y,c(f),x[1],i,p,f[1])$
是与容许短正合列
$$
\xymatrix{y\ar[r] & c(f) \ar[r] & x[1]}
$$
相联系的三角，其中锥 $c(f)$ 如引理 \ref{dga-lemma-get-triangle} 中所定义。
\end{lemma}

\begin{proof}
这由公理 (C) 得出。
\end{proof}

\noindent
下述引理是引理 \ref{dga-lemma-rotate-triangle} 的对应版本。

\begin{lemma}
\label{dga-lemma-cone-rotate-isom}
在情形 \ref{dga-situation-ABC} 中，设 $\alpha:x\to y$ 和 $\beta:y\to z$
给出 $\text{Comp}(\mathcal{A})$ 中的容许短正合列
$$
\xymatrix{
x \ar[r] &
y\ar[r] &
z
}
$$
。令 $(x,y,z,\alpha,\beta,\delta)$ 为 $K(\mathcal{A})$ 中相联系的三角。
那么三角
$$
(z[-1], x, y, \delta[-1], \alpha, \beta)
\quad\text{且}\quad
(z[-1], x, c(\delta[-1]), \delta[-1], i, p)
$$
同构。
\end{lemma}

\begin{proof}
有如下形式的图表
$$
\xymatrix{
z[-1]\ar[r]^{\delta[-1]}\ar[d]^1 &
x\ar@<0.5ex>[r]^{\alpha}\ar[d]^1 &
y\ar@<0.5ex>[r]^{\beta}\ar@{.>}[d]\ar@<0.5ex>[l]^{\tilde{\alpha}} &
z\ar[d]^1\ar@<0.5ex>[l]^{\tilde\beta} \\
z[-1] \ar[r]^{\delta[-1]} &
x\ar@<0.5ex>[r]^i &
c(\delta[-1]) \ar@<0.5ex>[r]^p\ar@<0.5ex>[l]^{\tilde i} &
z\ar@<0.5ex>[l]^{\tilde p}
}
$$
，其中 $\alpha,\beta,i,p$ 的分裂分别由
$\tilde\alpha,\tilde\beta,\tilde i,\tilde p$ 给出。把态射
$y\to c(\delta[-1])$ 定义为 $i\tilde\alpha+\tilde p\beta$，把态射
$c(\delta[-1])\to y$ 定义为 $\alpha\tilde i+\tilde\beta p$。先验证它们
确实是 $\text{Comp}(\mathcal{A})$ 中的态射。由引理
\ref{dga-lemma-get-triangle} 中的恒等式，有关系
$\delta[-1]=\tilde\alpha d(\tilde\beta)=-d(\tilde\alpha)\tilde\beta$
以及 $\delta[-1]=\tilde i d(\tilde p)$。于是
\begin{align*}
d(\tilde{\alpha})
& =
d(\tilde{\alpha})\tilde{\beta}\beta \\
& =
-\delta[-1]\beta
\end{align*}
，其中第一个等号使用了引理 \ref{dga-lemma-get-triangle} 的公式 (6)，第二个
等号使用了前述关系。类似地，得到 $d(\tilde p)=i\delta[-1]$。因此
\begin{align*}
d(i\tilde{\alpha} + \tilde{p}\beta)
& =
d(i)\tilde{\alpha} + id(\tilde{\alpha}) +
d(\tilde{p})\beta + \tilde{p}d(\beta) \\
& =
id(\tilde{\alpha}) + d(\tilde{p})\beta \\
& =
-i\delta[-1]\beta + i\delta[-1]\beta \\
& =
0
\end{align*}
，故 $i\tilde\alpha+\tilde p\beta$ 确实是 $\text{Comp}(\mathcal{A})$
中的态射。类似计算表明，$\alpha\tilde i+\tilde\beta p$ 也是
$\text{Comp}(\mathcal{A})$ 中的态射。显然这些态射组成上述交换图表。计算得
\begin{align*}
(i\tilde{\alpha} + \tilde{p}\beta)(\alpha \tilde{i} + \tilde{\beta} p)
& =
i\tilde{\alpha}\alpha\tilde{i} + i\tilde{\alpha}\tilde{\beta}p
+ \tilde{p}\beta\alpha\tilde{i} + \tilde{p}\beta\tilde{\beta}p \\
& =
i\tilde{i} + \tilde{p}p \\
& =
1_{c(\delta[-1])}
\end{align*}
，其中自由使用了引理 \ref{dga-lemma-get-triangle} 中的恒等式。类似地，计算得
$(\alpha \tilde{i} + \tilde{\beta} p)(i\tilde{\alpha} + \tilde{p}\beta) = 1_y$,
，故 $y\cong c(\delta[-1])$。因此，所讨论的两个三角同构。
\end{proof}

\noindent
下述引理是引理 \ref{dga-lemma-third-isomorphism} 的对应版本。

\begin{lemma}
\label{dga-lemma-analogue-third-isomorphism}
在情形 \ref{dga-situation-ABC} 中，设 $f_1:x_1\to y_1$ 和
$f_2:x_2\to y_2$ 是 $\text{Comp}(\mathcal{A})$ 中的态射。设
$$
(a,b,c): (x_1,y_1,c(f_1), f_1, i_1, p_1) \to (x_2,y_2, c(f_2), f_2, i_1, p_1)
$$
是 $K(\mathcal{A})$ 中任意三角态射。若 $a$ 和 $b$ 是同伦等价，则 $c$
也是同伦等价。
\end{lemma}

\begin{proof}
由于 $a,b$ 是同伦等价，它们在 $K(\mathcal{A})$ 中可逆。以
$a^{-1},b^{-1}$ 表示它们在 $K(\mathcal{A})$ 中的逆，得到交换图表
$$
\xymatrix{
x_2\ar[d]^{a^{-1}}\ar[r]^{f_2} &
y_2\ar[d]^{b^{-1}}\ar[r]^{i_2} &
c(f_2)\ar[d]^{c'} \\
x_1\ar[r]^{f_1} &
y_1 \ar[r]^{i_1} &
c(f_1)
}
$$
，其中映射 $c'$ 通过把引理 \ref{dga-lemma-cone} 应用于上图左侧交换方块而定义。
由于图表在 $K(\mathcal{A})$ 中交换，由引理 \ref{dga-lemma-triseq}，只需证明：
给定 $K(\mathcal{A})$ 中的三角态射
$(1,1,c): (x,y,c(f),f,i,p)\to (x,y,c(f),f,i,p)$
，映射 $c$ 是 $K(\mathcal{A})$ 中的同构。在 $K(\mathcal{A})$ 中有交换图表
$$
\vcenter{
\xymatrix{
y\ar[d]^{1}\ar[r] &
c(f)\ar[d]^{c}\ar[r] &
x[1]\ar[d]^{1} \\
y\ar[r] &
c(f) \ar[r] &
x[1]
}
}
\quad\Rightarrow\quad
\vcenter{
\xymatrix{
y\ar[d]^{0}\ar[r] &
c(f)\ar[d]^{c-1}\ar[r] &
x[1]\ar[d]^{0} \\
y\ar[r] &
c(f) \ar[r] &
x[1]
}
}
$$
由于各行都是容许短正合列，由引理 \ref{dga-lemma-triseq} 得到恒等式
$(c-1)^2=0$。由此推出 $2-c$ 是 $c$ 在 $K(\mathcal{A})$ 中的逆，
因而 $c$ 是同构。
\end{proof}

\noindent
下述引理是引理 \ref{dga-lemma-the-same-up-to-isomorphisms} 的对应版本。

\begin{lemma}
\label{dga-lemma-cone-homotopy}
在情形 \ref{dga-situation-ABC} 中：
\begin{enumerate}
\item 给定容许短正合列 $x\xrightarrow{\alpha}y\xrightarrow{\beta}z$。
那么存在同伦等价 $e:C(\alpha)\to z$，使图表
\begin{equation}
\label{dga-equation-cone-isom-triangle}
\vcenter{
\xymatrix{
x\ar[r]^{\alpha}\ar[d] &
y\ar[r]^{b}\ar[d] &
C(\alpha)\ar[r]^{-c}\ar@{.>}[d]^{e} &
x[1]\ar[d] \\
x\ar[r]^{\alpha} &
y\ar[r]^{\beta} &
z\ar[r]^{\delta} & x[1]
}
}
\end{equation}
在 $K(\mathcal{A})$ 中定义三角同构。这里
$y\xrightarrow{b}C(\alpha)\xrightarrow{c}x[1]$ 是按公理 (C) 给出的容许
短正合列。
\item 给定 $\text{Comp}(\mathcal{A})$ 中的态射 $\alpha:x\to y$，令
$x\xrightarrow{\tilde\alpha}\tilde y\to y$ 为引理 \ref{dga-lemma-factor}
中给出的分解，其中容许单态射 $x\xrightarrow{\tilde\alpha}y$ 延拓为容许短正合列
$$
\xymatrix{
x \ar[r]^{\tilde{\alpha}} &
\tilde{y} \ar[r] & z
}
$$
。那么存在三角同构
$$
\xymatrix{
x \ar[r]^{\tilde{\alpha}} \ar[d] &
\tilde{y} \ar[r] \ar[d] &
z \ar[r]^{\delta} \ar@{.>}[d]^{e} &
x[1] \ar[d] \\
x \ar[r]^{\alpha} &
y \ar[r] &
C(\alpha) \ar[r]^{-c} &
x[1]
}
$$
，其中上方三角是与列 $x\xrightarrow{\tilde\alpha}\tilde y\to z$ 相联系的三角。
\end{enumerate}
\end{lemma}

\begin{proof}
对 (1)，考虑下述更完整的图表，其中\emph{不}改变 $c$ 的符号：
$$
\xymatrix{
x\ar@<0.5ex>[r]^{\alpha} \ar[d] &
y\ar@<0.5ex>[l]^{\pi} \ar@<0.5ex>[r]^{b}\ar[d] &
C(\alpha)\ar@<0.5ex>[l]^{p} \ar@<0.5ex>[r]^{c}\ar@{.>}@<0.5ex>[d]^{e} &
x[1]\ar@<0.5ex>[l]^{\sigma} \ar[d]\ar@<0.5ex>[r]^{\alpha} &
y[1]\ar@<0.5ex>[l]^{\pi} \\
x\ar@<0.5ex>[r]^{\alpha} &
y\ar@<0.5ex>[r]^{\beta} \ar@<0.5ex>[l]^{\pi} &
z\ar[r]^{\delta}\ar@<0.5ex>[l]^{s} \ar@{.>}@<0.5ex>[u]^{f} &
x[1]
}
$$
其中容许短正合列 $x\xrightarrow{\alpha}y\xrightarrow{\beta}z$ 带有分裂
$\pi,s$，容许短正合列
$y\xrightarrow{b}C(\alpha)\xrightarrow{c}x[1]$ 带有分裂 $p,\sigma$。
注意（在移位下识别 Hom 集），由引理 \ref{dga-lemma-get-triangle} 中的构造，
$$
\alpha = pd(\sigma) = -d(p)\sigma,\quad
\delta = \pi d(s) = -d(\pi)s
$$
。

\medskip\noindent
定义 $e=\beta p$ 和 $f=bs-\sigma\delta$。先验证它们是
$\text{Comp}(\mathcal{A})$ 中的态射。为证明 $d(e)=\beta d(p)$ 为零，
只需证明 $\beta d(p)b$ 和 $\beta d(p)\sigma$ 均为零，而
$$
\beta d(p)b = \beta d(pb) = \beta d(1_y) = 0,\quad
\beta d(p)\sigma = -\beta\alpha = 0
$$
。类似地，为验证 $d(f)=bd(s)-d(\sigma)\delta$ 为零，只需验证它与 $p$
及 $c$ 的后复合都为零，而
\begin{align*}
pbd(s) - pd(\sigma)\delta
= &
d(s)-\alpha\delta = d(s)-\alpha\pi d(s) = 0 \\
cbd(s)-cd(\sigma)\delta
= &
-cd(\sigma)\delta = -d(c\sigma)\delta = 0
\end{align*}
。图表 \ref{dga-equation-cone-isom-triangle} 左侧两个方块的交换性直接由定义得出。
在证明右侧方块在同伦意义下交换之前，先验证 $e$ 是同伦等价。显然，
$$
ef=\beta p (bs-\sigma\delta)=\beta s=1_z
$$
。为验证 $fe$ 与 $1_{C(\alpha)}$ 同伦，先注意
$$
b\alpha = bpd(\alpha) = d(\sigma),\quad
\alpha c = -d(p)\sigma c = -d(p),\quad
d(\pi)p = d(\pi)s\beta p = -\delta\beta p
$$
。利用这些恒等式，计算得
\begin{align*}
1_{C(\alpha)} = &
bp + \sigma c
\quad (\text{由 }y \xrightarrow{b} C(\alpha) \xrightarrow{c} x[1]) \\
= &
b(\alpha\pi + s\beta)p + \sigma(\pi\alpha)c
\quad (\text{由 }x \xrightarrow{\alpha} y \xrightarrow{\beta} z) \\
= &
d(\sigma)\pi p + bs\beta p - \sigma\pi d(p)
\quad (\text{由上述前两个恒等式}) \\
= &
d(\sigma)\pi p + bs\beta p - \sigma\delta\beta p
+ \sigma\delta\beta p - \sigma\pi d(p) \\
= &
(bs - \sigma\delta)\beta p + d(\sigma)\pi p
- \sigma d(\pi)p - \sigma\pi d(p)\quad
(\text{由上述第三个恒等式}) \\
= &
fe + d(\sigma \pi p)
\end{align*}
，因为 $\sigma\in\Hom^{-1}(x,C(\alpha))$（参见引理
\ref{dga-lemma-id-cone-null} 的证明）。因此 $e,f$ 互为同伦逆。最后，为验证
图表 \ref{dga-equation-cone-isom-triangle} 的右侧方块在同伦意义下交换，
只须验证 $-cf=\delta$。这由下式得出：
$$
-cf = -c(bs-\sigma\delta) = c\sigma\delta = \delta
$$
，因为 $cb=0$。

\medskip\noindent
对 (2)，考虑引理 \ref{dga-lemma-factor} 中给出的分解
$x\xrightarrow{\tilde\alpha}\tilde y\to y$，故第二个态射是同伦等价。
由引理 \ref{dga-lemma-cone} 和 \ref{dga-lemma-analogue-third-isomorphism}，
下列两个三角之间存在三角同构：
$$
x \xrightarrow{\alpha} y \to C(\alpha) \to x[1]
\quad\text{且}\quad
x \xrightarrow{\tilde{\alpha}} \tilde{y} \to C(\tilde{\alpha}) \to x[1]
$$
。由于可以复合三角同构，通过以 $\tilde\alpha$ 替代 $\alpha$、以
$\tilde y$ 替代 $y$、以 $C(\tilde\alpha)$ 替代 $C(\alpha)$，可假设
$\alpha$ 是容许单态射。在这种情形下，结论由 (1) 得出。
\end{proof}

\noindent
下述引理是引理 \ref{dga-lemma-homotopy-category-pre-triangulated} 的对应版本。

\begin{lemma}
\label{dga-lemma-analogue-homotopy-category-pre-triangulated}
在情形 \ref{dga-situation-ABC} 中，同伦范畴 $K(\mathcal{A})$ 连同其自然平移
函子和有区别三角，是预三角范畴。
\end{lemma}

\begin{proof}
逐一验证 TR1、TR2 和 TR3。

\medskip\noindent
TR1 的证明。由定义，每个与有区别三角同构的三角都是有区别的。由于
$$
\xymatrix{x\ar[r]^{1_x} & x\ar[r] & 0}
$$
是容许短正合列，$(x,x,0,1_x,0,0)$ 是有区别三角。此外，给定
$\text{Comp}(\mathcal{A})$ 中的态射 $\alpha:x\to y$，由引理
\ref{dga-lemma-cone-homotopy}，三角 $(x,y,c(\alpha),\alpha,i,-p)$ 是有区别的。

\medskip\noindent
TR2 的证明。设 $(x,y,z,\alpha,\beta,\gamma)$ 是三角，并假设
$(y,z,x[1],\beta,\gamma,-\alpha[1])$ 是有区别的。则存在容许短正合列
$0\to x'\to y'\to z'\to0$，使相联系的三角
$(x',y',z',\alpha',\beta',\gamma')$ 同构于
$(y,z,x[1],\beta,\gamma,-\alpha[1])$。旋转后，得出
$(x,y,z,\alpha,\beta,\gamma)$ 同构于
$(z'[-1],x',y',\gamma'[-1],\alpha',\beta')$。由引理
\ref{dga-lemma-cone-rotate-isom}，后者同构于
$(z'[-1],x',c(\gamma'[-1]),\gamma'[-1],i,p)$。按下图所示改变符号并复合
这两个同构：
$$
\xymatrix@C=3pc{
x\ar[r]^{\alpha}\ar[d] &
y\ar[r]^{\beta}\ar[d] &
z\ar[r]^{\gamma}\ar[d] &
x[1]\ar[d] \\
z'[-1]\ar[r]^{-\gamma'[-1]}\ar[d]_{-1_{z'[-1]}} &
x \ar[r]^{\alpha'}\ar@{=}[d] &
y' \ar[r]^{\beta'} \ar[d] &
z'\ar[d]^{-1_{z'}} \\
z'[-1]\ar[r]^{\gamma'[-1]} &
x \ar[r]^{\alpha'} &
c(\gamma'[-1]) \ar[r]^{-p} &
z'
}
$$
由引理 \ref{dga-lemma-cone-homotopy} (2)，推出
$(x,y,z,\alpha,\beta,\gamma)$ 是有区别的。反过来，假设
$(x,y,z,\alpha,\beta,\gamma)$ 是有区别的。由引理
\ref{dga-lemma-cone-homotopy} (1)，它同构于形如
$(x',y',c(\alpha'),\alpha',i,-p)$ 的三角，其中
$\alpha':x'\to y'$ 是 $\text{Comp}(\mathcal{A})$ 中的某个态射。
旋转后的三角 $(y,z,x[1],\beta,\gamma,-\alpha[1])$ 同构于
$(y',c(\alpha'),x'[1],i,-p,-\alpha[1])$，而后者同构于
$(y',c(\alpha'),x'[1],i,p,\alpha[1])$。由引理
\ref{dga-lemma-restate-axiom-c}，这个三角是有区别的，从而
$(y,z,x[1],\beta,\gamma,-\alpha[1])$ 是有区别的。

\medskip\noindent
TR3 的证明。设 $(x,y,z,\alpha,\beta,\gamma)$ 和
$(x',y',z',\alpha',\beta',\gamma')$ 是 $\text{Comp}(\mathcal{A})$
中的有区别三角，且 $f:x\to x'$、$g:y\to y'$ 是满足
$\alpha'\circ f=g\circ\alpha$ 的态射。由引理
\ref{dga-lemma-cone-homotopy}，可假设
$(x,y,z,\alpha,\beta,\gamma)= (x,y,c(\alpha),\alpha, i, -p)$
，且
$(x',y',z',\alpha',\beta',\gamma')=(x',y',c(\alpha'),\alpha',i',-p')$。
现应用引理 \ref{dga-lemma-cone} 即得结论。
\end{proof}

\noindent
下述引理是引理 \ref{dga-lemma-two-split-injections} 的对应版本。

\begin{lemma}
\label{dga-lemma-dgc-analogue-tr4}
在情形 \ref{dga-situation-ABC} 中，给定 $\mathcal{A}$ 中的容许单态射
$x\xrightarrow{\alpha}y$、$y\xrightarrow{\beta}z$，存在有区别三角
$(x,y,q_1,\alpha,p_1,\delta_1)$、$(x,z,q_2,\beta\alpha,p_2,\delta_2)$
和 $(y,z,q_3,\beta,p_3,\delta_3)$，使 TR4 成立。
\end{lemma}

\begin{proof}
给定容许单态射 $x\xrightarrow{\alpha}y$ 和 $y\xrightarrow{\beta}z$，
通过将它们延拓为容许短正合列，可以得到有区别三角
$$
\xymatrix{
x\ar@<0.5ex>[r]^{\alpha} &
y\ar@<0.5ex>[l]^{\pi_1} \ar@<0.5ex>[r]^{p_1} &
q_1 \ar[r]^{\delta_1} \ar@<0.5ex>[l]^{s_1} &
x[1]
}
$$
$$
\xymatrix{
x\ar@<0.5ex>[r]^{\beta\alpha} &
z\ar@<0.5ex>[l]^{\pi_1\pi_3} \ar@<0.5ex>[r]^{p_2} &
q_2 \ar[r]^{\delta_2} \ar@<0.5ex>[l]^{s_2} &
x[1]
}
$$
$$
\xymatrix{
y\ar@<0.5ex>[r]^{\beta} &
z\ar@<0.5ex>[l]^{\pi_3} \ar@<0.5ex>[r]^{p_3} &
q_3 \ar[r]^{\delta_3} \ar@<0.5ex>[l]^{s_3} &
x[1]
}
$$
。在这些图表中，映射 $\delta_i$ 类似于引理 \ref{dga-lemma-get-triangle}
中定义的映射，定义为 $\delta_i=\pi_i d(s_i)$。它们构成下述实线交换图表
$$
\xymatrix@C=5pc@R=3pc{
x\ar@<0.5ex>[r]^{\alpha} \ar@<0.5ex>[dr]^{\beta\alpha} &
y\ar@<0.5ex>[d]^{\beta} \ar@<0.5ex>[l]^{\pi_1} \ar@<0.5ex>[r]^{p_1} &
q_1 \ar[r]^{\delta_1} \ar@<0.5ex>[l]^{s_1} \ar@{.>}[dd]^{p_2\beta s_1} &
x[1] \\
 &
z \ar@<0.5ex>[u]^{\pi_3}\ar@<0.5ex>[d]^{p_3}
\ar@<0.5ex>[dr]^{p_2} \ar@<0.5ex>[ul]^{\pi_1\pi_3} & & \\
 &
q_3\ar@<0.5ex>[u]^{s_3} \ar[d]^{\delta_3} &
q_2 \ar@{.>}[l]^{p_3s_2} \ar@<0.5ex>[ul]^{s_2} \ar[dr]^{\delta_2} \\
 &
y[1] & & x[1]}
$$
，其中虚线箭头按图中所示定义。显然，它们的复合
$p_3s_2p_2\beta s_1=0$，因为 $s_2p_2=0$。断言两者都是
$\text{Comp}(\mathcal{A})$ 中的态射。可利用引理
\ref{dga-lemma-get-triangle} 中的等式验证这一点：
$$
d(p_2\beta s_1) = p_2\beta d(s_1) = p_2\beta\alpha\pi_1 d(s_1) = 0
$$
，因为 $p_2\beta\alpha=0$；并且
$$
d(p_3s_2) = p_3d(s_2) = p_3\beta\alpha\pi_1\pi_3 d(s_2) = 0
$$
，因为 $p_3\beta=0$。为验证 $q_1\to q_2\to q_3$ 是容许短正合列，
还须证明在底层分次范畴中 $q_2=q_1\oplus q_3$，且上述两个态射分别是
余投影和投影。为此，注意在底层分次范畴 $\mathcal{C}$ 中有
$$
y = x\oplus q_1,\quad
z = y\oplus q_3 = x\oplus q_1\oplus q_3
$$
，其中 $\pi_1\pi_3$ 给出到第一个因子的投影态射
$x\oplus q_1\oplus q_3\to z$。由 $\mathcal{A}$ 的公理 (A)，
$\mathcal{C}$ 是加性范畴。因此可应用《同调代数》引理
\ref{homology-lemma-additive-cat-biproduct-kernel}，得出
$$
\Ker(\pi_1\pi_3) = q_1\oplus q_3
$$
在 $\mathcal{C}$ 中成立。再把《同调代数》引理
\ref{homology-lemma-additive-cat-biproduct-kernel} 应用于
$z=x\oplus q_2$，得到 $\Ker(\pi_1\pi_3)=q_2$。故在 $\mathcal{C}$ 中
$q_2\cong q_1\oplus q_3$。显然，上述虚线态射给出余投影和投影。

\medskip\noindent
最后，还须验证容许短正合列 $q_1\to q_2\to q_3$ 诱导的态射
$\delta:q_3\to q_1[1]$ 与 $p_1\delta_3$ 一致。由引理
\ref{dga-lemma-get-triangle} 中的构造，态射 $\delta$ 由下式给出：
\begin{align*}
p_1\pi_3s_2d(p_2s_3)
= &
p_1\pi_3s_2p_2d(s_3) \\
= &
p_1\pi_3(1-\beta\alpha\pi_1\pi_3)d(s_3) \\
= &
p_1\pi_3d(s_3)\quad (\text{因为 }\pi_3\beta = 0) \\
= &
p_1\delta_3
\end{align*}
，这正是所需等式。证明完成。
\end{proof}

\noindent
汇集以上结果，终于得到命题
\ref{dga-proposition-homotopy-category-triangulated} 的对应版本。

\begin{proposition}
\label{dga-proposition-ABC-homotopy-category-triangulated}
在情形 \ref{dga-situation-ABC} 中，同伦范畴 $K(\mathcal{A})$ 连同其自然平移
函子和有区别三角，是三角范畴。
\end{proposition}

\begin{proof}
由引理 \ref{dga-lemma-analogue-homotopy-category-pre-triangulated}，
$K(\mathcal{A})$ 是预三角范畴。把引理
\ref{dga-lemma-analogue-sequence-maps-split}、\ref{dga-lemma-dgc-analogue-tr4}
与《导出范畴》引理 \ref{derived-lemma-easier-axiom-four} 结合，得出
$K(\mathcal{A})$ 是三角范畴。
\end{proof}

\begin{lemma}
\label{dga-lemma-functor-between-ABC}
设 $R$ 是环，$F:\mathcal{A}\to\mathcal{B}$ 是满足公理 (A)、(B)、(C)
的 $R$ 上微分分次范畴之间的函子，并且 $F(x[1])=F(x)[1]$。
那么 $F$ 诱导三角范畴之间的正合函子
$K(\mathcal{A})\to K(\mathcal{B})$。
\end{lemma}

\begin{proof}
确实，若 $x\to y\to z$ 是 $\text{Comp}(\mathcal{A})$ 中的容许短正合列，
则 $F(x)\to F(y)\to F(z)$ 是 $\text{Comp}(\mathcal{B})$ 中的容许短正合列。
此外，引理 \ref{dga-lemma-get-triangle} 中构造的“边界”态射
$\delta=\pi\text{d}(s):z\to x[1]$ 产生态射
$F(\delta):F(z)\to F(x[1])=F(x)[1]$；它等于容许短正合列
$F(x)\to F(y)\to F(z)$ 的边界映射 $F(\pi)\text{d}(F(s))$。
\end{proof}





\section{双模}
\label{dga-section-bimodules}

\noindent
继续第 \ref{dga-section-tensor-product} 节开始的讨论。

\begin{definition}
\label{dga-definition-bimodule}
双模。设 $R$ 是环。
\begin{enumerate}
\item 设 $A,B$ 是 $R$-代数。{\it $(A,B)$-双模}是带有 $R$-双线性映射
$$
A \times M \to M, (a, x) \mapsto ax
\quad\text{且}\quad
M \times B \to M, (x, b) \mapsto xb
$$
的 $R$-模 $M$，并满足
\begin{enumerate}
\item $a'(ax)=(a'a)x$ 且 $(xb)b'=x(bb')$；
\item $a(xb)=(ax)b$；以及
\item $1x=x=x1$。
\end{enumerate}
\item 设 $A,B$ 是 $\mathbf{Z}$-分次 $R$-代数。{\it 分次 $(A,B)$-双模}
是带有分次 $M=\bigoplus M^n$ 的 $(A,B)$-双模 $M$，使得
$A^nM^m\subset M^{n+m}$ 且 $M^nB^m\subset M^{n+m}$。
\item 设 $A,B$ 是微分分次 $R$-代数。{\it 微分分次 $(A,B)$-双模}
是带有次数为 $1$ 的齐次微分 $\text{d}:M\to M$ 的分次 $(A,B)$-双模，
使得对齐次元素 $a\in A$、$x\in M$、$b\in B$，有
$\text{d}(ax)=\text{d}(a)x+(-1)^{\deg(a)}a\text{d}(x)$ 且
$\text{d}(xb)=\text{d}(x)b+(-1)^{\deg(x)}x\text{d}(b)$。
\end{enumerate}
\end{definition}

\noindent
注意，微分分次 $(A,B)$-双模 $M$ 与下述数据是一回事：$M$ 是右微分分次
$B$-模，同时也是左微分分次 $A$-模，两个分次和微分彼此一致，并且
$A$-模结构与 $B$-模结构交换。精确表述如下。

\begin{lemma}
\label{dga-lemma-what-makes-a-bimodule-dg}
设 $R$ 是环，$(A,\text{d})$ 和 $(B,\text{d})$ 是 $R$ 上的微分分次代数，
$M$ 是右微分分次 $B$-模。$M$ 上与给定微分分次 $B$-模结构相容的
$(A,B)$-双模结构，与微分分次 $R$-代数同态
$$
A
\longrightarrow
\Hom_{\text{Mod}^{dg}_{(B, \text{d})}}(M, M)
$$
之间有一一对应。
\end{lemma}

\begin{proof}
设 $\mu:A\times M\to M$ 在 $M$ 的底层 $R$-模复形 $M^\bullet$ 上定义
左微分分次 $A$-模结构。由引理 \ref{dga-lemma-left-module-structure}，结构
$\mu$ 对应于微分分次 $R$-代数的映射
$\gamma:A\to\Hom^\bullet(M^\bullet,M^\bullet)$。本引理的断言恰是：
$\mu$ 与 $B$-作用交换，当且仅当 $\gamma$ 的像落在
$$
\Hom_{\text{Mod}^{dg}_{(B, \text{d})}}(M, M) \subset
\Hom^\bullet(M^\bullet, M^\bullet)
$$
中。详细计算略。
\end{proof}

\noindent
设 $M$ 是微分分次 $(A,B)$-双模。回忆第 \ref{dga-section-left-modules} 节：
左微分分次 $A$-模结构对应于右微分分次 $A^{opp}$-模结构。由于 $A$-模
结构与 $B$-模结构交换，这赋予 $M$ 微分分次
$A^{opp}\otimes_R B$-模结构：
$$
x \cdot (a \otimes b) = (-1)^{\deg(a)\deg(x)} axb
$$
。反过来，若有微分分次 $A^{opp}\otimes_R B$-模 $M$，则可利用上式得到
微分分次 $(A,B)$-双模。

\begin{lemma}
\label{dga-lemma-bimodule-over-tensor}
设 $R$ 是环，$(A,\text{d})$ 和 $(B,\text{d})$ 是 $R$ 上的微分分次代数。
上述构造定义范畴等价
$$
\begin{matrix}
\text{微分分次}\\
(A, B)\text{-双模}
\end{matrix}
\longleftrightarrow
\begin{matrix}
\text{右微分分次}\\
A^{opp} \otimes_R B\text{-模}
\end{matrix}
$$
\end{lemma}

\begin{proof}
由上述讨论立即得出。
\end{proof}

\noindent
设 $R$ 是环，$(A,\text{d})$ 和 $(B,\text{d})$ 是微分分次 $R$-代数。
设 $P$ 是微分分次 $(A,B)$-双模。若存在滤过
$$
0 = F_{-1}P \subset F_0P \subset F_1P \subset \ldots \subset P
$$
，它由微分分次 $(A,B)$-双模组成，并且
\begin{enumerate}
\item $P=\bigcup F_pP$；
\item 包含映射 $F_iP\to F_{i+1}P$ 作为分次 $(A,B)$-双模映射是分裂的；
\item 作为微分分次 $(A,B)$-双模，商 $F_{i+1}P/F_iP$ 同构于若干个
$(A\otimes_R B)[k]$ 的直和，
\end{enumerate}
则称 $P$ {\it 具有性质 (P)}。

\begin{lemma}
\label{dga-lemma-bimodule-resolve}
设 $R$ 是环，$(A,\text{d})$ 和 $(B,\text{d})$ 是微分分次 $R$-代数，
$M$ 是微分分次 $(A,B)$-双模。存在微分分次 $(A,B)$-双模同态 $P\to M$，
它是拟同构，并且 $P$ 具有上面定义的性质 (P)。
\end{lemma}

\begin{proof}
由引理 \ref{dga-lemma-bimodule-over-tensor} 和 \ref{dga-lemma-resolve} 立即得出。
\end{proof}

\begin{lemma}
\label{dga-lemma-bimodule-property-P-sequence}
设 $R$ 是环，$(A,\text{d})$ 和 $(B,\text{d})$ 是微分分次 $R$-代数。
设 $P$ 是具有性质 (P) 的微分分次 $(A,B)$-双模，其相应滤过为
$F_\bullet$。那么得到短正合列
$$
0 \to
\bigoplus\nolimits F_iP \to
\bigoplus\nolimits F_iP \to P \to 0
$$
，它由微分分次 $(A,B)$-双模组成，并且作为分次 $(A,B)$-双模列是分裂的。
\end{lemma}

\begin{proof}
由引理 \ref{dga-lemma-bimodule-over-tensor} 和
\ref{dga-lemma-property-P-sequence} 立即得出。
\end{proof}















\section{双模与张量积}
\label{dga-section-bimodules-tensor}

\noindent
设 $R$ 是环，$A,B$ 是 $R$-代数，$M$ 是右 $A$-模，$N$ 是
$(A,B)$-双模。那么 $M\otimes_A N$ 是右 $B$-模。

\medskip\noindent
若在上一段的情形中，$A,B$ 是 $\mathbf{Z}$-分次代数，$M$ 是分次
$A$-模，$N$ 是分次 $(A,B)$-双模，则 $M\otimes_A N$ 是右分次 $B$-模。
该构造关于 $M$ 具有函子性，并定义例 \ref{dga-example-gm-gr-cat} 中分次范畴
之间的函子
$$
- \otimes_A N :
\text{Mod}^{gr}_A
\longrightarrow
\text{Mod}^{gr}_B
$$
。确实，若 $M,M'$ 是分次 $A$-模，$f:M\to M'$ 是次数 $n$ 齐次的
$A$-模同态，则
$f\otimes\text{id}_N:M\otimes_A N\to M'\otimes_A N$ 是次数 $n$
齐次的 $B$-模同态。

\medskip\noindent
若在上一段的情形中，$(A,\text{d})$ 和 $(B,\text{d})$ 是微分分次代数，
$M$ 是微分分次 $A$-模，$N$ 是微分分次 $(A,B)$-双模，则
$M\otimes_A N$ 是右微分分次 $B$-模。

\begin{lemma}
\label{dga-lemma-tensor}
设 $R$ 是环，$(A,\text{d})$ 和 $(B,\text{d})$ 是 $R$ 上的微分分次代数，
$N$ 是微分分次 $(A,B)$-双模。那么 $M\mapsto M\otimes_A N$ 定义微分分次
范畴之间的函子
$$
- \otimes_A N :
\text{Mod}^{dg}_{(A, \text{d})}
\longrightarrow
\text{Mod}^{dg}_{(B, \text{d})}
$$
。应用引理 \ref{dga-lemma-functorial}，此函子诱导函子
$$
\text{Mod}_{(A, \text{d})} \to \text{Mod}_{(B, \text{d})}
\quad\text{以及}\quad
K(\text{Mod}_{(A, \text{d})}) \to K(\text{Mod}_{(B, \text{d})})
$$
。
\end{lemma}

\begin{proof}
上面已经看到该构造如何定义底层分次范畴之间的函子。因此，只需证明该构造
与微分相容。设 $M,M'$ 是微分分次 $A$-模，$f:M\to M'$ 是次数 $n$
齐次的 $A$-模同态。那么
$$
\text{d}(f) = \text{d}_{M'} \circ f - (-1)^n f \circ \text{d}_M
$$
。另一方面，
$$
\text{d}(f \otimes \text{id}_N) =
\text{d}_{M' \otimes_A N} \circ
(f \otimes \text{id}_N)
- (-1)^n
(f \otimes \text{id}_N) \circ \text{d}_{M \otimes_A N}
$$
。把它应用于元素 $x\otimes y$，其中 $x\in M$、$y\in N$ 是齐次元素，得到
\begin{align*}
\text{d}(f \otimes \text{id}_N)(x \otimes y)
= &
\text{d}_{M'}(f(x)) \otimes y + (-1)^{n + \deg(x)}f(x) \otimes \text{d}_N(y) \\
& - (-1)^n f(\text{d}_M(x)) \otimes y
- (-1)^{n + \deg(x)}f(x) \otimes \text{d}_N(y) \\
= &
\text{d}(f) (x \otimes y)
\end{align*}
故 $\text{d}(f)\otimes\text{id}_N=
\text{d}(f\otimes\text{id}_N)$。证明完成。
\end{proof}

\begin{remark}
\label{dga-remark-shift-tensor-no-sign}
设 $R$ 是环，$(A,\text{d})$ 和 $(B,\text{d})$ 是 $R$ 上的微分分次代数，
$N$ 是微分分次 $(A,B)$-双模，$M$ 是右微分分次 $A$-模。那么对每个
$k\in\mathbf{Z}$，有右微分分次 $B$-模同构
$$
(M \otimes_A N)[k] \longrightarrow
M[k] \otimes_A N
$$
，其定义不引入符号；见《代数进阶》第
\ref{more-algebra-section-sign-rules} 节。
\end{remark}

\noindent
若有环 $R$、$R$-代数 $A,B,C$、右 $A$-模 $M$、$(A,B)$-双模 $N$
以及 $(B,C)$-双模 $N'$，则 $N\otimes_B N'$ 是 $(A,C)$-双模，并且
$$
(M \otimes_A N) \otimes_B N' = M \otimes_A (N \otimes_B N')
$$
。这个等式在分次情形和微分分次情形中仍成立。符号法则见《代数进阶》第
\ref{more-algebra-section-sign-rules} 节。






\section{双模与内部 Hom}
\label{dga-section-bimodules-hom}

\noindent
设 $R$ 是环。若 $A$ 是 $R$-代数（见第 \ref{dga-section-conventions} 节中的约定），
而 $M,M'$ 是右 $A$-模，则像通常一样定义
$$
\Hom_A(M, M') = \{f : M \to M' \mid f \text{ 是 }A\text{-线性的}\}
$$
。

\medskip\noindent
设 $R$ 是环，$A,B$ 是 $R$-代数，$N$ 是 $(A,B)$-双模，$N'$ 是右
$B$-模。在这种情形下，通过前复合把
$$
\Hom_B(N, N')
$$
看作右 $A$-模。

\medskip\noindent
设 $R$ 是环，$A,B$ 是 $\mathbf{Z}$-分次 $R$-代数，$N$ 是分次
$(A,B)$-双模，$N'$ 是右分次 $B$-模。在这种情形下，通过前复合把例
\ref{dga-example-gm-gr-cat} 中定义的分次 $R$-模
$$
\Hom_{\text{Mod}^{gr}_B}(N, N')
$$
看作右分次 $A$-模。该构造关于 $N'$ 具有函子性，并定义例
\ref{dga-example-gm-gr-cat} 中分次范畴之间的函子
$$
\Hom_{\text{Mod}^{gr}_B}(N, -) :
\text{Mod}^{gr}_B
\longrightarrow
\text{Mod}^{gr}_A
$$
。确实，若 $N_1,N_2$ 是分次 $B$-模，而 $f:N_1\to N_2$ 是次数 $n$
齐次的 $B$-模同态，则诱导映射
$\Hom_{\text{Mod}^{gr}_B}(N,N_1)\to
\Hom_{\text{Mod}^{gr}_B}(N,N_2)$ 是次数 $n$ 齐次的 $A$-模同态。

\medskip\noindent
设 $R$ 是环，$A,B$ 是微分 $\mathbf{Z}$-分次 $R$-代数，$N$ 是微分分次
$(A,B)$-双模，$N'$ 是右微分分次 $B$-模。在这种情形下，通过如分次情形
中的前复合，把例 \ref{dga-example-dgm-dg-cat} 中定义的微分分次 $R$-模
$$
\Hom_{\text{Mod}^{dg}_{(B, \text{d})}}(N, N')
$$
看作右微分分次 $A$-模。这与微分相容，因为乘法是复合
$$
\Hom_{\text{Mod}^{dg}_B}(N, N') \otimes_R A \to
\Hom_{\text{Mod}^{dg}_B}(N, N') \otimes_R
\Hom_{\text{Mod}^{dg}_B}(N, N) \to
\Hom_{\text{Mod}^{dg}_B}(N, N')
$$
。第一个箭头使用引理 \ref{dga-lemma-what-makes-a-bimodule-dg} 中的映射，
第二个箭头是微分分次范畴 $\text{Mod}^{dg}_{(B, \text{d})}$ 中的复合。

\begin{lemma}
\label{dga-lemma-hom}
设 $R$ 是环，$(A,\text{d})$ 和 $(B,\text{d})$ 是 $R$ 上的微分分次代数，
$N$ 是微分分次 $(A,B)$-双模。上述构造定义微分分次范畴之间的函子
$$
\Hom_{\text{Mod}^{dg}_{(B, \text{d})}}(N, -) :
\text{Mod}^{dg}_{(B, \text{d})}
\longrightarrow
\text{Mod}^{dg}_{(A, \text{d})}
$$
。应用引理 \ref{dga-lemma-functorial}，此函子诱导函子
$$
\text{Mod}_{(B, \text{d})} \to \text{Mod}_{(A, \text{d})}
\quad\text{以及}\quad
K(\text{Mod}_{(B, \text{d})}) \to K(\text{Mod}_{(A, \text{d})})
$$
。
\end{lemma}

\begin{proof}
上面已经看到该构造如何定义底层分次范畴之间的函子。因此，只需证明该构造
与微分相容。设 $N_1,N_2$ 是微分分次 $B$-模。记
$$
H_{12} = \Hom_{\text{Mod}^{dg}_{(B, \text{d})}}(N_1, N_2),\quad
H_1 = \Hom_{\text{Mod}^{dg}_{(B, \text{d})}}(N, N_1),\quad
H_2 = \Hom_{\text{Mod}^{dg}_{(B, \text{d})}}(N, N_2)
$$
。考虑微分分次范畴 $\text{Mod}^{dg}_{(B, \text{d})}$ 中的复合
$$
c : H_{12} \otimes_R H_1 \longrightarrow H_2
$$
。设 $f:N_1\to N_2$ 是次数 $n$ 齐次的 $B$-模同态；换言之，
$f\in H_{12}^n$。引理中的函子把 $f$ 映到
$c_f:H_1\to H_2$、$g\mapsto c(f,g)$。$\text{d}(f)$ 也类似。另一方面，
$$
\Hom_{\text{Mod}^{dg}_{(A, \text{d})}}(H_1, H_2)
$$
上的微分把 $c_f$ 映到
$\text{d}_{H_2}\circ c_f-(-1)^n c_f\circ\text{d}_{H_1}$。由于 $c$ 是
$R$-模复形之间的态射，有
$\text{d}c(f,g)=c(\text{d}f,g)+(-1)^n c(f,\text{d}g)$。因此
\begin{align*}
(\text{d}c_f)(g)
& =
\text{d}c(f,g) - (-1)^n c(f, \text{d}g) \\
& =
c(\text{d}f, g) + (-1)^n c(f, \text{d}g)  - (-1)^n c(f, \text{d}g) \\
& =
c(\text{d}f, g) = c_{\text{d}f}(g)
\end{align*}
。证明完成。
\end{proof}

\begin{remark}
\label{dga-remark-shift-hom-no-sign}
设 $R$ 是环，$(A,\text{d})$ 和 $(B,\text{d})$ 是 $R$ 上的微分分次代数，
$N$ 是微分分次 $(A,B)$-双模，$N'$ 是右微分分次 $B$-模。那么对每个
$k\in\mathbf{Z}$，有右微分分次 $A$-模同构
$$
\Hom_{\text{Mod}^{gr}_B}(N, N')[k]
\longrightarrow
\Hom_{\text{Mod}^{gr}_B}(N, N'[k])
$$
，其定义不引入符号；见《代数进阶》第
\ref{more-algebra-section-sign-rules} 节。
\end{remark}

\begin{lemma}
\label{dga-lemma-tensor-hom-adjunction}
设 $R$ 是环，$A,B$ 是 $R$-代数，$M$ 是右 $A$-模，$N$ 是
$(A,B)$-双模，$N'$ 是右 $B$-模。那么有 $R$-模的典范同构
$$
\Hom_B(M \otimes_A N, N') = \Hom_A(M, \Hom_B(N, N'))
$$
。若 $A,B,M,N,N'$ 带有彼此相容的分次，则有分次 $R$-模的典范同构
$$
\Hom_{\text{Mod}_B^{gr}}(M \otimes_A N, N') =
\Hom_{\text{Mod}_A^{gr}}(M, \Hom_{\text{Mod}_B^{gr}}(N, N'))
$$
。若 $A,B,M,N,N'$ 带有彼此相容的微分分次结构，则有 $R$-模复形的典范同构
$$
\Hom_{\text{Mod}^{dg}_{(B, \text{d})}}(M \otimes_A N, N') =
\Hom_{\text{Mod}^{dg}_{(A, \text{d})}}(M,
\Hom_{\text{Mod}^{dg}_{(B, \text{d})}}(N, N'))
$$
。
\end{lemma}

\begin{proof}
略。提示：在非分次情形下，把两边都解释为右侧 $B$-线性的 $A$-双线性映射
$\psi:M\times N\to N'$。在（微分）分次情形下，使用《代数进阶》引理
\ref{more-algebra-lemma-compose} 的同构，并验证它与模结构相容。或者，
使用引理 \ref{dga-lemma-characterize-hom} 的同构，并证明它与 $B$-模结构相容。
\end{proof}







\section{导出 Hom}
\label{dga-section-restriction}

\noindent
本节与《代数进阶》第 \ref{more-algebra-section-RHom} 节对应。

\medskip\noindent
设 $R$ 是环，$(A,\text{d})$ 和 $(B,\text{d})$ 是 $R$ 上的微分分次代数，
$N$ 是微分分次 $(A,B)$-双模。考虑第 \ref{dga-section-bimodules-hom} 节的函子
\begin{equation}
\label{dga-equation-restriction}
\Hom_{\text{Mod}^{dg}_{(B, \text{d})}}(N, -) :
\text{Mod}_{(B, \text{d})}
\longrightarrow
\text{Mod}_{(A, \text{d})}
\end{equation}
。

\begin{lemma}
\label{dga-lemma-restriction-homotopy}
函子 (\ref{dga-equation-restriction}) 定义三角范畴之间的正合函子
$K(\text{Mod}_{(B, \text{d})})\to K(\text{Mod}_{(A, \text{d})})$。
\end{lemma}

\begin{proof}
通过引理 \ref{dga-lemma-hom} 和评注 \ref{dga-remark-shift-hom-no-sign}，
这由引理 \ref{dga-lemma-functor-between-ABC} 的一般原理得出。
\end{proof}

\noindent
回忆，有三角范畴之间的正合函子
$$
\Hom_{\text{Mod}^{dg}_{(B, \text{d})}}(N, -) :
K(\text{Mod}_{(B, \text{d})}) \to K(\text{Mod}_{(A, \text{d})})
$$
；见引理 \ref{dga-lemma-restriction-homotopy}。考虑图表
$$
\xymatrix{
K(\text{Mod}_{(B, \text{d})}) \ar[d] \ar[rr]_{\text{见上文}} \ar[rrd]_F & &
K(\text{Mod}_{(A, \text{d})}) \ar[d] \\
D(B, \text{d}) \ar@{..>}[rr] & &
D(A, \text{d})
}
$$
希望把虚线箭头构造为复合 $F$ 的{\it 右导出函子}。
（{\it 警告}：在多数有意义的情形中，图表并不交换。）具体地，在《导出范畴》
第 \ref{derived-section-derived-functors} 节的一般框架中，希望计算 $F$
关于 $K(\text{Mod}_{(A, \text{d})})$ 中拟同构乘法系的右导出函子。

\begin{lemma}
\label{dga-lemma-derived-restriction}
在上述情形中，$F$ 的右导出函子存在。记为
$R\Hom(N,-):D(B,\text{d})\to D(A,\text{d})$。
\end{lemma}

\begin{proof}
使用《导出范畴》引理 \ref{derived-lemma-find-existence-computes} 来证明。
取具有性质 (I) 的对象作为对象族 $\mathcal{I}$。性质 (1) 已在引理
\ref{dga-lemma-right-resolution} 中证明。性质 (2) 成立，因为若
$s:I\to I'$ 是具有性质 (I) 的模之间的拟同构，则由引理
\ref{dga-lemma-hom-derived}，$s$ 是同伦等价。
\end{proof}

\begin{lemma}
\label{dga-lemma-functoriality-derived-restriction}
设 $R$ 是环，$(A,\text{d})$ 和 $(B,\text{d})$ 是微分分次 $R$-代数，
$f:N\to N'$ 是微分分次 $(A,B)$-双模同态。那么 $f$ 诱导函子态射
$$
- \circ f : R\Hom(N', -) \longrightarrow R\Hom(N, -)
$$
。若 $f$ 是拟同构，则 $f\circ-$ 是函子同构。
\end{lemma}

\begin{proof}
记 $\mathcal{B}=\text{Mod}^{dg}_{(B, \text{d})}$ 为微分分次 $B$-模的
微分分次范畴；见例 \ref{dga-example-dgm-dg-cat}。设 $I$ 是具有性质 (I) 的
微分分次 $B$-模。那么
$f\circ-:\Hom_\mathcal{B}(N',I)\to\Hom_\mathcal{B}(N,I)$ 是微分分次
$A$-模之间的映射，并且关于 $I$ 具有函子性。由于函子
$R\Hom(N',-)$ 和 $R\Hom(N,-)$ 通过把 $\Hom_\mathcal{B}$ 应用于具有
性质 (I) 的对象来计算（引理 \ref{dga-lemma-derived-restriction}），便得到所示的
函子变换。

\medskip\noindent
假设 $f$ 是拟同构。设 $F_\bullet$ 是 $I$ 上的给定滤过。由于
$I=\lim I/F_pI$，有
$\Hom_\mathcal{B}(N', I) = \lim \Hom_\mathcal{B}(N', I/F_pI)$，并且
$\Hom_\mathcal{B}(N,I)=\lim\Hom_\mathcal{B}(N,I/F_pI)$。由于系统
$I/F_pI$ 的过渡映射作为分次模映射是分裂的，系统
$\Hom_\mathcal{B}(N',I/F_pI)$ 和 $\Hom_\mathcal{B}(N,I/F_pI)$ 的过渡
映射都是满射。因此，把 $\Hom_\mathcal{B}(N',I)$（分别地，
$\Hom_\mathcal{B}(N,I)$）看作阿贝尔群复形，它计算复形系统
$\Hom_\mathcal{B}(N',I/F_pI)$（分别地，
$\Hom_\mathcal{B}(N,I/F_pI)$）的 $R\lim$。见《代数进阶》引理
\ref{more-algebra-lemma-compute-Rlim}。因此只需证明每个映射
$$
\Hom_\mathcal{B}(N', I/F_pI) \to \Hom_\mathcal{B}(N, I/F_pI)
$$
都是拟同构。由于满射 $I/F_{p+1}I\to I/F_pI$ 作为分次 $B$-模映射是
分裂的，故
$$
0 \to \Hom_\mathcal{B}(N', F_pI/F_{p + 1}I) \to
\Hom_\mathcal{B}(N', I/F_{p + 1}I) \to
\Hom_\mathcal{B}(N', I/F_pI) \to 0
$$
是微分分次 $A$-模的短正合列。对 $N$ 有类似的列，并且 $f$ 诱导短正合列
之间的映射。因此，对 $p$ 归纳（从 $p=0$ 开始，此时 $I/F_0I=0$），
得出只需证明映射
$\Hom_\mathcal{B}(N', F_pI/F_{p + 1}I) \to \Hom_\mathcal{B}(N, F_pI/F_{p + 1}I)$
是拟同构。由于 $F_pI/F_{p+1}I$ 是若干个 $A^\vee$ 的移位的直积，只需证明
$\Hom_\mathcal{B}(N', B^\vee[k]) \to \Hom_\mathcal{B}(N, B^\vee[k])$
是拟同构。由引理 \ref{dga-lemma-hom-into-shift-dual-free}，只需证明
$(N')^\vee\to N^\vee$ 是拟同构。这成立，因为 $f$ 是拟同构，且
$(\ )^\vee$ 是正合函子。
\end{proof}

\begin{lemma}
\label{dga-lemma-derived-restriction-exts}
设 $(A,\text{d})$ 和 $(B,\text{d})$ 是环 $R$ 上的微分分次代数，$N$
是微分分次 $(A,B)$-双模。那么对每个 $n\in\mathbf{Z}$，有 $R$-模同构
$$
H^n(R\Hom(N, M)) = \Ext^n_{D(B, \text{d})}(N, M)
$$
，它关于 $M$ 具有函子性。关于引理
\ref{dga-lemma-functoriality-derived-restriction} 中所述的运算，它也关于 $N$
具有函子性。
\end{lemma}

\begin{proof}
在引理 \ref{dga-lemma-derived-restriction} 的证明中已经看到
$$
R\Hom(N, M) = 
\Hom_{\text{Mod}^{dg}_{(B, \text{d})}}(N, I)
$$
作为微分分次 $A$-模成立，其中 $M\to I$ 是从 $M$ 到具有性质 (I) 的微分
分次 $B$-模的拟同构。因此，由引理 \ref{dga-lemma-hom-derived}，该复形具有
正确的上同调模。这些等同的函子性讨论从略。
\end{proof}

\begin{lemma}
\label{dga-lemma-compute-derived-restriction}
设 $R$ 是环，$(A,\text{d})$ 和 $(B,\text{d})$ 是微分分次 $R$-代数，
$N$ 是微分分次 $(A,B)$-双模。若
$\Hom_{D(B, \text{d})}(N, N') = \Hom_{K(\text{Mod}_{(B, \text{d})})}(N, N')$
对所有 $N'\in K(B,\text{d})$ 成立（例如，若 $N$ 作为微分分次 $B$-模
具有性质 (P)），则在 $D(B,\text{d})$ 中关于 $M$ 函子地有
$$
R\Hom(N, M) = \Hom_{\text{Mod}^{dg}_{(B, \text{d})}}(N, M)
$$
。
\end{lemma}

\begin{proof}
由构造（引理 \ref{dga-lemma-derived-restriction}），为求 $R\Hom(N,M)$，
选择拟同构 $M\to I$，其中 $I$ 是具有性质 (I) 的微分分次 $B$-模，并置
$R\Hom(N, M) = \Hom_{\text{Mod}^{dg}_{(B, \text{d})}}(N, I)$.
由假设，$M\to I$ 诱导的映射
$$
\Hom_{\text{Mod}^{dg}_{(B, \text{d})}}(N, M) \longrightarrow
\Hom_{\text{Mod}^{dg}_{(B, \text{d})}}(N, I)
$$
是拟同构；见例 \ref{dga-example-dgm-dg-cat} 中的讨论。这证明了本引理。
若 $N$ 作为 $B$-模具有性质 (P)，则由引理 \ref{dga-lemma-hom-derived}，
假设成立。
\end{proof}




\section{导出 Hom 的一种变体}
\label{dga-section-variant}

\noindent
设 $\mathcal{A}$ 是阿贝尔范畴。考虑 $\mathcal{A}$ 的复形所组成的微分分次
范畴 $\text{Comp}^{dg}(\mathcal{A})$；见例
\ref{dga-example-category-complexes}。设 $K^\bullet$ 是 $\mathcal{A}$ 的复形。置
$$
(E, \text{d}) = \Hom_{\text{Comp}^{dg}(\mathcal{A})}(K^\bullet, K^\bullet)
$$
，并考虑引理 \ref{dga-lemma-construction} 中微分分次范畴之间的函子
$$
\text{Comp}^{dg}(\mathcal{A}) \longrightarrow \text{Mod}^{dg}_{(E, \text{d})},
\quad
X^\bullet
\longmapsto
\Hom_{\text{Comp}^{dg}(\mathcal{A})}(K^\bullet, X^\bullet)
$$
。

\begin{lemma}
\label{dga-lemma-existence-of-derived}
在上述情形中，若
$\Hom(K^\bullet,-):K(\mathcal{A})\to D(\textit{Ab})$ 的右导出函子
$R\Hom(K^\bullet,-)$ 在 $D(\mathcal{A})$ 上处处有定义，则得到三角范畴之间
的典范正合函子
$$
R\Hom(K^\bullet, -) : D(\mathcal{A}) \longrightarrow D(E, \text{d})
$$
；取相应的阿贝尔群复形后，它化为通常的函子。
\end{lemma}

\begin{proof}
注意，由引理 \ref{dga-lemma-construction}，有相应函子
$K(\mathcal{A})\to K(\text{Mod}_{(E, \text{d})})$。断言它是三角范畴之间
的正合函子。确实，设 $f:A^\bullet\to B^\bullet$ 是 $\mathcal{A}$ 的
复形映射。计算表明
$$
\Hom_{\text{Comp}^{dg}(\mathcal{A})}(K^\bullet, C(f)^\bullet)
=
C\left(
\Hom_{\text{Comp}^{dg}(\mathcal{A})}(K^\bullet, A^\bullet) \to
\Hom_{\text{Comp}^{dg}(\mathcal{A})}(K^\bullet, B^\bullet)
\right)
$$
，其中右侧是本章前文在 $\text{Mod}_{(E, \text{d})}$ 中定义的锥。这表明
该函子与锥相容，因而与有区别三角相容。设 $X^\bullet$ 是
$K(\mathcal{A})$ 的对象。考虑拟同构 $s:X^\bullet\to Y^\bullet$ 的范畴。
已知函子
$(s : X^\bullet \to Y^\bullet) \mapsto \Hom_\mathcal{A}(K^\bullet, Y^\bullet)$
在 $D(\textit{Ab})$ 中看时本质上是常值的。但由于遗忘函子
$D(E,\text{d})\to D(\textit{Ab})$ 与取上同调相容，在
$D(E,\text{d})$ 中也有同样结论。这证明了本引理。
\end{proof}

\noindent
{\bf 警告：}虽然本引理按其表述成立，并可能按此形式发挥作用，但除非对所有
$n\in\mathbf{Z}$ 都有
$H^n(E) = \Ext^n_{D(\mathcal{A})}(K^\bullet, K^\bullet)$
，否则微分代数 $E$ 不是“正确”的微分代数。





\section{导出张量积}
\label{dga-section-base-change}

\noindent
本节与《代数进阶》第
\ref{more-algebra-section-derived-base-change} 节对应。

\medskip\noindent
设 $R$ 是环，$(A,\text{d})$ 和 $(B,\text{d})$ 是 $R$ 上的微分分次代数，
$N$ 是微分分次 $(A,B)$-双模。考虑第 \ref{dga-section-bimodules-tensor} 节
定义的函子
\begin{equation}
\label{dga-equation-bc}
\text{Mod}_{(A, \text{d})}
\longrightarrow
\text{Mod}_{(B, \text{d})},\quad
M \longmapsto M \otimes_A N
\end{equation}
。

\begin{lemma}
\label{dga-lemma-bc-homotopy}
函子 (\ref{dga-equation-bc}) 定义三角范畴之间的正合函子
$K(\text{Mod}_{(A, \text{d})})\to K(\text{Mod}_{(B, \text{d})})$。
\end{lemma}

\begin{proof}
通过引理 \ref{dga-lemma-tensor} 和评注 \ref{dga-remark-shift-tensor-no-sign}，
这由引理 \ref{dga-lemma-functor-between-ABC} 的一般原理得出。
\end{proof}

\noindent
此时可以考虑图表
$$
\xymatrix{
K(\text{Mod}_{(A, \text{d})}) \ar[d] \ar[rr]_{- \otimes_A N} \ar[rrd]_F & &
K(\text{Mod}_{(B, \text{d})}) \ar[d] \\
D(A, \text{d}) \ar@{..>}[rr] & &
D(B, \text{d})
}
$$
下文构造的虚线箭头将是复合 $F$ 的{\it 左导出函子}。
（{\it 警告}：该图表并不交换。）具体地，在《导出范畴》第
\ref{derived-section-derived-functors} 节的一般框架中，希望计算 $F$
关于 $K(\text{Mod}_{(A, \text{d})})$ 中拟同构乘法系的左导出函子。

\begin{lemma}
\label{dga-lemma-derived-bc}
在上述情形中，$F$ 的左导出函子存在。记为
$-\otimes_A^\mathbf{L}N:D(A,\text{d})\to D(B,\text{d})$。
\end{lemma}

\begin{proof}
使用《导出范畴》引理 \ref{derived-lemma-find-existence-computes} 来证明。
取具有性质 (P) 的对象作为对象族 $\mathcal{P}$。性质 (1) 已在引理
\ref{dga-lemma-resolve} 中证明。性质 (2) 成立，因为若 $s:P\to P'$ 是具有
性质 (P) 的模之间的拟同构，则由引理 \ref{dga-lemma-hom-derived}，$s$
是同伦等价。
\end{proof}

\begin{lemma}
\label{dga-lemma-functoriality-bc}
设 $R$ 是环，$(A,\text{d})$ 和 $(B,\text{d})$ 是微分分次 $R$-代数，
$f:N\to N'$ 是微分分次 $(A,B)$-双模同态。那么 $f$ 诱导函子态射
$$
1\otimes f :
- \otimes_A^\mathbf{L} N
\longrightarrow
- \otimes_A^\mathbf{L} N'
$$
。若 $f$ 是拟同构，则 $1\otimes f$ 是函子同构。
\end{lemma}

\begin{proof}
设 $M$ 是具有性质 (P) 的微分分次 $A$-模。那么
$1\otimes f:M\otimes_A N\to M\otimes_A N'$ 是微分分次 $B$-模之间的映射，
并且关于 $M$ 具有函子性。由于函子 $-\otimes_A^\mathbf{L}N$ 和
$-\otimes_A^\mathbf{L}N'$ 通过在具有性质 (P) 的对象上取张量来计算
（引理 \ref{dga-lemma-derived-bc}），便得到所示的函子变换。

\medskip\noindent
假设 $f$ 是拟同构。设 $F_\bullet$ 是 $M$ 上的给定滤过。注意
$M \otimes_A N = \colim F_i(M) \otimes_A N$，并且
$M \otimes_A N' = \colim F_i(M) \otimes_A N'$。因此只需证明
$F_n(M) \otimes_A N \to F_n(M) \otimes_A N'$
是拟同构（滤过余极限是正合的；见《代数》引理
\ref{algebra-lemma-directed-colimit-exact}）。由于包含映射
$F_n(M)\to F_{n+1}(M)$ 作为分次 $A$-模映射是分裂的，故
$$
0 \to F_n(M) \otimes_A N \to F_{n + 1}(M) \otimes_A N \to
F_{n + 1}(M)/F_n(M) \otimes_A N \to 0
$$
是微分分次 $B$-模的短正合列。对 $N'$ 有类似的列，并且 $f$ 诱导短正合列
之间的映射。因此，对 $n$ 归纳（从 $n=-1$ 开始，此时
$F_{-1}(M)=0$），得出只需证明映射
$F_{n+1}(M)/F_n(M)\otimes_A N\to F_{n+1}(M)/F_n(M)\otimes_A N'$
是拟同构。这成立，因为 $F_{n+1}(M)/F_n(M)$ 是若干个 $A$ 的移位的直和，
而对 $A[k]$ 结论成立，原因是 $f:N\to N'$ 是拟同构。
\end{proof}

\begin{lemma}
\label{dga-lemma-compute-bc}
设 $R$ 是环，$(A,\text{d})$ 和 $(B,\text{d})$ 是微分分次 $R$-代数。
设 $N$ 是微分分次 $(A,B)$-双模，并且作为左微分分次 $A$-模具有性质 (P)。
那么对所有微分分次 $A$-模 $M$，$M\otimes_A^\mathbf{L}N$ 由
$M\otimes_A N$ 计算。
\end{lemma}

\begin{proof}
设 $f:M\to M'$ 是微分分次 $A$-模之间的拟同构。断言
$f\otimes\text{id}:M\otimes_A N\to M'\otimes_A N$ 是拟同构。若断言成立，
则由引理 \ref{dga-lemma-derived-bc} 证明中导出张量积的构造，得到所需结论。
映射 $f\otimes\text{id}$ 的构造只依赖于 $N$ 上的左微分分次 $A$-模结构。
此外，
$M\otimes_A N=N\otimes_{A^{opp}}M=N\otimes_{A^{opp}}^\mathbf{L}M$，
因为 $N$ 作为微分分次 $A^{opp}$-模具有性质 (P)。因此，断言由引理
\ref{dga-lemma-functoriality-bc} 得出。
\end{proof}

\begin{lemma}
\label{dga-lemma-tensor-hom-adjoint}
设 $R$ 是环，$(A,\text{d})$ 和 $(B,\text{d})$ 是微分分次 $R$-代数，
$N$ 是微分分次 $(A,B)$-双模。那么引理 \ref{dga-lemma-derived-bc} 中的函子
$$
- \otimes_A^\mathbf{L} N : D(A, \text{d}) \longrightarrow D(B, \text{d})
$$
是引理 \ref{dga-lemma-derived-restriction} 中函子
$$
R\Hom(N, -) : D(B, \text{d}) \longrightarrow D(A, \text{d})
$$
的左伴随。
\end{lemma}

\begin{proof}
这由《导出范畴》引理
\ref{derived-lemma-pre-derived-adjoint-functors-general}，以及引理
\ref{dga-lemma-tensor-hom-adjunction} 中
$-\otimes_A N$ 与 $\Hom_{\text{Mod}^{dg}_{(B, \text{d})}}(N,-)$
互为伴随这一事实得出。
\end{proof}

\begin{example}
\label{dga-example-map-hom-tensor}
设 $R$ 是环，$(A,\text{d})\to(B,\text{d})$ 是微分分次 $R$-代数同态。
可以把 $B$ 看作微分分次 $(A,B)$-双模，并得到函子
$$
- \otimes_A B : D(A, \text{d}) \longrightarrow D(B, \text{d})
$$
。由引理 \ref{dga-lemma-tensor-hom-adjoint}，它的左伴随是函子
$R\Hom(B,-)$。对微分分次 $B$-模 $N$，以 $N_A$ 表示通过 $A\to B$
限制得到的微分分次 $A$-模。那么显然有关于 $B$-模 $N$ 函子的典范同构
$$
\Hom_{\text{Mod}^{dg}_{(B, \text{d})}}(B, N) \longrightarrow N_A,\quad
f \longmapsto f(1)
$$
。因此，$R\Hom(B,-)$ 是限制函子，并得到
$$
\Hom_{D(A, \text{d})}(M, N_A) =
\Hom_{D(B, \text{d})}(M \otimes^\mathbf{L}_A B, N)
$$
；它对 $M$ 与 $N$ 双函子地成立，正如交换环的情形。最后注意，限制也是张量函子，
因为 $N_A = N \otimes_B {}_BB_A = N \otimes_B^\mathbf{L} {}_BB_A$，
其中 ${}_BB_A$ 表示把 $B$ 看作微分分次 $(B,A)$-双模。
\end{example}

\begin{lemma}
\label{dga-lemma-tensor-with-compact-fully-faithful}
沿用引理 \ref{dga-lemma-tensor-hom-adjoint} 的记号与假设。假设
\begin{enumerate}
\item $N$ 在 $D(B,\text{d})$ 中定义一个紧对象；
\item 对每个 $k\in\mathbf{Z}$，映射
$H^k(A) \to \Hom_{D(B, \text{d})}(N, N[k])$ 都是同构。
\end{enumerate}
则函子 $-\otimes_A^\mathbf{L}N$ 是全忠实的。
\end{lemma}

\begin{proof}
由引理 \ref{dga-lemma-tensor-hom-adjoint}，我们的函子有由
$R\Hom(N,-)$ 给出的左伴随。由《范畴》引理
\ref{categories-lemma-adjoint-fully-faithful}，只需证明：对任意微分分次
$A$-模 $M$，映射
$$
M \longrightarrow R\Hom(N, M \otimes_A^\mathbf{L} N)
$$
在 $D(A,\text{d})$ 中是同构。为此，只需证明
$$
H^n(M) \longrightarrow
\text{Ext}^n_{D(B, \text{d})}(N, M \otimes_A^\mathbf{L} N)
$$
是同构，见引理 \ref{dga-lemma-derived-restriction-exts}。由于 $N$ 是紧对象，
右端与直和交换。因此由评注 \ref{dga-remark-P-resolution}，只需在
$M=A[k]$ 时证明此映射为同构。由评注
\ref{dga-remark-shift-tensor-no-sign}，
$(A[k]\otimes_A^\mathbf{L}N)=N[k]$；而对 $N$ 的假设 (2)
正是说结论在这些情形下成立。
\end{proof}

\begin{lemma}
\label{dga-lemma-base-change-K-flat}
设 $R\to R'$ 是环映射，$(A,\text{d})$ 是微分分次 $R$-代数，
而 $(A',\text{d})$ 是其基变换，即 $A'=A\otimes_RR'$。
若 $A$ 作为 $R$-模复形是 K-平坦的，则
\begin{enumerate}
\item $- \otimes_A^\mathbf{L} A' : D(A, \text{d}) \to D(A', \text{d})$
等于函子
$$
K(A, \text{d}) \longrightarrow K(A', \text{d}),\quad
M \longmapsto M \otimes_R R'
$$
的右导出函子；
\item 图表
$$
\xymatrix{
D(A, \text{d}) \ar[r]_{- \otimes_A^\mathbf{L} A'} \ar[d]_{\text{限制}} &
D(A', \text{d}) \ar[d]^{\text{限制}} \\
D(R) \ar[r]^{- \otimes_R^\mathbf{L} R'} & D(R')
}
$$
交换；
\item 若 $M$ 作为 $R$-模复形是 K-平坦的，则微分分次 $A'$-模
$M\otimes_RR'$ 表示 $M\otimes_A^\mathbf{L}A'$。
\end{enumerate}
\end{lemma}

\begin{proof}
对任意微分分次 $A$-模 $M$，都有典范映射
$$
c_M : M \otimes_R R' \longrightarrow M \otimes_A A'
$$
。设 $P$ 是具有性质 (P) 的微分分次 $A$-模。我们断言 $c_P$ 是同构，
并且 $P$ 作为 $R$-模复形是 K-平坦的。利用引理
\ref{dga-lemma-derived-bc} 和《代数进阶》第
\ref{more-algebra-section-derived-tensor-product} 节中导出张量积的定义，
这将通过形式论证证明引理所述全部结论。

\medskip\noindent
设 $F_\bullet$ 是表明 $P$ 具有性质 (P) 的滤过。注意 $c_A$ 是同构，
且依假设 $A$ 作为 $R$-模复形是 K-平坦的。因此，对 $A$ 的平移的直和
同样成立（处理直和时也可使用《代数进阶》引理
\ref{more-algebra-lemma-colimit-K-flat}）。于是，这对复形
$F_{p+1}P/F_pP$ 成立。由于短正合列
$$
0 \to F_pP \to F_{p + 1}P \to F_{p + 1}P/F_pP \to 0
$$
作为分次模列是分裂正合的，可归纳得出：对每个 $p$，$c_{F_pP}$ 都是同构，
且 $F_pP$ 作为 $R$-模复形是 K-平坦的（使用《代数进阶》引理
\ref{more-algebra-lemma-K-flat-two-out-of-three}）。最后，由
$P=\colim F_pP$ 可知 $c_P$ 是同构，且 $P$ 作为 $R$-模复形是
K-平坦的（使用《代数进阶》引理
\ref{more-algebra-lemma-colimit-K-flat}）。
\end{proof}

\begin{lemma}
\label{dga-lemma-RHom-is-tensor}
设 $R$ 是环，$(A,\text{d})$ 与 $(B,\text{d})$ 是微分分次
$R$-代数，$T$ 是微分分次 $(A,B)$-双模。假设
\begin{enumerate}
\item $T$ 在 $D(B,\text{d})$ 中定义一个紧对象；
\item $S = \Hom_{\text{Mod}^{dg}_{(B, \text{d})}}(T, B)$
在 $D(A,\text{d})$ 中表示 $R\Hom(T,B)$。
\end{enumerate}
则 $S$ 具有微分分次 $(B,A)$-双模结构，并有同构
$$
N \otimes_B^\mathbf{L} S \longrightarrow R\Hom(T, N)
$$
，它关于 $D(B,\text{d})$ 中的 $N$ 是函子的。
\end{lemma}

\begin{proof}
记 $\mathcal{B}=\text{Mod}^{dg}_{(B,\text{d})}$。$S$ 上的右
$A$-模结构来自映射 $A\to\Hom_\mathcal{B}(T,T)$ 以及例
\ref{dga-example-dgm-dg-cat} 中定义的复合
$\Hom_\mathcal{B}(T, B) \otimes \Hom_\mathcal{B}(T, T)
\to \Hom_\mathcal{B}(T, B)$。再次使用这一乘法，可得映射
$$
c_N :
N \otimes_B S = \Hom_\mathcal{B}(B, N) \otimes_B \Hom_\mathcal{B}(T, B)
\longrightarrow
\Hom_\mathcal{B}(T, N)
$$
，它关于 $N$ 是函子的。给定 $N$，可选取拟同构
$P\to N\to I$，其中 $P$（相应地 $I$）是具有性质 (P)（相应地 (I)）
的微分分次 $B$-模。于是利用 $c_N$，在分别表示
$S\otimes_B^\mathbf{L}N$ 与 $R\Hom(T,N)$ 的对象之间得到映射
$P\otimes_BS\to\Hom_\mathcal{B}(T,I)$。显然，这定义了引理所述的
函子变换 $c$。

\medskip\noindent
为证明 $c$ 是函子同构，可假设 $N$ 是具有性质 (P) 的微分分次
$B$-模。由于 $T$ 在 $D(B,\text{d})$ 中定义紧对象，且箭头两端都定义
三角范畴的正合函子，利用引理 \ref{dga-lemma-property-P-sequence}，
可约化到 $N$ 具有有限滤过、其分次片是若干 $B[k]$ 的直和的情形。
再次利用箭头两端都是三角范畴的正合函子以及 $T$ 的紧性，可进一步
约化到 $N=B[k]$ 的情形。假设 (2) 恰好断言此时 $c$ 是同构。
\end{proof}






\section{导出张量积的复合}
\label{dga-section-compose-tensor-functors}

\noindent
我们建议读者略过本节。

\medskip\noindent
设 $R$ 是环，$(A,\text{d})$、$(B,\text{d})$ 和 $(C,\text{d})$
是微分分次 $R$-代数。设 $N$ 是微分分次 $(A,B)$-双模，$N'$ 是微分分次
$(B,C)$-模。以 $N_B$ 表示把双模 $N$ 看作微分分次 $B$-模
（忘掉 $A$-结构）所得的模。存在典范映射
\begin{equation}
\label{dga-equation-plain-versus-derived}
N_B \otimes_B^\mathbf{L} N'
\longrightarrow
(N \otimes_B N')_C
\end{equation}
于 $D(C,\text{d})$ 中。这里，$(N\otimes_BN')_C$ 表示把
$(A,C)$-双模 $N\otimes_BN'$ 看作微分分次 $C$-模所得的模。
事实上，此映射来自如下事实：导出张量积总有到普通张量积的映射
（因为它是左导出函子）。

\begin{lemma}
\label{dga-lemma-compose-tensor-functors-general}
设 $R$ 是环，$(A,\text{d})$、$(B,\text{d})$ 和 $(C,\text{d})$
是微分分次 $R$-代数。设 $N$ 是微分分次 $(A,B)$-双模，$N'$ 是微分分次
$(B,C)$-模。假设 (\ref{dga-equation-plain-versus-derived}) 是同构。
则复合
$$
\xymatrix{
D(A, \text{d}) \ar[rr]^{- \otimes_A^\mathbf{L} N} & &
D(B, \text{d}) \ar[rr]^{- \otimes_B^\mathbf{L} N'} & &
D(C, \text{d})
}
$$
同构于 $-\otimes_A^\mathbf{L}N''$，其中 $N''=N\otimes_BN'$
被看作 $(A,C)$-双模。
\end{lemma}

\begin{proof}
我们定义函子变换
$$
(- \otimes_A^\mathbf{L} N) \otimes_B^\mathbf{L} N'
\longrightarrow
- \otimes_A^\mathbf{L} N''
$$
。为此，设 $M$ 是具有性质 (P) 的微分分次 $A$-模。按照引理
\ref{dga-lemma-derived-bc} 的证明中函子 $-\otimes_A^\mathbf{L}N''$
的构造，普通张量积 $M\otimes_AN''$ 在 $D(C,\text{d})$ 中表示
$M\otimes_A^\mathbf{L}N''$。于是可写成
$$
M \otimes_A N'' =
M \otimes_A (N \otimes_B N') =
(M \otimes_A N) \otimes_B N'
$$
。模 $M\otimes_AN$ 在 $D(B,\text{d})$ 中表示
$M\otimes_A^\mathbf{L}N$。选取拟同构 $Q\to M\otimes_AN$，
其中 $Q$ 是具有性质 (P) 的微分分次 $B$-模。则 $Q\otimes_BN'$
在 $D(C,\text{d})$ 中表示
$(M\otimes_A^\mathbf{L}N)\otimes_B^\mathbf{L}N'$。
因此，可由下式定义所需映射：
$$
(M \otimes_A^\mathbf{L} N) \otimes_B^\mathbf{L} N' =
Q \otimes_B N' \to
M \otimes_A N \otimes_B N' =
M \otimes_A^\mathbf{L} N''
$$
。此映射的构造关于 $M$ 是函子的，并与有区别三角及直和相容；细节从略。
考虑 $D(A,\text{d})$ 的对象 $M$ 的性质 $T$，其含义是此映射为同构。则
\begin{enumerate}
\item 若每个 $M_i$ 都满足 $T$，则 $\bigoplus M_i$ 满足 $T$；
\item 对一个有区别三角，若三个对象中的两个满足 $T$，则第三个也满足；
\item $A[k]$ 满足 $T$，因为此时所得映射是
(\ref{dga-equation-plain-versus-derived}) 的一个平移，而我们已经假设后者是同构。
\end{enumerate}
因此由评注 \ref{dga-remark-P-resolution}，性质 $T$ 恒成立，证明完毕。
\end{proof}

\noindent
设 $R$ 是环，$(A,\text{d})$、$(B,\text{d})$ 和 $(C,\text{d})$
是微分分次 $R$-代数。暂以 $(A\otimes_RB)_B$ 表示把微分分次代数
$A\otimes_RB$ 看作（右）微分分次 $B$-模所得的模；以
${}_B(B\otimes_RC)_C$ 表示把微分分次代数 $B\otimes_RC$
看作微分分次 $(B,C)$-双模所得的双模。于是存在典范映射
\begin{equation}
\label{dga-equation-plain-versus-derived-algebras}
(A \otimes_R B)_B \otimes_B^\mathbf{L} {}_B(B \otimes_R C)_C
\longrightarrow
(A \otimes_R B \otimes_R C)_C
\end{equation}
于 $D(C,\text{d})$ 中，其中 $(A\otimes_RB\otimes_RC)_C$
表示把微分分次 $R$-代数 $A\otimes_RB\otimes_RC$ 看作（右）
微分分次 $C$-模所得的模。事实上，此映射来自等同
$$
(A \otimes_R B)_B \otimes_B {}_B(B \otimes_R C)_C =
(A \otimes_R B \otimes_R C)_C
$$
以及导出张量积总有到普通张量积的映射这一事实
（因为它是左导出函子）。

\begin{lemma}
\label{dga-lemma-compose-tensor-functors-general-algebra}
设 $R$ 是环，$(A,\text{d})$、$(B,\text{d})$ 和 $(C,\text{d})$
是微分分次 $R$-代数。假设
(\ref{dga-equation-plain-versus-derived-algebras}) 是同构。设 $N$ 是微分分次
$(A,B)$-双模，$N'$ 是微分分次 $(B,C)$-双模。则复合
$$
\xymatrix{
D(A, \text{d}) \ar[rr]^{- \otimes_A^\mathbf{L} N} & &
D(B, \text{d}) \ar[rr]^{- \otimes_B^\mathbf{L} N'} & &
D(C, \text{d})
}
$$
同构于 $-\otimes_A^\mathbf{L}N''$，其中微分分次 $(A,C)$-双模
$N''$ 将在证明中描述。
\end{lemma}

\begin{proof}
由引理 \ref{dga-lemma-functoriality-bc}，可把 $N$ 与 $N'$ 替换为与之
拟同构的双模。因此可假设 $N$（相应地 $N'$）作为微分分次
$(A,B)$-双模（相应地 $(B,C)$-双模）具有性质 (P)，见引理
\ref{dga-lemma-bimodule-resolve}。我们断言，对 $(A,C)$-双模
$N''=N\otimes_BN'$，引理成立。为证明这一点，只需证明
$$
N_B \otimes_B^\mathbf{L} N' \longrightarrow (N \otimes_B N')_C
$$
在 $D(C,\text{d})$ 中是同构，见引理
\ref{dga-lemma-compose-tensor-functors-general}。

\medskip\noindent
设 $F_\bullet$ 是双模性质 (P) 中 $N$ 上的滤过。由引理
\ref{dga-lemma-bimodule-property-P-sequence}，存在微分分次
$(A,B)$-双模的短正合列
$$
0 \to
\bigoplus\nolimits F_iN \to
\bigoplus\nolimits F_iN \to N \to 0
$$
，它作为分次 $(A,B)$-双模列是分裂的。因而它更是微分分次 $B$-模的
容许短正合列，并给出 $D(B,\text{d})$ 中的有区别三角
$$
\bigoplus\nolimits F_iN_B \to
\bigoplus\nolimits F_iN_B \to N_B \to
\bigoplus\nolimits F_iN_B[1]
$$
。利用 $-\otimes_B^\mathbf{L}N'$ 是三角范畴的正合函子且与直和交换，
并利用 $-\otimes_BN'$ 把容许正合列变为容许正合列且与直和交换，
可将问题约化为证明：对每个 $p$，
$$
(F_pN)_B \otimes_B^\mathbf{L} N' \longrightarrow (F_pN)_B \otimes_B N'
$$
都是拟同构。对 $(A,B)$-双模的短正合列
$$
0 \to F_pN \to F_{p + 1}N \to F_{p + 1}N/F_pN \to 0
$$
重复此论证；这些列作为分次 $(A,B)$-双模列是分裂的，故可约化为对
$F_{p+1}N/F_pN$ 证明同一断言。由于这些模是
$(A\otimes_RB)_B$ 的平移的直和，又可约化为证明
$$
(A \otimes_R B)_B \otimes_B^\mathbf{L} N'
\longrightarrow
(A \otimes_R B)_B \otimes_B N'
$$
是拟同构。

\medskip\noindent
选取双模性质 (P) 中 $N'$ 上的滤过 $F_\bullet$。选取微分分次
$B$-模的拟同构 $P\to(A\otimes_RB)_B$，其中 $P$ 具有性质 (P)。
由引理 \ref{dga-lemma-derived-bc} 中的构造，$P\otimes_BN'$ 在
$D(C,\text{d})$ 中表示
$(A\otimes_RB)_B\otimes_B^\mathbf{L}N'$；因此我们必须证明
$P\otimes_BN'\to(A\otimes_RB)_B\otimes_BN'$ 是拟同构。
由于 $N'=\colim F_pN'$，只需证明
$P\otimes_BF_pN'\to(A\otimes_RB)_B\otimes_BF_pN'$
是拟同构。利用短正合列
$0 \to F_pN' \to F_{p + 1}N' \to F_{p + 1}N'/F_pN' \to 0$
作为分次 $(B,C)$-双模列是分裂的，可约化为证明：对每个 $p$，
$P \otimes_B F_{p + 1}N'/F_pN' \to
(A \otimes_R B)_B \otimes_B F_{p + 1}N'/F_pN'$
是拟同构。最后，利用 $F_{p+1}N'/F_pN'$ 是
${}_B(B\otimes_RC)_C$ 的平移的直和，可知只需证明
$$
P \otimes_B {}_B(B \otimes_R C)_C \to
(A \otimes_R B)_B \otimes_B {}_B(B \otimes_R C)_C
$$
是拟同构。由于 $P\to(A\otimes_RB)_B$ 是由满足性质 (P) 的模
给出的分辨，此微分分次 $C$-模映射在 $D(C,\text{d})$ 中表示态射
(\ref{dga-equation-plain-versus-derived-algebras})，证明完毕。
\end{proof}

\begin{lemma}
\label{dga-lemma-compose-tensor-functors}
设 $R$ 是环，$(A,\text{d})$、$(B,\text{d})$ 和 $(C,\text{d})$
是微分分次 $R$-代数。若 $C$ 作为 $R$-模复形是 K-平坦的，则
(\ref{dga-equation-plain-versus-derived-algebras}) 是同构，并且引理
\ref{dga-lemma-compose-tensor-functors-general-algebra} 的结论成立。
\end{lemma}

\begin{proof}
选取微分分次 $B$-模的拟同构 $P\to(A\otimes_RB)_B$，其中 $P$
具有性质 (P)。于是只需证明
$$
P \otimes_B (B \otimes_R C) \longrightarrow
(A \otimes_R B) \otimes_B (B \otimes_R C)
$$
是拟同构。等价地，我们考察
$$
P \otimes_R C \longrightarrow
A \otimes_R B \otimes_R C
$$
。若 $C$ 作为 $R$-模复形是 K-平坦的，则由《代数进阶》引理
\ref{more-algebra-lemma-K-flat-quasi-isomorphism}，这是拟同构。
\end{proof}





\section{导出张量积的一种变体}
\label{dga-section-variant-base-change}

\noindent
设 $(\mathcal{C},\mathcal{O})$ 是环化位点。于是有函子
$$
\text{Comp}(\mathcal{O}) \to K(\mathcal{O}) \to D(\mathcal{O})
$$
；如上所见，还有微分分次增强 $\text{Comp}^{dg}(\mathcal{O})$。
具体地说，这是例 \ref{dga-example-category-complexes} 中与阿贝尔范畴
$\textit{Mod}(\mathcal{O})$ 相关联的微分分次范畴。设 $K^\bullet$
是 $\mathcal{O}$-模复形，换言之，是 $\text{Comp}^{dg}(\mathcal{O})$
的对象。令
$$
(E, \text{d}) =
\Hom_{\text{Comp}^{dg}(\mathcal{O})}(K^\bullet, K^\bullet)
$$
。这是微分分次 $\mathbf{Z}$-代数。我们断言，在这种情形下有导出基变换
的类似物。

\begin{lemma}
\label{dga-lemma-tensor-with-complex}
在上述情形下，存在微分分次范畴的函子
$$
- \otimes_E K^\bullet :
\text{Mod}^{dg}_{(E, \text{d})}
\longrightarrow
\text{Comp}^{dg}(\mathcal{O})
$$
。此函子把 $E$ 映到 $K^\bullet$，并与直和交换。
\end{lemma}

\begin{proof}
设 $M$ 是微分分次 $E$-模。对 $\mathcal{C}$ 的每个对象 $U$，复形
$K^\bullet(U)$ 既是左微分分次 $E$-模，也是右 $\mathcal{O}(U)$-模。
两个作用交换，故得到一个双模。因此，由第 \ref{dga-section-tensor-product}
节与第 \ref{dga-section-bimodules} 节中的构造，可以形成张量积
$$
M \otimes_E K^\bullet(U)
$$
；它是微分分次 $\mathcal{O}(U)$-模，即 $\mathcal{O}(U)$-模复形。
此构造关于 $U$ 是函子的，故可层化而得到 $\mathcal{O}$-模复形，记为
$$
M \otimes_E K^\bullet
$$
。此外，对每个 $U$，由引理 \ref{dga-lemma-tensor}，此构造确定微分分次
范畴的函子
$\text{Mod}^{dg}_{(E, \text{d})} \to \text{Comp}^{dg}(\mathcal{O}(U))$
。因此，显然得到引理所述函子。
\end{proof}

\begin{lemma}
\label{dga-lemma-tensor-with-complex-homotopy}
引理 \ref{dga-lemma-tensor-with-complex} 的函子定义三角范畴的正合函子
$K(\text{Mod}_{(E, \text{d})}) \to K(\mathcal{O})$.
\end{lemma}

\begin{proof}
由引理 \ref{dga-lemma-functorial}，该函子诱导同伦范畴之间的函子。
我们必须证明 $-\otimes_EK^\bullet$ 把有区别三角变为有区别三角。
设 $0\to K\to L\to M\to0$ 是微分分次 $E$-模的容许短正合列。
设 $s:M\to L$ 是分次 $E$-模同态，并且是 $L\to M$ 的左逆。
则 $s$ 定义分次 $\mathcal{O}$-模的映射
$M\otimes_EK^\bullet\to L\otimes_EK^\bullet$
（即保持 $\mathcal{O}$-模结构与分次，但不保持微分），它是
$L\otimes_EK^\bullet\to M\otimes_EK^\bullet$ 的左逆。因此
$$
0 \to K \otimes_E K^\bullet \to L \otimes_E K^\bullet \to
M \otimes_E K^\bullet \to 0
$$
是逐项分裂的复形短正合列，也就是说，它在 $K(\mathcal{O})$ 中定义
一个有区别三角。
\end{proof}

\begin{lemma}
\label{dga-lemma-tensor-with-complex-derived}
引理 \ref{dga-lemma-tensor-with-complex-homotopy} 的函子
$K(\text{Mod}_{(E, \text{d})})\to K(\mathcal{O})$ 有定义在整个
$D(E,\text{d})$ 上的左导出版本，记为
$- \otimes_E^\mathbf{L} K^\bullet : D(E, \text{d}) \to D(\mathcal{O})$.
\end{lemma}

\begin{proof}
使用《导出范畴》引理 \ref{derived-lemma-find-existence-computes}
来证明。取具有性质 (P) 的对象为对象族 $\mathcal{P}$。
性质 (1) 已在引理 \ref{dga-lemma-resolve} 中证明。性质 (2) 成立，
因为若 $s:P\to P'$ 是具有性质 (P) 的模之间的拟同构，则由引理
\ref{dga-lemma-hom-derived}，$s$ 是同伦等价。
\end{proof}

\begin{lemma}
\label{dga-lemma-upgrade-tensor-with-complex-derived}
设 $R$ 是环，$\mathcal{C}$ 是位点，$\mathcal{O}$ 是交换
$R$-代数层，$K^\bullet$ 是 $\mathcal{O}$-模复形。引理
\ref{dga-lemma-tensor-with-complex-derived} 的函子具有如下性质：对
$D(E,\text{d})$ 中的任意 $M,N$，在 $D(R)$ 中有典范映射
$$
R\Hom(M, N)
\longrightarrow
R\Hom_\mathcal{O}(M \otimes_E^\mathbf{L} K^\bullet,
N \otimes_E^\mathbf{L} K^\bullet)
$$
；它在上同调模上给出由函子 $-\otimes_E^\mathbf{L}K^\bullet$ 诱导的映射
$$
\Ext^n_{D(E, \text{d})}(M, N) \to
\Ext^n_{D(\mathcal{O})}
(M \otimes_E^\mathbf{L} K^\bullet, N \otimes_E^\mathbf{L} K^\bullet)
$$
。
\end{lemma}

\begin{proof}
箭头右端是《位点上的上同调》第
\ref{sites-cohomology-section-global-RHom} 节引入的全局导出 Hom，
它具有正确的上同调模。对左端，把 $M$ 看作 $(R,A)$-双模，并使用第
\ref{dga-section-restriction} 节引入的导出 $\Hom$；它也具有正确的上同调模。
为证明引理，可假设 $M,N$ 都是具有性质 (P) 的微分分次 $E$-模；
由引理 \ref{dga-lemma-functoriality-derived-restriction}，这不改变箭头左端。
由引理 \ref{dga-lemma-compute-derived-restriction}，箭头左端于是变为
$\Hom_{\text{Mod}^{dg}_{(B, \text{d})}}(M, N)$.
在引理 \ref{dga-lemma-tensor-with-complex}、
\ref{dga-lemma-tensor-with-complex-homotopy} 与
\ref{dga-lemma-tensor-with-complex-derived}
中，我们已经构造微分分次范畴的函子
$$
- \otimes_E K^\bullet :
\text{Mod}^{dg}_{(E, \text{d})}
\longrightarrow
\text{Comp}^{dg}(\mathcal{O})
$$
，并已证明，在具有性质 (P) 的微分分次 $E$-模上取此函子的值即可计算
$-\otimes_E^\mathbf{L}K^\bullet$。因此得到 $R$-模复形的映射
$$
\Hom_{\text{Mod}^{dg}_{(B, \text{d})}}(M, N)
\longrightarrow
\Hom_{\text{Comp}^{dg}(\mathcal{O})}
(M \otimes_E K^\bullet, N \otimes_E K^\bullet)
$$
。对任意 $\mathcal{O}$-模复形 $\mathcal{F}^\bullet$、
$\mathcal{G}^\bullet$，存在典范映射
$$
\Hom_{\text{Comp}^{dg}(\mathcal{O})}
(\mathcal{F}^\bullet, \mathcal{G}^\bullet) =
\Gamma(\mathcal{C},
\SheafHom^\bullet(\mathcal{F}^\bullet, \mathcal{G}^\bullet))
\longrightarrow
R\Hom_\mathcal{O}(\mathcal{F}^\bullet, \mathcal{G}^\bullet).
$$
合并这些映射，即得引理所需映射。
\end{proof}

\begin{lemma}
\label{dga-lemma-tensor-with-complex-hom-adjoint}
设 $(\mathcal{C},\mathcal{O})$ 是环化位点，$K^\bullet$ 是
$\mathcal{O}$-模复形。则引理 \ref{dga-lemma-tensor-with-complex-derived}
的函子
$$
- \otimes_E^\mathbf{L} K^\bullet :
D(E, \text{d})
\longrightarrow
D(\mathcal{O})
$$
是引理 \ref{dga-lemma-existence-of-derived} 的函子
$$
R\Hom(K^\bullet, -) : D(\mathcal{O}) \longrightarrow D(E, \text{d})
$$
的左伴随。
\end{lemma}

\begin{proof}
此断言的含义是，存在关于 $M$ 与 $L^\bullet$ 双函子的等式
$$
\Hom_{D(E, \text{d})}(M, R\Hom(K^\bullet, L^\bullet)) =
\Hom_{D(\mathcal{O})}(M \otimes^\mathbf{L}_E K^\bullet, L^\bullet)
$$
。为说明这一点，可把 $M$ 替换为具有性质 (P) 的微分分次 $E$-模 $P$，
也可把 $L^\bullet$ 替换为 K-内射的 $\mathcal{O}$-模复形
$I^\bullet$。结合引理 \ref{dga-lemma-hom-derived} 与陈述中所引诸引理给出的
导出函子计算，上式化为
$$
\Hom_{K(\text{Mod}_{(E, \text{d})})}
(P, \Hom_\mathcal{B}(K^\bullet, I^\bullet)) =
\Hom_{K(\mathcal{O})}(P \otimes_E K^\bullet, I^\bullet)
$$
，其中 $\mathcal{B}=\text{Comp}^{dg}(\mathcal{O})$。存在从右端到左端、
关于 $P$ 与 $I^\bullet$ 函子的求值映射（细节从略）。选取性质 (P)
定义中 $P$ 上的滤过 $F_\bullet$。由引理
\ref{dga-lemma-property-P-sequence}，并利用等式两端在
$K(\text{Mod}_{(E,\text{d})})$ 上都是关于 $P$ 的同调函子，
可约化到把 $P$ 换成微分分次 $E$-模 $\bigoplus F_iP$ 的情形。
由于两端都把变量 $P$ 中的直和变为直积，可再约化到 $P$ 是某个
微分分次 $E$-模 $F_iP$ 的情形。每个 $F_iP$ 都有由容许单态射给出的
有限滤过，其分次片是分次投射 $E$-模，故可约化到 $P$ 是分次投射
$E$-模的情形。此时显然有
$$
\Hom_{\text{Mod}^{dg}_{(E, \text{d})}}
(P, \Hom_\mathcal{B}(K^\bullet, I^\bullet)) =
\Hom_{\text{Comp}^{dg}(\mathcal{O})}(P \otimes_E K^\bullet, I^\bullet)
$$
作为分次 $\mathbf{Z}$-模的等式（因为此断言可约化到 $P=E[k]$
这一显然的情形）。该同构与微分相容，故得结论。
\end{proof}

\begin{lemma}
\label{dga-lemma-fully-faithful-in-compact-case}
设 $(\mathcal{C},\mathcal{O})$ 是环化位点，$K^\bullet$ 是
$\mathcal{O}$-模复形。假设
\begin{enumerate}
\item $K^\bullet$ 表示 $D(\mathcal{O})$ 的一个紧对象；
\item $E = \Hom_{\text{Comp}^{dg}(\mathcal{O})}(K^\bullet, K^\bullet)$
计算 $K^\bullet$ 在 $D(\mathcal{O})$ 中的 Ext 群。
\end{enumerate}
则引理 \ref{dga-lemma-tensor-with-complex-derived} 的函子
$$
- \otimes_E^\mathbf{L} K^\bullet :
D(E, \text{d})
\longrightarrow
D(\mathcal{O})
$$
是全忠实的。
\end{lemma}

\begin{proof}
由引理 \ref{dga-lemma-tensor-with-complex-hom-adjoint}，我们的函子有由
$R\Hom(K^\bullet,-)$ 给出的左伴随。因此只需证明：对微分分次
$E$-模 $M$，映射
$$
H^0(M) \longrightarrow
\Hom_{D(\mathcal{O})}(K^\bullet, M \otimes_E^\mathbf{L} K^\bullet)
$$
是同构。可假设 $M=P$ 是具有性质 (P) 的微分分次 $E$-模。
由于 $K^\bullet$ 定义紧对象，利用引理
\ref{dga-lemma-property-P-sequence}，可约化到 $P$ 具有有限滤过且其分次片
是若干 $E[k]$ 的直和的情形。再次利用紧性，可约化到 $P=E[k]$
的情形。对 $K^\bullet$ 的假设正是说结论在这些情形下成立。
\end{proof}











\section{紧对象的刻画}
\label{dga-section-compact}

\noindent
加性范畴的紧对象定义于《导出范畴》定义
\ref{derived-definition-compact-object}。本节刻画微分分次代数的导出范畴
中的紧对象。

\begin{remark}
\label{dga-remark-source-graded-projective}
设 $(A,\text{d})$ 是微分分次代数。能否刻画满足下列条件的微分分次
$A$-模 $P$：对每个微分分次 $A$-模 $M$，都有
$$
\Hom_{K(A, \text{d})}(P, M) =  \Hom_{D(A, \text{d})}(P, M)
$$
？设 $\mathcal{D}\subset K(A,\text{d})$ 是以满足上述条件的对象 $P$
为对象的满子范畴。则 $\mathcal{D}$ 是 $K(A,\text{d})$ 的严格满饱和
三角子范畴。若 $P$ 作为分次 $A$-模是投射的，则要判定 $P$ 是否为
$\mathcal{D}$ 的对象，只需检验：每当 $M$ 非循环时，
$\Hom_{K(A,\text{d})}(P,M)=0$。然而，一般来说，仅假设 $P$
作为分次 $A$-模是投射的还不够。例如，取 $A=R=k[\epsilon]$，其中
$k$ 是域，$k[\epsilon]=k[x]/(x^2)$ 是对偶数环。令 $P$ 为如下对象：
对每个 $n\in\mathbf{Z}$ 都有 $P^n=R$，且微分由乘以 $\epsilon$
给出。于是 $\text{id}_P\in\Hom_{K(A,\text{d})}(P,P)$ 是非零元素，
但 $P$ 非循环。
\end{remark}

\begin{remark}
\label{dga-remark-graded-projective-is-compact}
设 $(A,\text{d})$ 是微分分次代数。若微分分次 $A$-模 $M$ 作为右
$A$-模由有限多个元素生成，就称 $M$ 是{\it 有限的}。若微分分次
$A$-模 $P$ 是有限分次投射的，可以问：$P$ 是否给出
$D(A,\text{d})$ 的紧对象？推测一般并非如此，但我们不知道反例。
然而，若 $P$ 还是评注 \ref{dga-remark-source-graded-projective} 的范畴
$\mathcal{D}$ 的对象，则答案是肯定的（这来自如下事实：
$D(A,\text{d})$ 中的直和由模的直和给出；细节从略）。
\end{remark}

\begin{lemma}
\label{dga-lemma-factor-through-nicer}
设 $(A,\text{d})$ 是微分分次代数，$E$ 是 $D(A,\text{d})$ 的紧对象。
设 $P$ 是微分分次 $A$-模，它具有有限滤过
$$
0 = F_{-1}P \subset F_0P \subset F_1P \subset \ldots \subset F_nP = P
$$
，其各项为微分分次子模，并且对某些集合 $J_i$ 与整数 $k_{i,j}$ 有
$$
F_{i + 1}P/F_iP \cong \bigoplus\nolimits_{j \in J_i} A[k_{i, j}]
$$
作为微分分次 $A$-模的同构。设 $E\to P$ 是 $D(A,\text{d})$ 的态射。
则存在微分分次子模 $P'\subset P$，使得对某些有限子集
$J'_i\subset J_i$，$F_{i+1}P\cap P'/(F_iP\cap P')$ 等于
$\bigoplus_{j\in J'_i}A[k_{i,j}]$，并且 $E\to P$ 通过 $P'$ 分解。
\end{lemma}

\begin{proof}
我们对 $-1\leq m\leq n$ 向下归纳，证明存在微分分次子模
$P'\subset P$ 使得
\begin{enumerate}
\item $F_mP\subset P'$；
\item 对 $i\geq m$，商 $F_{i+1}P\cap P'/(F_iP\cap P')$ 同构于
$\bigoplus_{j\in J'_i}A[k_{i,j}]$，其中 $J'_i\subset J_i$ 为有限子集；
\item $E\to P$ 通过 $P'$ 分解。
\end{enumerate}
初始情形是 $m=n$，此时可取 $P'=P$。

\medskip\noindent
归纳步骤。假设 $P'$ 对 $m$ 有效。对 $i\geq m$ 与 $j\in J'_i$，
取次数为 $k_{i,j}$ 的齐次元素
$x_{i,j}\in F_{i+1}P\cap P'$，使其在
$F_{i+1}P\cap P'/(F_iP\cap P')$ 中的像是对应于 $j\in J_i$
的直和项的生成元。诸 $x_{i,j}$ 作为 $A$-模生成 $P'/F_mP$。写成
$$
\text{d}(x_{i, j}) = \sum x_{i', j'} a_{i, j}^{i', j'} + y_{i, j}
$$
，其中 $y_{i,j}\in F_mP$ 且 $a_{i,j}^{i',j'}\in A$。存在有限子集
$J'_{m-1}\subset J_{m-1}$，使每个 $y_{i,j}$ 都映到
$F_mP/F_{m-1}P$ 的子模
$\bigoplus_{j\in J'_{m-1}}A[k_{m-1,j}]$ 中。令
$P''\subset F_mP$ 为该子模在映射 $F_mP\to F_mP/F_{m-1}P$
下的逆像。于是可见 $A$-子模
$$
P'' + \sum x_{i, j}A
$$
是所求类型的微分分次子模。此外，
$$
P'/(P'' + \sum x_{i, j}A) =
\bigoplus\nolimits_{j \in J_{m - 1} \setminus J'_{m - 1}} A[k_{m - 1, j}]
$$
。由于 $E$ 是紧的，给定映射 $E\to P'$ 与商映射的复合通过上面所示模
的一个有限直和子项分解。因此，扩大 $J'_{m-1}$ 后，可假设
$E\to P'$ 通过 $P''+\sum x_{i,j}A$ 分解，正合所需。
\end{proof}

\noindent
并非 $D(A,\text{d})$ 的每个紧对象都来自有限分次投射的微分分次
$A$-模；见《例》第 \ref{examples-section-interesting-compact} 节。

\begin{proposition}
\label{dga-proposition-compact}
设 $(A,\text{d})$ 是微分分次代数，$E$ 是 $D(A,\text{d})$ 的对象。
则下列条件等价：
\begin{enumerate}
\item $E$ 是紧对象；
\item $E$ 是 $D(A,\text{d})$ 的某个对象的直和项，该对象由微分分次模
$P$ 表示；$P$ 有由微分分次子模组成的有限滤过 $F_\bullet$，并且
$F_iP/F_{i-1}P$ 是 $A$ 的平移的有限直和。
\end{enumerate}
\end{proposition}

\begin{proof}
假设 $E$ 是紧的。由引理 \ref{dga-lemma-resolve}，可假设 $E$ 由具有性质
(P) 的微分分次 $A$-模 $P$ 表示。考虑由引理
\ref{dga-lemma-property-P-sequence} 的容许短正合列给出的有区别三角
$$
\bigoplus F_iP \to \bigoplus F_iP \to P
\xrightarrow{\delta} \bigoplus F_iP[1]
$$
。由于 $E$ 是紧的，对某些态射
$\delta_i:P\to F_iP[1]$，有
$\delta = \sum_{i = 1, \ldots, n} \delta_i$
。$\delta$ 与映射
$\bigoplus F_iP[1]\to\bigoplus F_iP[1]$ 的复合为零
（《导出范畴》引理 \ref{derived-lemma-composition-zero}），故
$\delta=0$（这是因为 $\bigoplus F_iP\to\bigoplus F_iP$ 在直和项
$F_iP$ 上是恒等映射与到 $F_{i-1}P$ 的包含映射之差）。因此，由
《导出范畴》引理 \ref{derived-lemma-split}，$E$ 上的恒等映射在
$D(A,\text{d})$ 中通过 $\bigoplus F_iP$ 分解。再次利用 $P$ 的紧性，
可知映射 $E\to\bigoplus F_iP$ 对某个 $n$ 通过
$\bigoplus_{i=1,\ldots,n}F_iP$ 分解。换言之，$E$ 上的恒等映射通过
$\bigoplus_{i=1,\ldots,n}F_iP$ 分解。由引理
\ref{dga-lemma-factor-through-nicer}，$E$ 的恒等映射可分解为
$E\to P\to E$，其中 $P$ 如引理陈述的第 (2) 项所述。于是已经证明
(1) 蕴含 (2)。

\medskip\noindent
假设 (2)。由《导出范畴》引理
\ref{derived-lemma-compact-objects-subcategory}，只需证明 $P$ 给出紧对象。
注意 $P$ 具有性质 (P)，故由引理 \ref{dga-lemma-hom-derived}，对任意微分
分次模 $M$ 都有
$$
\Hom_{D(A, \text{d})}(P, M) = \Hom_{K(A, \text{d})}(P, M)
$$
。由于 $D(A,\text{d})$ 中的直和由分次模的直和给出（引理
\ref{dga-lemma-derived-products}），问题约化为证明
$\Hom_{K(A,\text{d})}(P,M)$ 与直和交换。利用 $K(A,\text{d})$
是三角范畴、$\Hom$ 在第一变量中是上同调函子以及 $P$ 上的滤过，
可约化到 $P$ 是 $A$ 的平移的有限直和的情形，进而约化到显然的情形
$P=A[k]$。
\end{proof}

\begin{lemma}
\label{dga-lemma-compact-implies-bounded}
设 $(A,\text{d})$ 是微分分次代数。对 $D(A,\text{d})$ 的每个紧对象
$E$，存在整数 $a\leq b$，使得只要对 $i\in[a,b]$ 有
$H^i(M)=0$，就有 $\Hom_{D(A,\text{d})}(E,M)=0$。
\end{lemma}

\begin{proof}
注意，$D(A,\text{d})$ 中存在这样一对整数的对象所成的子范畴，
是 $D(A,\text{d})$ 的饱和严格满三角子范畴。因此由命题
\ref{dga-proposition-compact}，只需在 $E$ 由如下微分分次模 $P$ 表示时证明：
$P$ 有由微分分次子模组成的有限滤过 $F_\bullet$，且
$F_iP/F_{i-1}P$ 是 $A$ 的平移的有限直和。利用与三角的相容性，
只需对 $P=A$ 证明。此时
$\Hom_{D(A,\text{d})}(A,M)=H^0(M)$，并且取 $a=b=0$ 即得结论。
\end{proof}

\noindent
若 $(A,\text{d})$ 只是置于次数 $0$ 且微分为零的代数，或者更一般地，
只分布于有限多个次数，则确实可以得到对紧对象更精确的描述。

\begin{lemma}
\label{dga-lemma-compact}
设 $(A,\text{d})$ 是微分分次代数，并假设当 $|n|\gg0$ 时
$A^n=0$。设 $E$ 是 $D(A,\text{d})$ 的对象。下列条件等价：
\begin{enumerate}
\item $E$ 是紧对象；
\item $E$ 可由微分分次 $A$-模 $P$ 表示，其中 $P$ 作为分次 $A$-模
是有限投射的，并且对每个微分分次 $A$-模 $M$ 都满足
$\Hom_{K(A, \text{d})}(P, M) = \Hom_{D(A, \text{d})}(P, M)$
。
\end{enumerate}
\end{lemma}

\begin{proof}
设 $\mathcal{D}\subset K(A,\text{d})$ 是评注
\ref{dga-remark-source-graded-projective} 中讨论的三角子范畴。设 $P$
是 $\mathcal{D}$ 的对象，并且作为分次 $A$-模是有限投射的。则由评注
\ref{dga-remark-graded-projective-is-compact}，$P$ 表示
$D(A,\text{d})$ 的紧对象。

\medskip\noindent
为证明逆命题，设 $E$ 是 $D(A,\text{d})$ 的紧对象。固定引理
\ref{dga-lemma-compact-implies-bounded} 中的 $a\leq b$。必要时减小 $a$
并增大 $b$，还可假设当 $i\notin[a,b]$ 时 $H^i(E)=0$
（这由命题 \ref{dga-proposition-compact} 及对 $A$ 的假设得出）。此外，
固定整数 $c>0$，使得当 $|n|\geq c$ 时 $A^n=0$。

\medskip\noindent
由命题 \ref{dga-proposition-compact}，可见在 $D(A,\text{d})$ 中，$E$
是某个微分分次 $A$-模 $P$ 的直和项；$P$ 有由微分分次子模组成的有限
滤过 $F_\bullet$，且 $F_iP/F_{i-1}P$ 是 $A$ 的平移的有限直和。
特别地，$P$ 具有性质 (P)，并且由引理 \ref{dga-lemma-hom-derived}，
对任意微分分次模 $M$ 有
$\Hom_{D(A, \text{d})}(P, M) = \Hom_{K(A, \text{d})}(P, M)$
。换言之，$P$ 是评注 \ref{dga-remark-source-graded-projective} 中讨论的
三角子范畴 $\mathcal{D}\subset K(A,\text{d})$ 的对象。注意 $P$
作为分次 $A$-模是有限自由的。

\medskip\noindent
选取 $n>0$ 使 $b+4c-n<a$。以微分分次 $A$-模的自同态
$\varphi:P\to P$ 表示到 $E$ 上的投影算子。考虑 $K(A,\text{d})$
中的有区别三角
$$
P \xrightarrow{1 - \varphi} P \to C \to P[1]
$$
，其中 $C$ 是第一个箭头的锥。则 $C$ 是 $\mathcal{D}$ 的对象，
在 $D(A,\text{d})$ 中有 $C\cong E\oplus E[1]$，并且 $C$
是有限分次自由 $A$-模。接着考虑 $K(A,\text{d})$ 中的有区别三角
$$
C[1] \to C \to C' \to C[2]
$$
，其中 $C'$ 是某个态射 $C[1]\to C$ 的锥，而该态射表示
$D(A,\text{d})$ 中的复合
$$
C[1] \cong E[1] \oplus E[2] \to E[1] \to E \oplus E[1] \cong C
$$
。于是可见 $C'$ 表示 $E\oplus E[2]$。依此继续，可找到微分分次
$A$-模 $P$，它是 $\mathcal{D}$ 的对象，作为分次 $A$-模是有限自由的，
并表示 $E\oplus E[n]$。

\medskip\noindent
选取 $P$ 作为 $A$-模的齐次元素基 $x_i$（$i\in I$）。令
$d_i=\deg(x_i)$。令 $P_1$ 为由满足 $d_i\leq a-c-1$ 的
$x_i$ 与 $\text{d}(x_i)$ 生成的 $P$ 的 $A$-子模；令 $P_2$
为由满足 $d_i\geq b-n+c$ 的 $x_i$ 与 $\text{d}(x_i)$ 生成的
$P$ 的 $A$-子模。注意
\begin{enumerate}
\item $P_1$ 与 $P_2$ 是 $P$ 的微分分次子模；
\item 当 $t\geq a$ 时，$P_1^t=0$；
\item 当 $t\leq a-2c$ 时，$P_1^t=P^t$；
\item 当 $t\leq b-n$ 时，$P_2^t=0$；
\item 当 $t\geq b-n+2c$ 时，$P_2^t=P^t$。
\end{enumerate}
由于按 $n$ 的选取有 $b-n+2c\geq a-2c$，得到微分分次 $A$-模的
短正合列
$$
0 \to P_1 \cap P_2 \to P_1 \oplus P_2 \xrightarrow{\pi} P \to 0
$$
。由于 $P$ 作为分次 $A$-模是投射的，这是容许短正合列
（引理 \ref{dga-lemma-target-graded-projective}）。因此在
$K(A,\text{d})$ 中得到边界映射
$\delta:P\to(P_1\cap P_2)[1]$，见引理 \ref{dga-lemma-admissible-ses}。
由于 $P=E\oplus E[n]$，且 $P_1\cap P_2$ 集中于次数 $(b-n,a)$，
可知
$\Hom_{D(A,\text{d})}(E\oplus E[n],(P_1\cap P_2)[1])$ 为零。
又因 $P$ 是 $\mathcal{D}$ 的对象，故 $\delta$ 作为
$K(A,\text{d})$ 中的态射为零。由《导出范畴》引理
\ref{derived-lemma-split}，可找到映射 $s:P\to P_1\oplus P_2$，使得
$\pi \circ s = \text{id}_P + \text{d}h + h\text{d}$
，其中 $h:P\to P$ 的次数为 $-1$。由于 $P_1\oplus P_2\to P$
是满射，且 $P$ 作为分次 $A$-模是投射的，可以选取 $h$ 的齐次提升
$\tilde h:P\to P_1\oplus P_2$。把 $s$ 改为
$s+\text{d}\tilde h+\tilde h\text{d}$，即可得到
$\pi\circ s=\text{id}_P$。这意味着得到直和分解
$P=s^{-1}(P_1)\oplus s^{-1}(P_2)$。由于 $s^{-1}(P_2)$ 在次数
$\geq b-n+2c$ 上等于 $P$，可见 $s^{-1}(P_2)\to P\to E$
是拟同构，即 $D(A,\text{d})$ 中的同构。证明完毕。
\end{proof}








\section{导出范畴的等价}
\label{dga-section-equivalence}

\noindent
设 $R$ 是环，$(A,\text{d})$ 与 $(B,\text{d})$ 是微分分次 $R$-代数。
一个自然出现的问题是：$D(A,\text{d})$ 与 $D(B,\text{d})$ 作为
$R$-线性三角范畴等价意味着什么？这是相当微妙的问题，而且事实将表明，
它并不总是正确的问题。尽管如此，本节仍汇集若干保证这种等价成立的条件。

\medskip\noindent
我们强烈建议读者参阅关于此主题的开创性论文 \cite{Rickard}。

\begin{lemma}
\label{dga-lemma-qis-equivalence}
设 $R$ 是环，$(A,\text{d})\to(B,\text{d})$ 是微分分次 $R$-代数的
同态，并在上同调代数上诱导同构。则
$$
- \otimes_A^\mathbf{L} B : D(A, \text{d}) \to D(B, \text{d})
$$
给出三角范畴的 $R$-线性等价，其拟逆为限制函子 $N\mapsto N_A$。
\end{lemma}

\begin{proof}
由引理 \ref{dga-lemma-tensor-with-compact-fully-faithful}，函子
$M\mapsto M\otimes_A^\mathbf{L}B$ 是全忠实的。由引理
\ref{dga-lemma-tensor-hom-adjoint}，函子
$N\mapsto R\Hom(B,N)=N_A$ 是右伴随，见例
\ref{dga-example-map-hom-tensor}。显然 $R\Hom(B,-)$ 的核为零。
结论遂由《导出范畴》引理
\ref{derived-lemma-fully-faithful-adjoint-kernel-zero} 得出。
\end{proof}

\noindent
分析上述证明可见，下列推广无需额外工作即可得到。

\begin{lemma}
\label{dga-lemma-tilting-equivalence}
设 $R$ 是环，$(A,\text{d})$ 与 $(B,\text{d})$ 是微分分次 $R$-代数，
$N$ 是微分分次 $(A,B)$-双模。假设
\begin{enumerate}
\item $N$ 在 $D(B,\text{d})$ 中定义紧对象；
\item 若 $N'\in D(B,\text{d})$，且对每个 $n\in\mathbf{Z}$ 有
$\Hom_{D(B,\text{d})}(N,N'[n])=0$，则 $N'=0$；
\item 对每个 $k\in\mathbf{Z}$，映射
$H^k(A)\to\Hom_{D(B,\text{d})}(N,N[k])$ 都是同构。
\end{enumerate}
则
$$
- \otimes_A^\mathbf{L} N : D(A, \text{d}) \to D(B, \text{d})
$$
给出三角范畴的 $R$-线性等价。
\end{lemma}

\begin{proof}
由引理 \ref{dga-lemma-tensor-with-compact-fully-faithful}，函子
$M\mapsto M\otimes_A^\mathbf{L}N$ 是全忠实的。由引理
\ref{dga-lemma-tensor-hom-adjoint}，函子 $N'\mapsto R\Hom(N,N')$
是右伴随。由假设 (3)，$R\Hom(N,-)$ 的核为零。因此结论由
《导出范畴》引理 \ref{derived-lemma-fully-faithful-adjoint-kernel-zero}
得出。
\end{proof}

\begin{remark}
\label{dga-remark-tilting-equivalence}
在引理 \ref{dga-lemma-tilting-equivalence} 中，可把条件 (2) 替换为
$N$ 是 $D_{compact}(B,d)$ 的经典生成元这一条件，见《导出范畴》命题
\ref{derived-proposition-generator-versus-classical-generator}。
此外，若已知 $R\Hom(N,B)$ 是 $D(A,\text{d})$ 的紧对象，则只需检验
$N$ 是 $D_{compact}(B,\text{d})$ 的弱生成元。证明从略；如果 Stacks
project 今后用到此结论，我们会在这里补上证明。
\end{remark}

\noindent
有时把下述引理中的 $B$-模 $P$ 称为“$(A,B)$-倾斜复形”。

\begin{lemma}
\label{dga-lemma-rickard}
设 $R$ 是环，$(A,\text{d})$ 与 $(B,\text{d})$ 是微分分次 $R$-代数。
假设 $A=H^0(A)$。下列条件等价：
\begin{enumerate}
\item $D(A,\text{d})$ 与 $D(B,\text{d})$ 作为 $R$-线性三角范畴等价；
\item 存在 $D(B,\text{d})$ 的对象 $P$，使得
\begin{enumerate}
\item $P$ 是 $D(B,\text{d})$ 的紧对象；
\item 若 $N\in D(B,\text{d})$，且对每个 $i\in\mathbf{Z}$ 有
$\Hom_{D(B,\text{d})}(P,N[i])=0$，则 $N=0$；
\item 当 $i\neq0$ 时，$\Hom_{D(B,\text{d})}(P,P[i])=0$；当 $i=0$
时，该群等于 $A$。
\end{enumerate}
\end{enumerate}
由 (2) 构造的等价 $D(A,\text{d})\to D(B,\text{d})$ 把 $A$ 映到 $P$。
\end{lemma}

\begin{proof}
设 $F:D(A,\text{d})\to D(B,\text{d})$ 是等价。则 $F$ 把紧对象映到
紧对象。因此 $P=F(A)$ 是紧的，即 (2)(a) 成立。条件 (2)(b) 与
(2)(c) 由 $F$ 是等价立即得出。

\medskip\noindent
设 $P$ 是 (2) 中的对象，以具有性质 (P) 的微分分次模表示 $P$。令
$$
(E, \text{d}) = \Hom_{\text{Mod}^{dg}_{(B, \text{d})}}(P, P)
$$
。由引理 \ref{dga-lemma-hom-derived} 与假设 (2)(c)，有 $H^0(E)=A$，
且当 $k\neq0$ 时 $H^k(E)=0$。把 $P$ 看作 $(E,B)$-双模，并使用引理
\ref{dga-lemma-tilting-equivalence} 与假设 (2)(b)，得到等价
$$
D(E, \text{d}) \to D(B, \text{d})
$$
，它把 $E$ 映到 $P$。令 $E'\subset E$ 为如下微分分次 $R$-子代数：
$$
(E')^i = \left\{
\begin{matrix}
E^i & \text{若 }i < 0 \\
\Ker(E^0 \to E^1) & \text{若 }i = 0 \\
0 & \text{若 }i > 0
\end{matrix}
\right.
$$
于是有微分分次代数的拟同构
$(A,\text{d})\leftarrow(E',\text{d})\rightarrow(E,\text{d})$。
因此由引理 \ref{dga-lemma-qis-equivalence} 得到等价
$$
D(A, \text{d}) \leftarrow D(E', \text{d}) \rightarrow D(E, \text{d})
\rightarrow D(B, \text{d})
$$
。
\end{proof}

\begin{remark}
\label{dga-remark-lift-equivalence-to-dga}
设 $R$ 是环，$(A,\text{d})$ 与 $(B,\text{d})$ 是微分分次 $R$-代数。
设给定三角范畴的 $R$-线性等价
$$
F : D(A, \text{d}) \longrightarrow D(B, \text{d})
$$
。令 $N=F(A)$。则 $N$ 是微分分次 $B$-模。由于 $F$ 是等价，且
$A$ 是 $D(A,\text{d})$ 的紧对象，可知 $N$ 是 $D(B,\text{d})$
的紧对象。由于 $A$ 生成 $D(A,\text{d})$ 且 $F$ 是等价，可见 $N$
生成 $D(B,\text{d})$。最后，
$H^k(A)=\Hom_{D(A,\text{d})}(A,A[k])$，而由 $F$ 是等价可知，
对每个 $k$，$F$ 诱导同构
$H^k(A)=\Hom_{D(B,\text{d})}(N,N[k])$。要得出存在由引理
\ref{dga-lemma-tilting-equivalence} 的构造产生的等价
$D(A,\text{d})\to D(B,\text{d})$，所需的只是在 $N$ 上给出一个
与微分相容、并与给定右 $B$-模结构交换的左 $A$-模结构。事实上，
只需在把 $N$ 替换为与之拟同构的微分分次 $B$-模后完成此事。
在某些情形下可以构造这一模结构。例如，若假设 $F$ 可提升为微分分次函子
$$
F^{dg} :
\text{Mod}^{dg}_{(A, \text{d})}
\longrightarrow
\text{Mod}^{dg}_{(B, \text{d})}
$$
（记号见例 \ref{dga-example-dgm-dg-cat}），它位于相应微分分次范畴之间，
则上述要求成立。另一种情形将在下面的命题中讨论。
\end{remark}

\begin{proposition}
\label{dga-proposition-rickard}
设 $R$ 是环，$(A,\text{d})$ 与 $(B,\text{d})$ 是微分分次 $R$-代数，
$F:D(A,\text{d})\to D(B,\text{d})$ 是三角范畴的 $R$-线性等价。
假设
\begin{enumerate}
\item $A=H^0(A)$；
\item $B$ 作为 $R$-模复形是 K-平坦的。
\end{enumerate}
则存在引理 \ref{dga-lemma-tilting-equivalence} 所述的 $(A,B)$-双模 $N$。
\end{proposition}

\begin{proof}
如上面的评注 \ref{dga-remark-lift-equivalence-to-dga}，在
$D(B,\text{d})$ 中令 $N=F(A)$。可假设 $N$ 是具有性质 (P) 的微分分次
$B$-模。令
$$
(E, \text{d}) = \Hom_{\text{Mod}^{dg}_{(B, \text{d})}}(N, N)
$$
。由引理 \ref{dga-lemma-hom-derived}，有 $H^0(E)=A$，且当 $k\neq0$
时 $H^k(E)=0$。此外，由评注 \ref{dga-remark-lift-equivalence-to-dga}
中的讨论及引理 \ref{dga-lemma-tilting-equivalence}，可见 $N$ 作为
$(E,B)$-双模诱导等价
$-\otimes_E^\mathbf{L}N:D(E,\text{d})\to D(B,\text{d})$。
令 $E'\subset E$ 为如下微分分次 $R$-子代数：
$$
(E')^i = \left\{
\begin{matrix}
E^i & \text{若 }i < 0 \\
\Ker(E^0 \to E^1) & \text{若 }i = 0 \\
0 & \text{若 }i > 0
\end{matrix}
\right.
$$
于是有微分分次代数的拟同构
$(A,\text{d})\leftarrow(E',\text{d})\rightarrow(E,\text{d})$。
因此由引理 \ref{dga-lemma-qis-equivalence} 得到等价
$$
D(A, \text{d}) \leftarrow D(E', \text{d}) \rightarrow D(E, \text{d})
\rightarrow D(B, \text{d})
$$
。注意，左侧箭头的拟逆 $D(A,\text{d})\to D(E',\text{d})$ 由
$M\mapsto M\otimes_A^\mathbf{L}A$ 给出，其中把 $A$ 看作
$(A,E')$-双模，见例 \ref{dga-example-map-hom-tensor}。另一方面，函子
$D(E',\text{d})\to D(B,\text{d})$ 由
$M\mapsto M\otimes_{E'}^\mathbf{L}N$ 给出，其中 $N$ 如上。
由引理 \ref{dga-lemma-compose-tensor-functors} 即得结论。
\end{proof}

\begin{remark}
\label{dga-remark-rickard}
设 $A,B,F,N$ 如命题 \ref{dga-proposition-rickard}。尚不清楚 $F$ 与函子
$G(-)=-\otimes_A^\mathbf{L}N$ 是否同构。按构造，在
$D(B,\text{d})$ 中有同构 $N=G(A)\to F(A)$。当 $M$ 是分次投射的
微分分次 $A$-模（例如 $A$ 的平移的直和）时，可以直接把它扩张为
函子同构 $G(M)\to F(M)$。进而，当 $M$ 是此类模之间某个映射的锥时，
可得 $G(M)\cong F(M)$。一般情形下是否还有更多结论，我们并不知道。
\end{remark}

\begin{lemma}
\label{dga-lemma-rickard-rings}
设 $R$ 是环，$A,B$ 是 $R$-代数。下列条件等价：
\begin{enumerate}
\item 存在三角范畴的 $R$-线性等价 $D(A)\to D(B)$；
\item 存在 $D(B)$ 的对象 $P$，使得
\begin{enumerate}
\item $P$ 可由有限投射 $B$-模组成的有限复形表示；
\item 若 $K\in D(B)$，且对每个 $i\in\mathbf{Z}$ 有
$\Ext^i_B(P,K)=0$，则 $K=0$；
\item 当 $i\neq0$ 时，$\Ext^i_B(P,P)=0$；当 $i=0$ 时，该群等于 $A$。
\end{enumerate}
\end{enumerate}
此外，若 $B$ 作为 $R$-模是平坦的，则上述条件还等价于
\begin{enumerate}
\item[(3)] 存在 $(A,B)$-双模 $N$，使得
$-\otimes_A^\mathbf{L}N:D(A)\to D(B)$ 是等价。
\end{enumerate}
\end{lemma}

\begin{proof}
(1) 与 (2) 的等价，是引理 \ref{dga-lemma-rickard} 结合引理
\ref{dga-lemma-compact} 对 $D(B)$ 的紧对象之刻画所得的特例
（略去一个小细节）。当 $B$ 是 $R$-平坦的时，与 (3) 的等价由命题
\ref{dga-proposition-rickard} 得出。
\end{proof}

\begin{remark}
\label{dga-remark-centers}
设 $R$ 是环，$A,B$ 是 $R$-代数。若 $D(A)$ 与 $D(B)$ 作为
$R$-线性三角范畴等价，则 $A$ 与 $B$ 的中心作为 $R$-代数同构。
特别地，若 $A,B$ 交换，则 $A\cong B$。这一颇为棘手的证明见
\cite[Proposition 9.2]{Rickard} 或 \cite[Proposition 6.3.2]{KZ}。
另一种途径或许是使用 Hochschild 上同调（见下面的评注）。
\end{remark}

\begin{remark}
\label{dga-remark-hochschild-cohomology}
设 $R$ 是环，$(A,\text{d})$ 与 $(B,\text{d})$ 是导出等价的微分分次
$R$-代数，即存在三角范畴的 $R$-线性等价
$D(A,\text{d})\to D(B,\text{d})$。我们想证明 $(A,\text{d})$ 与
$(B,\text{d})$ 的某些不变量相同。在许多情形下，可以更好地控制局面。
例如，可能存在形如
$$
- \otimes_A \Omega : D(A, \text{d}) \longrightarrow D(B, \text{d})
$$
的等价，其中 $\Omega$ 是某个微分分次 $(A,B)$-双模
（命题 \ref{dga-proposition-rickard} 的情形即如此；若等价来自几何构造，
也常常如此）。若该函子的拟逆又由
$$
- \otimes_B^\mathbf{L} \Omega' : D(B, \text{d}) \longrightarrow D(A, \text{d})
$$
给出，其中 $\Omega'$ 是微分分次 $(B,A)$-双模
（同前，这样的模 $\Omega'$ 在实际中常常存在），则可考虑双模导出范畴
之间的函子
$$
D(A^{opp} \otimes_R A, \text{d})
\longrightarrow
D(B^{opp} \otimes_R B, \text{d}),\quad
M \longmapsto \Omega' \otimes^\mathbf{L}_A M \otimes_A^\mathbf{L} \Omega
$$
（使用引理 \ref{dga-lemma-bimodule-over-tensor} 把双模化为右模）。注意，
此函子把 $(A,A)$-双模 $A$ 映到 $(B,B)$-双模 $B$。在适当条件下
（例如 $A,B,\Omega$ 在 $R$ 上的平坦性等），此函子也将是等价。
若确实如此，则得到 Hochschild 上同调群的同构
$$
HH^i(A, \text{d}) =
\Hom_{D(A^{opp} \otimes_R A, \text{d})}(A, A[i])
\longrightarrow
\Hom_{D(B^{opp} \otimes_R B, \text{d})}(B, B[i]) =
HH^i(B, \text{d}).
$$
例如，若 $A=H^0(A)$，则 $HH^0(A,\text{d})$ 等于 $A$ 的中心，
这就给出了评注 \ref{dga-remark-centers} 中所述结果的概念性证明。
如果今后用到本评注，我们将在此给出精确陈述与详细证明。
\end{remark}



\section{微分分次代数的分辨}
\label{dga-section-resolution-dgas}

\noindent
设 $R$ 是环。按我们的约定，集合 $S$ 上的自由 $R$-代数
$R\langle S\rangle$ 是如下代数：它以表达式
$$
s_1 s_2 \ldots s_n
$$
为 $R$-基，其中 $n\geq0$，而 $s_1,\ldots,s_n\in S$ 是 $S$
中元素的一个序列。乘法由连接给出：
$$
(s_1 s_2 \ldots s_n) \cdot (s'_1 s'_2 \ldots s'_m) =
s_1 \ldots s_n s'_1 \ldots s'_m
$$
。此代数由如下性质刻画：对每个 $R$-代数 $A$，映射
$$
\Mor_{R\text{-代数}}(R\langle S \rangle, A) \to
\text{Map}(S, A),\quad
\varphi \longmapsto (s \mapsto \varphi(s))
$$
是双射。

\medskip\noindent
在分次 $R$-代数范畴中，集合 $S$ 应带有分次，我们把它看作映射
$\deg:S\to\mathbf{Z}$。于是 $R\langle S\rangle$ 带有分次，使单项式
的次数为
$$
\deg(s_1 s_2 \ldots s_n) = \deg(s_1) + \ldots + \deg(s_n)
$$
。在此设定下，对每个分次 $R$-代数 $A$，映射
$$
\Mor_{\text{分次 }R\text{-代数}}(R\langle S \rangle, A) \to
\text{Map}_{\text{分次集合}}(S, A),\quad
\varphi \longmapsto (s \mapsto \varphi(s))
$$
是双射。

\medskip\noindent
若 $A$ 是分次 $R$-代数，$S$ 是分次集合，也可类似地构造
$A\langle S\rangle$。$A\langle S\rangle$ 的元素是形如
$$
a_0 s_1 a_1 s_2 \ldots a_{n - 1} s_n a_n
$$
的元素之和，其中 $a_i\in A$，并模去使这些表达式关于
$(a_0,\ldots,a_n)$ 为 $R$-多重线性的关系。因此，对 $S$ 中元素的
每个序列 $s_1,\ldots,s_n$，都有包含
$$
A \otimes_R \ldots \otimes_R A \subset A\langle S \rangle
$$
，而该代数是这些项的直和。由此定义，读者可证明映射
$$
\Mor_{\text{分次 }R\text{-代数}}(A\langle S \rangle, B) \to
\Mor_{\text{分次 }R\text{-代数}}(A, B) \times
\text{Map}_{\text{分次集合}}(S, B),
$$
把 $\varphi$ 映到 $(\varphi|_A,(s\mapsto\varphi(s)))$；
对每个分次 $R$-代数 $A$，该映射是双射。注意，若 $A$ 是自由分次
$R$-代数，则 $A\langle S\rangle$ 也是自由分次 $R$-代数。

\medskip\noindent
设 $A$ 是微分分次 $R$-代数，$S$ 是分次集合。再假设对每个
$s\in S$，给定齐次元素 $f_s\in A$，满足
$\deg(f_s)=\deg(s)+1$ 且 $\text{d}f_s=0$。则在
$A\langle S\rangle$ 上存在唯一的微分分次代数结构，使得
$\text{d}(s)=f_s$。例如，给定 $a,b,c\in A$ 与 $s,t\in S$，定义
\begin{align*}
\text{d}(asbtc)
& =
\text{d}(a)sbtc + (-1)^{\deg(a)}a f_s b t c +
(-1)^{\deg(a) + \deg(s)} as\text{d}(b)tc \\
& + (-1)^{\deg(a) + \deg(s) + \deg(b)} asb f_t c +
(-1)^{\deg(a) + \deg(s) + \deg(b) + \deg(t)} asbt\text{d}(c)
\end{align*}
。细节从略。

\begin{lemma}
\label{dga-lemma-K-flat-resolution}
设 $R$ 是环，$(B,\text{d})$ 是微分分次 $R$-代数。存在微分分次
$R$-代数的拟同构 $(A,\text{d})\to(B,\text{d})$，并具有下列性质：
\begin{enumerate}
\item $A$ 作为 $R$-模复形是 K-平坦的；
\item $A$ 是自由分次 $R$-代数。
\end{enumerate}
\end{lemma}

\begin{proof}
首先断言，可找到具有 (1)、(2) 且在上同调上诱导满射的
$(A_0,\text{d})\to(B,\text{d})$。具体地，取分次集合 $S$；
对每个 $s\in S$，取次数为 $\deg(s)$ 的齐次元素
$b_s\in\Ker(d:B\to B)$，使其在 $H^*(B)$ 中的类
$\overline b_s$ 作为 $R$-模生成 $H^*(B)$。可令
$A_0=R\langle S\rangle$，微分为零，并以 $s\mapsto b_s$
定义 $A_0\to B$。

\medskip\noindent
给定在上同调上诱导满射的 $A_0\to B$，归纳构造序列
$$
A_0 \to A_1 \to A_2 \to \ldots B
$$
。给定 $A_n\to B$，令 $S_n$ 为分次集合；对每个 $s\in S_n$，
令 $a_s\in\Ker(\text{d}:A_n\to A_n)$ 为次数
$\deg(s)+1$ 的齐次元素，其在 $H^*(A_n)$ 中的类
$\overline a_s$ 映到 $H^*(B)$ 中的零。取 $S_n$ 足够大，使诸
$\overline a_s$ 作为 $R$-模生成
$\Ker(H^*(A_n)\to H^*(B))$。然后令
$$
A_{n + 1} = A_n\langle S_n \rangle
$$
，其微分由 $\text{d}(s)=a_s$ 给出，见上面的讨论。于是每个
$(A_n,\text{d})$ 都满足 (1)、(2)；细节从略。映射
$H^*(A_n)\to H^*(B)$ 是满射，因为 $n=0$ 时已经如此。

\medskip\noindent
显然 $A=\colim A_n$ 是自由分次 $R$-代数。由《代数进阶》引理
\ref{more-algebra-lemma-colimit-K-flat}，它是 K-平坦的。映射
$H^*(A)\to H^*(B)$ 是同构：它是满射，也是单射。事实上，
$H^*(A)$ 的每个元素都来自某个 $H^*(A_n)$ 中的元素；若它在
$H^*(B)$ 中消失，则它在 $H^*(A_{n+1})$ 中消失，因而也在
$H^*(A)$ 中消失。
\end{proof}

\noindent
作为应用，我们证明引理
\ref{dga-lemma-compose-tensor-functors-general-algebra} 的“正确”版本。

\begin{lemma}
\label{dga-lemma-compose-tensor-functors-tor}
设 $R$ 是环，$(A,\text{d})$、$(B,\text{d})$ 和 $(C,\text{d})$
是微分分次 $R$-代数。假设 $A\otimes_RC$ 在 $D(R)$ 中表示
$A\otimes_R^\mathbf{L}C$。设 $N$ 是微分分次 $(A,B)$-双模，
$N'$ 是微分分次 $(B,C)$-双模。则复合
$$
\xymatrix{
D(A, \text{d}) \ar[rr]^{- \otimes_A^\mathbf{L} N} & &
D(B, \text{d}) \ar[rr]^{- \otimes_B^\mathbf{L} N'} & &
D(C, \text{d})
}
$$
同构于 $-\otimes_A^\mathbf{L}N''$，其中 $N''$ 是某个微分分次
$(A,C)$-双模。
\end{lemma}

\begin{proof}
利用引理 \ref{dga-lemma-K-flat-resolution}，选取拟同构
$(B',\text{d})\to(B,\text{d})$，其中 $B'$ 作为 $R$-模复形是
K-平坦的。由引理 \ref{dga-lemma-qis-equivalence}，函子
$-\otimes_{B'}^\mathbf{L}B:D(B',\text{d})\to D(B,\text{d})$
是等价，其拟逆由限制给出。注意，限制典范同构于函子
$-\otimes_B^\mathbf{L}B:D(B,\text{d})\to D(B',\text{d})$，
其中把 $B$ 看作 $(B,B')$-双模。因此，只需对下列复合证明引理：
$$
D(A) \to D(B) \to D(B'),\quad
D(B') \to D(B) \to D(C),\quad
D(A) \to D(B') \to D(C).
$$
第一个复合由引理 \ref{dga-lemma-compose-tensor-functors} 处理，因为 $B'$
作为 $R$-模复形是 K-平坦的。第二个复合成立，因为
$B \otimes_B^\mathbf{L} N' = N' = B \otimes_B N'$
，故可应用引理 \ref{dga-lemma-compose-tensor-functors-general}。于是约化到
$B$ 作为 $R$-模复形是 K-平坦的情形。

\medskip\noindent
假设 $B$ 作为 $R$-模复形是 K-平坦的。只需证明
(\ref{dga-equation-plain-versus-derived-algebras}) 是同构，见引理
\ref{dga-lemma-compose-tensor-functors-general-algebra}。选取拟同构
$L\to A$，其中 $L$ 是具有性质 (P) 的微分分次 $R$-模。显然
$P=L\otimes_RB$ 作为微分分次 $B$-模具有性质 (P)。因此必须证明
$P\to A\otimes_RB$ 诱导拟同构
$$
P \otimes_B (B \otimes_R C)
\longrightarrow
(A \otimes_R B) \otimes_B (B \otimes_R C)
$$
。可将其改写为
$$
P \otimes_R B \otimes_R C \longrightarrow A \otimes_R B \otimes_R C
$$
。由于 $B$ 作为 $R$-模复形是 K-平坦的，由《代数进阶》引理
\ref{more-algebra-lemma-K-flat-quasi-isomorphism}，只需证明
$$
P \otimes_R C \to A \otimes_R C
$$
是拟同构，而这恰是我们的假设。
\end{proof}

\noindent
下列引理其实不属于本节，但似乎没有更自然的安置位置。

\begin{lemma}
\label{dga-lemma-countable}
设 $(A,\text{d})$ 是微分分次代数，且对每个 $i$，$H^i(A)$ 都可数。
设 $M$ 是 $D(A,\text{d})$ 的对象。下列条件等价：
\begin{enumerate}
\item $M=\text{hocolim},E_n$，其中 $E_n$ 是 $D(A,\text{d})$ 的紧对象；
\item 对每个 $i$，$H^i(M)$ 都可数。
\end{enumerate}
\end{lemma}

\begin{proof}
假设 (1) 成立。由《导出范畴》引理
\ref{derived-lemma-cohomology-of-hocolim}，有
$H^i(M)=\colim H^i(E_n)$。因此只需证明对每个 $n$，$H^i(E_n)$
可数。由命题 \ref{dga-proposition-compact}，可见 $E_n$ 在
$D(A,\text{d})$ 中同构于某个微分分次模 $P$ 的直和项；$P$ 有由微分
分次子模组成的有限滤过 $F_\bullet$，且 $F_jP/F_{j-1}P$ 是 $A$
的平移的有限直和。依假设，群 $H^i(F_jP/F_{j-1}P)$ 可数。
对滤过长度归纳并使用上同调长正合列，即得 (2)。更有趣的是反向蕴含。

\medskip\noindent
我们断言，存在可数微分分次子代数 $A'\subset A$，使包含映射
$A'\to A$ 在上同调上定义同构。为构造 $A'$，选取可数微分分次子代数
$$
A_1 \subset A_2 \subset A_3 \subset \ldots
$$
，使得：(a) $H^i(A_1)\to H^i(A)$ 是满射；(b) 对 $n>1$，映射
$H^i(A_{n-1})\to H^i(A_n)$ 的核与映射
$H^i(A_{n-1})\to H^i(A)$ 的核相同。为构造 $A_1$，取生成 $A$
的上同调（作为环或作为分次阿贝尔群）的任意可数上链集
$S\subset A$，并令 $A_1$ 为由 $S$ 生成的 $A$ 的微分分次子代数。
给定 $A_{n-1}$ 后，为构造 $A_n$，对每个在 $H^i(A)$ 中映到零的上链
$a\in A_{n-1}^i$，选取 $s_a\in A^{i-1}$ 使
$\text{d}(s_a)=a$，并令 $A_n$ 为由 $A_{n-1}$ 与诸元素 $s_a$
生成的 $A$ 的微分分次子代数。最后取 $A'=\bigcup A_n$。

\medskip\noindent
由引理 \ref{dga-lemma-qis-equivalence}，限制映射
$D(A,\text{d})\to D(A',\text{d})$，$M\mapsto M_{A'}$，是等价。
由于 $M$ 与 $M_{A'}$ 的上同调群相同，只需对 $(A',\text{d})$
证明蕴含 (2) $\Rightarrow$ (1)。

\medskip\noindent
假设 $A$ 可数。用与上面完全相同类型的论证可见，对
$D(A,\text{d})$ 中的 $M$，下列两者等价：对每个 $i$，$H^i(M)$
可数；$M$ 可由可数微分分次模表示。因此，为证明蕴含
(2) $\Rightarrow$ (1)，可约化到下一段所述情形。

\medskip\noindent
假设 $A$ 可数，且 $M$ 是可数微分分次 $A$-模。我们断言，存在微分分次
$A$-模同态 $P\to M$，使得
\begin{enumerate}
\item $P\to M$ 是拟同构；
\item $P$ 具有性质 (P)；
\item $P$ 可数。
\end{enumerate}
查看引理 \ref{dga-lemma-resolve} 中构造 P-分辨的证明可见，只需说明在可数
微分分次模的设定下仍能证明引理 \ref{dga-lemma-good-quotient}。
这由该证明立即得出。

\medskip\noindent
假设 $A$ 可数，且 $M$ 是具有性质 (P) 的可数微分分次模。选取由微分分次
子模组成的滤过
$$
0 = F_{-1}P \subset F_0P \subset F_1P \subset \ldots \subset P
$$
，使得
\begin{enumerate}
\item $P=\bigcup F_pP$；
\item $F_iP\to F_{i+1}P$ 是容许单态射；
\item 对某些集合 $J_i$ 与整数 $k_j$，有微分分次模的同构
$F_iP/F_{i - 1}P \to \bigoplus_{j \in J_i} A[k_j]$
。
\end{enumerate}
当然，对每个 $i$，$J_i$ 都可数。对每个 $i$ 与 $j\in J_i$，
选取次数为 $k_j$ 的 $x_{i,j}\in F_iP$，使其在
$F_iP/F_{i-1}P$ 中的像生成对应于 $j$ 的直和项。

\medskip\noindent
断言：给定 $n$ 以及有限子集 $S_i\subset J_i$（$i=1,\ldots,n$），
存在有限子集 $S_i\subset T_i\subset J_i$（$i=1,\ldots,n$），使得
$P'=\bigoplus_{i\leq n}\bigoplus_{j\in T_i}Ax_{i,j}$ 是 $P$
的微分分次子模。这已在引理 \ref{dga-lemma-factor-through-nicer} 的证明中
说明，也容易直接证明：诸元素 $x_{i,j}$ 作为右 $A$-模自由生成 $P$。
$P$ 的结构表明
$$
\text{d}(x_{i, j}) = \sum\nolimits_{i' < i} x_{i', j'}a_{i', j'}
$$
，其中求和当然是有限的。因此，给定 $S_0,\ldots,S_n$，可先选取
$S_0\subset S'_0,\ldots,S_{n-1}\subset S'_{n-1}$，使得对每个
$j\in S_n$ 都有
$\text{d}(x_{n, j}) \in \bigoplus_{i' < n, j' \in S'_{i'}} x_{i', j'}A$
。然后对 $n$ 归纳，可以选取
$S'_0\subset T_0,\ldots,S'_{n-1}\subset T_{n-1}$，保证
$\bigoplus_{i'<n,j'\in T_{i'}}x_{i',j'}A$ 是微分分次 $A$-子模。
令 $T_n=S_n$，便得到所需的
$P'=\bigoplus_{i\leq n,j\in T_i}x_{i,j}A$。

\medskip\noindent
由此断言显然可知，$P=\bigcup P'_n$ 是上述 $P'_n$ 的可数递增并。
按构造，每个 $P'_n$ 都是具有性质 (P) 的微分分次模，其滤过有限，
且逐次商是 $A$ 的平移的有限直和。因此 $P'_n$ 定义
$D(A,\text{d})$ 的紧对象，例如见命题 \ref{dga-proposition-compact}。
由引理 \ref{dga-lemma-homotopy-colimit}，在 $D(A,\text{d})$ 中有
$P=\text{hocolim},P'_n$，故蕴含 (2) $\Rightarrow$ (1) 证明完毕。
\end{proof}
